\documentclass[11pt]{amsart}

% ============================================================================
% The cohomological Hall algebra of the resolved conifold.
% Date: 2026-04-23.
% ============================================================================

\usepackage{amsmath,amssymb,amsthm,mathtools,enumitem,hyperref,microtype}
\usepackage{tikz-cd}
\usepackage{booktabs}
\usepackage{stmaryrd}
\usepackage[margin=1in]{geometry}
\newtheorem{theorem}{Theorem}[section]
\newtheorem{proposition}[theorem]{Proposition}
\newtheorem{lemma}[theorem]{Lemma}
\newtheorem{corollary}[theorem]{Corollary}
\newtheorem{conjecture}[theorem]{Conjecture}
\theoremstyle{definition}
\newtheorem{definition}[theorem]{Definition}
\newtheorem{remark}[theorem]{Remark}
\newtheorem{example}[theorem]{Example}

\providecommand{\CoHA}{\mathrm{CoHA}}
\providecommand{\bY}{\mathbf{Y}}
\providecommand{\conBy}{Y_{\mathrm{con}}}
\providecommand{\cO}{\mathcal{O}}
\providecommand{\cA}{\mathcal{A}}
\providecommand{\cF}{\mathcal{F}}
\providecommand{\cG}{\mathcal{G}}
\providecommand{\cM}{\mathcal{M}}
\providecommand{\cN}{\mathcal{N}}
\providecommand{\cW}{\mathcal{W}}
\providecommand{\cS}{\mathcal{S}}
\providecommand{\cT}{\mathcal{T}}
\providecommand{\cU}{\mathcal{U}}
\providecommand{\cV}{\mathcal{V}}
\providecommand{\cX}{\mathcal{X}}
\providecommand{\cY}{\mathcal{Y}}
\providecommand{\cZ}{\mathcal{Z}}
\providecommand{\Hom}{\mathrm{Hom}}
\providecommand{\Ext}{\mathrm{Ext}}
\providecommand{\End}{\mathrm{End}}
\providecommand{\bZ}{\mathbb{Z}}
\providecommand{\bR}{\mathbb{R}}
\providecommand{\bC}{\mathbb{C}}
\providecommand{\bP}{\mathbb{P}}
\providecommand{\bL}{\mathbb{L}}
\providecommand{\bQ}{\mathbb{Q}}
\providecommand{\bN}{\mathbb{N}}
\providecommand{\bT}{\mathbb{T}}
\providecommand{\bH}{\mathbb{H}}
\providecommand{\Z}{\mathbb{Z}}
\providecommand{\R}{\mathbb{R}}
\providecommand{\C}{\mathbb{C}}
\providecommand{\N}{\mathbb{N}}
\providecommand{\Q}{\mathbb{Q}}
\providecommand{\fgl}{\mathfrak{gl}}
\providecommand{\fsl}{\mathfrak{sl}}
\providecommand{\fn}{\mathfrak{n}}
\providecommand{\fg}{\mathfrak{g}}
\providecommand{\fh}{\mathfrak{h}}
\providecommand{\fU}{\mathfrak{U}}
\providecommand{\fu}{\mathfrak{u}}
\providecommand{\fD}{\mathfrak{D}}
\providecommand{\frakg}{\mathfrak{g}}
\providecommand{\HH}{\mathrm{HH}}
\providecommand{\ClaimStatusTheorem}{\textbf{\textsf{[Theorem]}}}
\providecommand{\ClaimStatusProposition}{\textbf{\textsf{[Proposition]}}}
\providecommand{\ClaimStatusProvedElsewhere}{\textbf{\textsf{[Proved elsewhere]}}}
\providecommand{\ClaimStatusConjectured}{\textbf{\textsf{[Conjectural]}}}
\providecommand{\ClaimStatusOpen}{\textbf{\textsf{[Open]}}}
\providecommand{\ClaimStatusConventionDependent}{\textbf{\textsf{[Convention-dependent]}}}
\providecommand{\ClaimStatusPartialPrimaryLit}{\textbf{\textsf{[Primary-lit partial]}}}

\title{The cohomological Hall algebra of the resolved conifold}
\author{}
\date{2026-04-23}

\begin{document}
\maketitle

\begin{abstract}
\noindent The critical cohomological Hall algebra of the resolved conifold
$\bY = \mathrm{Tot}(\cO_{\bP^1}(-1) \oplus \cO_{\bP^1}(-1))$ admits five
independent constructions --- direct Klebanov--Witten super-shuffle, two-chart
\v{C}ech gluing of two $Y^+(\widehat{\fgl}_1)$ stalks, Van den Bergh tilting
transport from $D^b(\mathrm{Coh}\,\bY)$ to $D^b(J(Q_{\mathrm{con}},
W_{\mathrm{con}})\text{-mod})$, RTT super-Yangian via Gauss decomposition, and
chiral vertex-algebra realisation via the Gaiotto--Rap\v{c}\'ak corner
$Y_{0,1,1}[\psi]$ under the Li--Yamazaki toric-quiver dictionary (the chiral
identification is conjectural at non-generic $\psi$, proved on the self-dual
locus $\epsilon_1 + \epsilon_2 = 0$; see
Remark~\ref{rem:corner-assignment-scope}). All five agree on the leading
structure constant $(\varepsilon_1 \varepsilon_2)^{-1}$ at the $(1,1)$-weight
imaginary root; the subleading $c_{11}$ asymmetric correction is conjectural
(Conjecture~\ref{conj:hopf-pairing}, Remark~\ref{rem:pairing-checks}). The
Drinfeld double is the affine super-Yangian $Y(\widehat{\fgl}(1|1))$ at
rational, trigonometric, and toroidal levels; the elliptic lift is
quasi-Hopf via the Felder--Jimbo--Konno dynamical $R$-matrix. The
construction transfers to any local toric CY$_3$ without compact
$4$-cycles; a sharp obstruction prevents direct transport to compact CY$_3$,
where automorphic-form machinery (Borcherds products, Gritsenko--Nikulin)
replaces chart-wise shuffle presentations.
\end{abstract}

\tableofcontents

% ============================================================================
\section{Geometric datum}
\label{sec:geometric-datum}
% ============================================================================

The resolved conifold
\[
  \bY \;=\; \mathrm{Tot}(\pi\colon E \to \bP^1), \qquad
  E \;=\; \cO_{\bP^1}(-1) \oplus \cO_{\bP^1}(-1),
\]
is a smooth quasi-projective Calabi--Yau threefold of complex dimension three.
As a rank-$2$ holomorphic vector bundle over $\bP^1$ its topological Euler
number is $\chi_{\mathrm{top}}(\bY) = \chi_{\mathrm{top}}(\bP^1) = 2$; its
holomorphic Euler characteristic $\chi(\cO_{\bY})$ is infinite (global
sections of $\cO_{\bY}$ include all polynomials in the fibre coordinates).
The second homology $H_2(\bY; \bZ) = \bZ[\bP^1]$ is generated by the class of
the zero section, a rational curve of genus zero; the effective cone of curve
classes is $\bZ_{\geq 0}\beta$ with $\beta = [\bP^1]$.

Triviality of the canonical bundle is a first-principles consequence of the
total-space formula: for any vector bundle $E$ over a smooth base $B$ the
canonical bundle of $\mathrm{Tot}(E)$ is
$K_{\mathrm{Tot}(E)} = \pi^*\bigl(K_B \otimes \det(E^\vee)\bigr)$
(Hartshorne III.\,8, relative-cotangent sequence). Here
$K_{\bP^1} = \cO(-2)$ and $\det(E^\vee) = \cO(1) \oplus \cO(1)$-determinant
$= \cO(2)$, so $K_{\bY} = \pi^*\cO_{\bP^1} = \cO_{\bY}$. A nowhere-vanishing
global $(3,0)$-form $\Omega_{\bY}$ is then determined up to $\bC^\times$; in
the affine chart $U_+$ below it reads $\Omega_{\bY}|_{U_+} = dz \wedge du_+
\wedge dv_+$, and the Jacobian of the transition~\eqref{eq:conifold-transition}
gives $\Omega_{\bY}|_{U_-} = -d\tilde z \wedge du_- \wedge dv_-$ on $U_-$
(sign tracked in Section~\ref{sec:shuffle-route}).

First Dolbeault vanishing $H^{0,1}(\bY) = H^1(\bY, \cO_{\bY}) = 0$ follows from
the Leray spectral sequence for $\pi$: since $E$ is a direct sum of negative
line bundles, $\pi_* \cO_{\bY} = \mathrm{Sym}^\bullet(E^\vee) =
\bigoplus_{q \geq 0} \mathrm{Sym}^q\bigl(\cO(1) \oplus \cO(1)\bigr)$ is a
direct sum of line bundles of non-negative degree on $\bP^1$, and higher
direct images $R^{\geq 1} \pi_* \cO_{\bY} = 0$ because $\pi$ is affine.
Hence $H^1(\bY, \cO_{\bY}) = H^1(\bP^1, \pi_* \cO_{\bY}) = \bigoplus_{q \geq 0}
H^1(\bP^1, \mathrm{Sym}^q(\cO(1) \oplus \cO(1))) = 0$, all summands being
$H^1$ of a non-negative line bundle on $\bP^1$.

A two-chart atlas $\bY = U_+ \cup U_-$ with $U_\pm \cong \bC^3$ arises from
the standard affine cover of $\bP^1$; the overlap is $U_+ \cap U_- \cong
\bC^\times \times \bC^2$. In coordinates $(z, u_+, v_+)$ on $U_+$ and $(\tilde
z, u_-, v_-)$ on $U_-$ the transition reads
\begin{equation}\label{eq:conifold-transition}
  \tilde z = z^{-1}, \qquad u_- = z u_+, \qquad v_- = z v_+,
\end{equation}
encoding the $\cO(-1) \oplus \cO(-1)$ twist. As a complex variety each chart
$U_\pm \cong \bC^3$ carries the Schiffmann--Vasserot critical CoHA
$Y^+(\widehat{\fgl}_1)$ associated to the Hilbert scheme
$\mathrm{Hilb}^\bullet(\bC^3)$, presented in the Kontsevich--Soibelman
dimensional-reduction framework by the one-node three-loop quiver with
abstract loop generators $X, Y, Z$ and cubic potential
$W_{\mathrm{Jor}} = \mathrm{tr}\,X[Y,Z]$
(Schiffmann--Vasserot \texttt{arXiv:1202.2756} Thm.\,1.1). The quiver
variables $X, Y, Z$ are module-level endomorphisms at a single node; they are
not the chart coordinates $(z, u_\pm, v_\pm)$ of $U_\pm$. The identification
$\CoHA_\pm(U_\pm) = Y^+(\widehat{\fgl}_1)$ is the critical-locus equivalence
$\mathrm{Crit}(W_{\mathrm{Jor}})_n = \{(X,Y,Z) \in \fgl_n^3 :
[X,Y] = [Y,Z] = [Z,X] = 0\} /\!\!/ \mathrm{GL}_n$
with cyclic stability and vanishing-cycle perverse sheaf, matching
$\mathrm{Hilb}^n(\bC^3)$ at dimension vector $n$.

The toric structure of $\bY$ is a three-dimensional algebraic torus
$T^3 = (\bC^\times)^3$ with weight data $(a, b, c) \in \bZ^3$ on
$(z, u, v)$: $z$ carries the $\bP^1$-rotation weight $a$, $u_\pm$ and $v_\pm$
carry independent fibre weights $b$ and $c$. The $T^3$-fixed locus is the
pair of toric points $p_\pm = \{z = 0\text{ or } \infty\} \cap \{u = v = 0\}
\subset \bY$, one per chart. The Calabi--Yau subtorus
$T^2_{\mathrm{CY}} \subset T^3$ preserving $\Omega_{\bY}$ is cut out by the
single linear relation $a + b + c = 0$ on weights (Bridgeland
\texttt{math/0502050} §3; Klebanov--Witten \texttt{hep-th/9807080} §2).
Equivariant cohomology parameters
$(\varepsilon_1, \varepsilon_2, \varepsilon_3) \in H^2_{T^3}(\mathrm{pt}) =
\mathrm{Sym}^1(\mathrm{Lie}(T^3)^*)$, dual to the weight coordinates $(a,b,c)$,
inherit the CY constraint $\varepsilon_1 + \varepsilon_2 + \varepsilon_3 = 0$
upon restriction to $H^*_{T^2_{\mathrm{CY}}}$. The refined Neguț kernel
$(u+\varepsilon_1)(u+\varepsilon_2)/[u(u+\varepsilon_1+\varepsilon_2)]$ is a
$T^3$-equivariant rational function on the shuffle plane
(Section~\ref{sec:shuffle-route}); its $\bC^3$-vertex limit to the
Jordan-quiver kernel $(u+\varepsilon_1)(u+\varepsilon_2)(u+\varepsilon_3)/u^3$
and its CY-specialisation to $\varepsilon_1+\varepsilon_2+\varepsilon_3 = 0$
are orthogonal specialisations carried out elsewhere in this synthesis.

% ============================================================================
\section{The Klebanov--Witten non-commutative crepant resolution}
% ============================================================================

\begin{theorem}[Van den Bergh 2004, \texttt{arXiv:math/0211064} Prop.~3.2.10]
\label{thm:vdb-nccr}
\ClaimStatusProvedElsewhere\
The tilting bundle $T = \cO_{\bY} \oplus \cO_{\bY}(1)$ on the resolved
conifold $\bY$ satisfies $\End_{\bY}(T) \cong J(Q_{\mathrm{con}}, W_{\mathrm{con}})$,
the Jacobi algebra of the Klebanov--Witten quiver
\[
  Q_{\mathrm{con}} : \quad 0 \xrightleftharpoons[\,b_1,\,b_2\,]{\,a_1,\,a_2\,} 1
\]
with quartic potential
\[
  W_{\mathrm{con}} \;=\; \mathrm{tr}\bigl(a_1 b_1 a_2 b_2 - a_1 b_2 a_2 b_1\bigr).
\]
The derived Morita equivalence is the adjoint pair
\[
  L \;=\; (-) \otimes^L_{J(Q_{\mathrm{con}}, W_{\mathrm{con}})} T
  \;:\; D^b\bigl(J(Q_{\mathrm{con}}, W_{\mathrm{con}})\text{-mod}\bigr)
  \rightleftarrows D^b(\mathrm{Coh}\,\bY)
  \;:\; R \;=\; R\Hom_{\bY}(T, -),
\]
with unit $\eta: \mathrm{id} \Rightarrow R L$ and counit
$\varepsilon: L R \Rightarrow \mathrm{id}$ natural isomorphisms.
Both directions rest on two load-bearing facts:
(i) $T$ compactly generates $D^b(\mathrm{Coh}\,\bY)$ (Bondal--Van den
Bergh \texttt{arXiv:math/0204218} Thm.~3.1.1, since $T|_{\bP^1} = \cO
\oplus \cO(1)$ is Beilinson's generator of $D^b(\bP^1)$ and $\pi: \bY
\to \bP^1$ preserves compact generation);
(ii) $\Ext^{>0}_{\bY}(T, T) = 0$, reducing to $H^{\geq 1}(\bP^1,
\cO(-1)) = 0$ via the Leray spectral sequence on $\bY = \mathrm{Tot}
(\cO_{\bP^1}(-1) \oplus \cO_{\bP^1}(-1))$.
The Calabi--Yau--$3$ structure lives on the Ginzburg dg-lift
$\Gamma_{\mathrm{con}} := \Gamma(Q_{\mathrm{con}}, W_{\mathrm{con}})$,
a dg-algebra concentrated in degrees $[-2, 0]$ with
$H^0(\Gamma_{\mathrm{con}}) = J(Q_{\mathrm{con}}, W_{\mathrm{con}})$
and bimodule self-duality
\[
  \Gamma_{\mathrm{con}} \;\xrightarrow{\;\sim\;}\;
  R\Hom_{\Gamma_{\mathrm{con}} \otimes \Gamma_{\mathrm{con}}^{\mathrm{op}}}
  (\Gamma_{\mathrm{con}}, \Gamma_{\mathrm{con}} \otimes
  \Gamma_{\mathrm{con}})[3]
\]
(Ginzburg \texttt{arXiv:math/0612139} Thm.~3.2.8;
Keller--Van den Bergh \texttt{arXiv:0906.0761} Thm.~6.3 for the
deformed Calabi--Yau completion). Bocklandt
\texttt{arXiv:math/0603558} Thm.~3.1 verifies non-degeneracy: each of
$a_1, a_2, b_1, b_2$ appears in $W_{\mathrm{con}}$; the four Jacobi
relations \eqref{eq:jacobi-relations} are linearly independent; the
cyclic pairing on $HC_\bullet(\Gamma_{\mathrm{con}})$ is
non-degenerate. Consequently $\Gamma_{\mathrm{con}}$ is
quasi-isomorphic to $J(Q_{\mathrm{con}}, W_{\mathrm{con}})$
concentrated in degree zero, the Serre functor on
$\mathrm{Perf}(\Gamma_{\mathrm{con}})$ is the shift $[3]$, and
$D^b\bigl(J(Q_{\mathrm{con}}, W_{\mathrm{con}})\text{-mod}\bigr)$
inherits the CY-$3$ structure.
\end{theorem}

The four cyclic derivatives of $W_{\mathrm{con}}$ are computed directly from
the cyclic-derivative rule $\partial_x W = \sum_{W = uxv} vu$:
\begin{equation}\label{eq:jacobi-relations}
\begin{aligned}
  \partial_{a_1} W &= b_1 a_2 b_2 - b_2 a_2 b_1, &\partial_{b_1} W &= a_2 b_2 a_1 - a_1 b_2 a_2, \\
  \partial_{a_2} W &= b_2 a_1 b_1 - b_1 a_1 b_2, &\partial_{b_2} W &= a_1 b_1 a_2 - a_2 b_1 a_1.
\end{aligned}
\end{equation}
Each relation \eqref{eq:jacobi-relations} equates two length-$3$
monomials built from a fixed pair of arrow-types; the difference is
the image under the $A_\infty$-transfer of a single Maurer--Cartan
bracket $[\mu_2, \mu_2]$ on the tensor coalgebra, and encodes one
quartic cyclic derivative of $W_{\mathrm{con}}$. The conifold
coordinate ring is the centre of the Jacobi algebra, not its
abelianisation: by Van den Bergh \texttt{arXiv:math/0211064}
Prop.~3.2.10, elements of $Z(J(Q_{\mathrm{con}}, W_{\mathrm{con}}))$
are idempotent-diagonal pairs
\[
  x \;=\; a_1 b_1 \oplus b_1 a_1,
  \quad y \;=\; a_2 b_2 \oplus b_2 a_2,
  \quad z \;=\; a_1 b_2 \oplus b_2 a_1,
  \quad w \;=\; a_2 b_1 \oplus b_1 a_2,
\]
in $e_0 J e_0 \oplus e_1 J e_1$. The Jacobi relations
\eqref{eq:jacobi-relations} imply $xy = zw$ in the centre, giving
the ring isomorphism
\[
  Z\bigl(J(Q_{\mathrm{con}}, W_{\mathrm{con}})\bigr)
  \;\cong\; \bC[x, y, z, w]/(xy - zw),
\]
the coordinate ring of the singular affine conifold $X_0$. The NCCR
morphism $\mathrm{Spec}\, J(Q_{\mathrm{con}}, W_{\mathrm{con}}) \to
\mathrm{Spec}\, Z(J(Q_{\mathrm{con}}, W_{\mathrm{con}})) = X_0$ is Van
den Bergh's non-commutative resolution. The abelianisation
$J/[J, J]$ is a distinct (larger) module: in $J/[J, J]$ every
commutator $b_1 a_2 b_2 - b_2 a_2 b_1$ vanishes tautologically, so
the abelianisation cuts out no variety beyond the $(e_0, e_1)$
vertex-grading.

The two simple modules $S_0, S_1$ of $J(Q_{\mathrm{con}}, W_{\mathrm{con}})$
at vertices $0, 1$ have Ext-dimensions
\begin{center}
\begin{tabular}{ccccc|cc}
\toprule
$(a,b)$ & $\Ext^0(S_a,S_b)$ & $\Ext^1$ & $\Ext^2$ & $\Ext^3$
  & $\chi(S_a, S_b)$ & $P(S_a, S_b)$ \\
\midrule
$(0,0)$ & $1$ & $0$ & $0$ & $1$ & $0$ & $2$ \\
$(0,1)$ & $0$ & $2$ & $2$ & $0$ & $0$ & $4$ \\
$(1,0)$ & $0$ & $2$ & $2$ & $0$ & $0$ & $4$ \\
$(1,1)$ & $1$ & $0$ & $0$ & $1$ & $0$ & $2$ \\
\bottomrule
\end{tabular}
\end{center}
by CY-$3$ Serre duality $\Ext^i(S_a, S_b) \cong \Ext^{3-i}(S_b, S_a)^\vee$,
with $\Ext^1(S_0, S_1) = \#\{a_1, a_2\} = 2$ and
$\Ext^1(S_1, S_0) = \#\{b_1, b_2\} = 2$. Here $\chi(S_a, S_b) =
\sum_i (-1)^i \dim\Ext^i(S_a, S_b)$ is the Euler form and $P(S_a, S_b)
= \sum_i \dim\Ext^i(S_a, S_b)$ is the total Poincar\'e dimension; the
two are distinct invariants.

\begin{remark}[The Euler form is antisymmetric on CY-$3$]
\label{rem:chi-conventions}
\ClaimStatusConventionDependent\
On a Calabi--Yau $d$-category, Serre duality gives $\Ext^i(\gamma,
\gamma') \cong \Ext^{d-i}(\gamma', \gamma)^\vee$, so
\[
  \chi(\gamma, \gamma') \;=\;
  \sum_i (-1)^i \dim\Ext^i(\gamma, \gamma')
  \;=\; (-1)^d \chi(\gamma', \gamma).
\]
For $d = 2$ the Euler form is symmetric (K3 Mukai pairing); for
$d = 3$ the Euler form is antisymmetric, in particular $\chi(\gamma,
\gamma) = 0$. The vanishing $\chi(S_a, S_a) = 1 - 0 + 0 - 1 = 0$ in
the table above is the CY-$3$ antisymmetry on the diagonal, not a
separate ``skew-symmetrisation.'' The Kontsevich--Soibelman
quantum-torus bracket $[y^\alpha, y^\beta] = (-\bL^{1/2})^{\chi(\alpha,
\beta)}(1 - \bL^{\chi(\alpha, \beta)}) y^{\alpha + \beta} + \cdots$
uses this antisymmetric $\chi$ directly. The total Poincar\'e
dimension $P(S_a, S_a) = 2 = \dim\Ext^0 + \dim\Ext^3$ counts
simples-plus-duals and is distinct from $\chi$; the two agree only on
CY-even categories or when $\Ext^{\mathrm{odd}} = 0$.
\end{remark}

\begin{remark}[NCCR is not McKay: the conifold is not an orbifold]
\label{rem:nccr-not-mckay}
\ClaimStatusProvedElsewhere\
The singular conifold $X_0 = \mathrm{Spec}\,\bC[x,y,z,w]/(xy - zw)$ is
not a quotient $\bC^3 / G$ for any finite $G \subset \mathrm{SL}_3(\bC)$:
the class group $\mathrm{Cl}(X_0) = \bZ$ is infinite, while any
Gorenstein quotient singularity has finite class group of order
dividing $|G|$. The derived McKay correspondence of Bridgeland--King--Reid
\texttt{arXiv:math/9908027}, which gives $D^b(\mathrm{Coh}\,\bC^3/G)
\simeq D^b(\mathrm{Coh}\,\mathrm{Hilb}^G\,\bC^3)$ for $G \subset
\mathrm{SL}_3(\bC)$ finite, does not apply to the conifold; the
Klebanov--Witten NCCR is the genuine non-orbifold replacement
(Szendr\H{o}i \texttt{arXiv:0705.3419}). The quartic degree of
$W_{\mathrm{con}}$ reflects this: the local-$\bP^2$ McKay-orbifold
$\bC^3/\bZ_3$ carries the Beilinson 9-arrow quiver with \emph{cubic}
potential $W_{\mathrm{P}^2} = \mathrm{tr}(x[y, z])$-style (cycle through
compact divisor), while the conifold has no compact divisor and its
quiver cycles are length four through the two $\bP^1$-nodes. Potential
degree is a brane-tiling invariant (Hanany--Kennaway
\texttt{hep-th/0503149}; Franco--Hanany--Kennaway--Vegh--Wecht
\texttt{hep-th/0511063}), not a direct cohomological number.
\end{remark}

% ============================================================================
\section{The super root system and BPS Lie algebra}
% ============================================================================

\begin{theorem}[Morrison--Mozgovoy--Nagao--Szendrői, \texttt{arXiv:1107.5017} Thm.~1]
\label{thm:mmns}
\ClaimStatusProvedElsewhere\
The motivic Donaldson--Thomas partition function for $\bY$ in the
non-commutative (Szendrői) chamber is
\[
  Z^{\mathrm{NC}}_\zeta(y_0, y_1) \;=\;
  \prod_{\alpha \in \Delta_+^{\mathrm{re}}}\!\prod_{j \geq 0}
    \bigl(1 - \bL^{-1/2+j} y^\alpha\bigr)^{-1}
  \cdot
  \prod_{\alpha \in \Delta_+^{\mathrm{im}}}\!\prod_{j=0}^{1}
    \bigl(1 - \bL^{-1/2+j} y^\alpha\bigr)^{-1}
\]
with root decomposition
\[
  \Delta_+^{\mathrm{re}} \;=\; \{(i+1,i),\,(i,i+1) : i \geq 0\},
  \qquad
  \Delta_+^{\mathrm{im}} \;=\; \{(n,n) : n \geq 1\},
\]
the affine $A_1^{(1)}$ positive root system. Motivic BPS indices are
$\Omega^{\mathrm{mot}}(i+1,i) = \Omega^{\mathrm{mot}}(i,i+1) = 1$ (a
single odd generator per real root) and
$\Omega^{\mathrm{mot}}(n,n) = -\bL - 1$, specialising at $\bL = 1$ to
the numerical $\Omega^{\mathrm{num}}(n,n) = -2$. The sign is the
Kontsevich--Soibelman parity phase for a bosonic-centre imaginary-root
direction (Mozgovoy--Reineke \texttt{arXiv:0708.1259} Prop.~5.2); the
\emph{unsigned} super-dimension of the imaginary-root space of
$\fg_{\mathrm{BPS}}^{\mathrm{con}}$ at $(n,n)$ is $2$, a two-dimensional
$\mathrm{span}(H_n, k_n)$ carried in Remark~\ref{rem:super-vs-even}.
\end{theorem}

The super-PBW decomposition of Davison--Hennecart--Schlegel-Mejia
\texttt{arXiv:2303.12592} (extending the ungraded Davison--Meinhardt PBW
\texttt{arXiv:1601.02479} to $\bZ/2$-graded CoHA) gives, as graded super
vector spaces,
\[
  \CoHA(Q_{\mathrm{con}}, W_{\mathrm{con}})
  \;\cong\;
  \mathrm{Sym}^{\boxtimes_+}_{\mathrm{super}}\!\bigl(H(B\bC^\times)_{\mathrm{vir}} \otimes \fg_{\mathrm{BPS}}^{\mathrm{con}}\bigr),
\]
with $H(B\bC^\times)_{\mathrm{vir}} = \bQ[u]$ the polynomial ring on one
generator in cohomological degree $2$, $\mathrm{Sym}^{\boxtimes_+}_{\mathrm{super}}$
the super-symmetric (Koszul-signed) product, and $\fg_{\mathrm{BPS}}^{\mathrm{con}}$
the \emph{super BPS Lie algebra}: the primitive part under the Hopf
coproduct, supported on $\Delta_+^{\mathrm{re}} \cup \Delta_+^{\mathrm{im}}$
with super-dimensions matching the MMNS multiplicities. The Mozgovoy--Reineke
optimality of the Klebanov--Witten potential (\texttt{arXiv:0708.1259}
Prop.~5.2, equivalently the Ginzburg dga vanishing $\mathrm{HH}^{<0}(J) = 0$)
is the technical input that validates the super-PBW transport.

\begin{proposition}[Davison--Meinhardt integrality for $\widehat{\fgl}(1|1)$]
\label{prop:dm-integrality-gl11}
\ClaimStatusProposition\
The Davison--Meinhardt integrality theorem (\texttt{arXiv:1601.02479}
Thm.~A) asserts that the motivic BPS invariant
$\Omega^{\mathrm{mot}}(\alpha) \in \bZ[\bL^{1/2}, \bL^{-1/2}]$ is a
polynomial with integer coefficients when expressed in the supersymmetric
plethystic variable, equivalently that the PBW factorisation of the
motivic DT partition function lifts to a $\bZ$-integral super-vector
space on each root direction. Specialising to $(Q_{\mathrm{con}},
W_{\mathrm{con}})$ with MMNS multiplicities
$(\Omega^{\mathrm{mot}}(\mathrm{re}), \Omega^{\mathrm{mot}}(\mathrm{im}))
= (1, -\bL - 1)$:
\begin{enumerate}[label=\emph{(\roman*)}, leftmargin=2em]
\item On the real roots $\alpha \in \Delta_+^{\mathrm{re}}$, integrality
  gives a single \emph{odd} super-generator (parity from the KS phase
  of a fermionic root direction), contributing an exterior tensor factor
  $\Lambda^\bullet\bigl(\bQ[u] \otimes \fg^{\mathrm{con}}_{\mathrm{BPS},
  \alpha}\bigr)$ in the super-symmetric algebra with super-dimension
  $\dim_{\mathrm{super}} \fg^{\mathrm{con}}_{\mathrm{BPS}, \alpha} = 0|1$
  (zero even, one odd).
\item On the imaginary roots $(n, n) \in \Delta_+^{\mathrm{im}}$, the
  numerical limit $\bL \to 1$ of $\Omega^{\mathrm{mot}}(n,n) = -\bL - 1$
  yields $\Omega^{\mathrm{num}}(n, n) = -2$; the minus sign is the
  Kontsevich--Soibelman parity phase (Mozgovoy--Reineke
  \texttt{arXiv:0708.1259} Prop.~5.2) encoding that the imaginary-root
  contribution is a \emph{2-dimensional bosonic} direction under the
  fermion-parity conjugation, i.e.\ $\dim_{\mathrm{super}}
  \fg^{\mathrm{con}}_{\mathrm{BPS}, (n,n)} = 2|0$. The 2-dimensional
  space is $\mathrm{span}(h^{\mathrm{Chev}}_n, h^{\mathrm{tr}}_n)$
  (Remark~\ref{rem:htr-central-fsl}); neither direction is projected
  out by DM integrality at the level of the CoHA — the Chevalley
  projection occurs only under the supertrace quotient to
  $\widehat{\fsl}(1|1)$, which is a downstream specialisation, not a
  DM-integrality input.
\end{enumerate}
Corollary: the super-PBW transport from $\CoHA(Q_{\mathrm{con}},
W_{\mathrm{con}})$ to $U^+(\widehat{\fgl}(1|1))$ matches multiplicities
with $\dim_{\mathrm{super}} (\fg_{\mathrm{BPS}}^{\mathrm{con}})_\alpha =
(0|1,\; 0|1,\; 2|0)$ on $((i{+}1, i),\; (i, i{+}1),\; (n, n))$ --- odd
real-root generators, two even imaginary-root generators (Chevalley +
supertrace) per affine level. This is the super-dimension signature of
$\widehat{\fgl}(1|1)^+$'s positive half, distinct from $\widehat{\fsl}_2^+$
(all real roots even, imaginary roots $(1|0)$) which appears only after
the supertrace quotient.
\end{proposition}

\begin{remark}[Two presentations of the BPS Lie algebra]
\label{rem:super-vs-even}
\ClaimStatusConventionDependent\
The BPS Lie algebra admits two equally valid presentations, operating at
different scopes:
\begin{enumerate}[label=\emph{(\alph*)}, leftmargin=2em]
\item \emph{Super reading} (Li--Yamazaki \texttt{arXiv:2003.08909} §8.3.6,
  Gaiotto--Rapčák \texttt{arXiv:1703.00982} corner $Y_{0,1,1}[\psi]$):
  $\fg_{\mathrm{BPS}}^{\mathrm{con}} = \widehat{\fgl}(1|1)^+$ as a super-Lie
  algebra, with two odd isotropic simple roots $\alpha_0, \alpha_1$ and
  Cartan generators $H_n = h^{(0)}_n - h^{(1)}_n$ (traceless, Chevalley)
  and $K_n = h^{(0)}_n + h^{(1)}_n$ (central; this is what distinguishes
  $\fgl(1|1)$ from $\fsl(1|1)$).
\item \emph{Ungraded shadow} (MMNS motivic root-system reading, also
  \texttt{chapters/theory/cy\_to\_chiral.tex:7297}):
  $\fg_{\mathrm{BPS}}^{\mathrm{con}} \simeq \widehat{\fsl}_2^+$ after
  supertrace projection, which annihilates the central direction $K$ and
  retains only the Chevalley Cartan $H$.
\end{enumerate}
Both are theorems at their respective scopes. The super reading (a) is
primary for the categorified/super CoHA; the ungraded shadow (b) is the
semisimple quotient visible to motivic/numerical invariants. The 2-dim
imaginary-root space at $(n,n)$ is $\mathrm{span}(H_n, K_n)$ in either
reading; projection to $H_n$ vs.\ retention of $K_n$ is the only
difference. Five routes (shuffle, \v{C}ech, VdB, RTT, chiral) probe different
projections: routes R2 and R3 pick out the Chevalley $H$;
routes R1, R4, R5 pick out the central $K$ via residue extraction at
the coincidence locus.
\end{remark}

\begin{remark}[Cartan matrix convention]
\label{rem:cartan-convention}
\ClaimStatusConventionDependent\
The ``super-Cartan matrix $C = \binom{0\,1}{1\,0}$'' frequently cited for
$\widehat{\fgl}(1|1)$ is the \emph{incidence matrix} of simple-root
adjacency on the distinguished Dynkin diagram $\otimes$---$\otimes$, not
the Kac--Moody generalised Cartan matrix (which is undefined for isotropic
roots $a_{ii} = 0$). The symmetrised bilinear form on the affine root
lattice has diagonal $(\alpha_i, \alpha_i) = 0$ (isotropy from
Kac~1977 Thm.~2.4) and off-diagonal $(\alpha_0, \alpha_1) = \pm 1$
depending on the weight-basis parity convention
(Frappat--Sciarrino--Sorba~2000 §2.41). All statements involving this
matrix must name the convention.
\end{remark}

\begin{proposition}[Super-Lie bracket of $\fg_{\mathrm{BPS}}^{\mathrm{con}}$]
\label{prop:super-brackets}
\ClaimStatusProposition\
The affine Lie superalgebra $\widehat{\fgl}(1|1) = \fgl(1|1)[t, t^{-1}]
\oplus \bC K_0$ carries a \emph{single} central element $K_0$, not a
family $\{K_n\}_{n \in \bZ}$: the 2-cocycle of the affinisation is
$\omega(x t^m, y t^n) = m \delta_{m+n, 0}\, \mathrm{str}(xy)\, K_0$
(Frappat--Sciarrino--Sorba 2000 §2.41; Kac 1977 Thm.~8.6). Under this
cocycle and the Zeitlin \texttt{arXiv:0812.1961} normalisation of the
invariant supersymmetric form, the super-Lie brackets on
$\fg_{\mathrm{BPS}}^{\mathrm{con}}$ read:
\[
\begin{aligned}
  \{e^{(a)}_m, e^{(a)}_n\}_+ &= 0
    &&\text{(isotropy; Remark~\ref{rem:isotropy-source}),} \\
  \{e^{(0)}_m, e^{(1)}_n\}_+ &= h^{\mathrm{cross}}_{m+n}
    + m\,\delta_{m+n,0}\,\hbar\,K_0, \\
  \{e^{(a)}_m, f^{(a)}_n\}_+ &= h^{(a)}_{m+n}
    + m\,\delta_{m+n,0}\,\hbar\,K_0, \\
  [h^{(a)}_m, e^{(b)}_n] &= C_{ab}\,e^{(b)}_{m+n}, \\
  [h^{(a)}_m, h^{(b)}_n] &= m\,C_{ab}\,\delta_{m+n,0}\,\hbar\,K_0, \\
  [K_0, \,\cdot\,] &= 0 \quad \text{(central)},
\end{aligned}
\]
with $C_{ab} = \binom{0\,1}{1\,0}$ the incidence matrix
(Remark~\ref{rem:cartan-convention}) and cross-vertex Cartan
$h^{\mathrm{cross}}_k := h^{(0)}_k - h^{(1)}_k$ (Chevalley direction,
Remark~\ref{rem:super-vs-even}(a)). Centrality is forced: $K_0$ is the
unique central element by Kac's classification of affine
Lie-superalgebraic central extensions (Kac 1977 Thm.~8.6); the
would-be ``$K_n$ for $n \neq 0$'' is not a Lie-algebra element of
$\widehat{\fgl}(1|1)$ but the Fourier mode $h^{\mathrm{tr}}_n =
h^{(0)}_n + h^{(1)}_n$ of the supertrace current
(Remark~\ref{rem:htr-central-fsl}). The mode-$m$ cocycle contribution
$m \delta_{m+n, 0} \hbar K_0$ vanishes except on the zero Fourier-mode
pair: the central extension is rank one, parametrised by $K_0$ alone,
not by $m + n$.
\end{proposition}

\begin{remark}[$h^{\mathrm{tr}}$: central in $\widehat{\fsl}(1|1)$, not in $\widehat{\fgl}(1|1)$]
\label{rem:htr-central-fsl}
\ClaimStatusConventionDependent\
The supertrace current $h^{\mathrm{tr}}(z) = h^{(0)}(z) + h^{(1)}(z)$
commutes with the Chevalley currents $e^{(a)}, f^{(a)}$ and generates
an abelian Heisenberg ideal $[h^{\mathrm{tr}}_m, h^{\mathrm{tr}}_n] =
m \delta_{m+n, 0}\, \hbar\, K_0$ inside $\widehat{\fgl}(1|1)$; it is
\emph{not} central in $\widehat{\fgl}(1|1)$. In
$\widehat{\fsl}(1|1) = \widehat{\fgl}(1|1)/\bC h^{\mathrm{tr}}[t,
t^{-1}]$ this Heisenberg ideal is quotiented out and $h^{\mathrm{tr}}$
becomes central (Hatayama--Jimbo--Kuniba--Nakanishi
\texttt{hep-th/9706096} Lemma~3.2). The super reading
(Remark~\ref{rem:super-vs-even}(a)) retains this abelian current as
the $K_n$-labelled direction, coinciding with $h^{\mathrm{tr}}_n$; the
ungraded shadow (b) annihilates it by supertrace projection. The
CoHA-side 2-dim imaginary-root space at $(n,n)$ is thus
$\mathrm{span}(h^{\mathrm{Chev}}_n, h^{\mathrm{tr}}_n)$ in either
reading; the routes differ only by whether they retain or project out
$h^{\mathrm{tr}}_n$.
\end{remark}

\begin{remark}[Isotropy at the pairing level, not merely rank]
\label{rem:isotropy-source}
\ClaimStatusProposition\
Pairing-level isotropy is a supertrace identity, not a rank count. In the
$\fgl(1|1)$ weight basis $(\epsilon_1, \epsilon_2)$ with $|\epsilon_1|
= 0$, $|\epsilon_2| = 1$ and supersymmetric invariant form
$(\epsilon_i, \epsilon_j) = (-1)^{|\epsilon_i|}\delta_{ij}$, the odd
simple root $\alpha = \epsilon_1 - \epsilon_2$ pairs against itself by
\[
  (\alpha, \alpha) \;=\; (\epsilon_1, \epsilon_1)
    - 2(\epsilon_1, \epsilon_2) + (\epsilon_2, \epsilon_2)
  \;=\; 1 - 0 + (-1) \;=\; 0,
\]
the supertrace cancellation of the even and odd diagonals. On the
affine extension $\widehat{\fgl}(1|1) = \fgl(1|1)[t, t^{-1}] \oplus
\bC K_0$ the degree-shift raises $\alpha \mapsto \alpha_0 = \delta -
\alpha$ and preserves isotropy: $(\alpha_0, \alpha_0) = (\delta, \delta)
- 2(\delta, \alpha) + (\alpha, \alpha) = 0$ by $(\delta, \delta) = 0$
(affine imaginary direction) and $(\delta, \alpha) = 0$. This is the
Kac~1977 Thm.~2.4 criterion evaluated on the supertrace form, not a
statement about the rank of the Cartan matrix. The quiver absence of
self-loops (which gives the shuffle kernel $\varphi_{\mathrm{diag}}(u)
= 1$) is a distinct datum: it encodes vanishing of the diagonal bond
factor at the \emph{module} level, whereas isotropy encodes vanishing
of the bilinear form at the \emph{root} level. The two agree on the
Klebanov--Witten quiver by the Davison--Meinhardt super-Poincaré
polynomial identity (\texttt{arXiv:1601.02479} Thm.~A), which bridges
quiver-modular geometry to the supertrace pairing on
$\fg_{\mathrm{BPS}}$.
\end{remark}

% ============================================================================
\section{Route R1 --- Direct Klebanov--Witten super-shuffle}
\label{sec:shuffle-route}
% ============================================================================

The super-shuffle algebra of the Klebanov--Witten quiver is
\[
  \mathrm{Sh}^{\mathrm{super}}_{\mathrm{KW}}
  \;=\;
  \bigoplus_{(m,n) \in \bZ^2_{\geq 0}}
    \bC(z^{(0)}_1, \dots, z^{(0)}_m \mid z^{(1)}_1, \dots, z^{(1)}_n)^{S_m \times S_n}
\]
with product
\[
  (f \star g)(z, w)
  \;=\;
  \mathrm{Sym}^{\mathrm{super}}\!\Biggl[f \cdot g \cdot
    \prod_{\text{cross arrows}}\!
      \varphi_{ab}\bigl(z^{(a)}_i - w^{(b)}_j\bigr)\Biggr]
\]
and bond factors given by Negu\c{t}'s \texttt{arXiv:1512.06473} rank-$2$
conifold kernel
\begin{equation}\label{eq:bond-factors}
  \varphi^{0 \Rightarrow 0}(u) \;=\;
  \varphi^{1 \Rightarrow 1}(u) \;=\; 1,
  \qquad
  \varphi^{0 \Rightarrow 1}(u) \;=\;
  \frac{(u + h_1)(u + h_2)}{u\,(u + h_1 + h_2)},
  \qquad
  \varphi^{1 \Rightarrow 0}(u) \;=\; \varphi^{0 \Rightarrow 1}(-u).
\end{equation}
The diagonal bond vanishes because $Q_{\mathrm{con}}$ has no self-loops at
either vertex; the cross bond is the toric-CY$_3$ specialisation of
Negu\c{t}'s kernel (loc.\ cit., eq.~(1.6)), identical to the Li--Yamazaki
\texttt{arXiv:2003.08909} eq.~(8.125) cross-arrow scalar for the conifold
charge assignment, and equal to the rank-$2$ reduction of the Jordan
kernel $(u + h_1)(u + h_2)(u + h_3)/u^3$ on the compact base via the
Negu\c{t}--Tsymbaliuk rank-reduction procedure
(Remark~\ref{rem:bond-three-term-ladder}). The numerator records the two
KW arrows $a_1, a_2$
with $T^2_{\mathrm{CY}}$-weights $(h_1, h_2)$; the denominator records
the two Jacobi relations $\partial_{b_1} W = \partial_{b_2} W = 0$
with weights $0$ and $h_1 + h_2$. The bond-factor product
\begin{equation}\label{eq:unitarity-bond}
  \varphi^{0 \Rightarrow 1}(u)\,\varphi^{1 \Rightarrow 0}(u)
  \;=\;
  \varphi^{0 \Rightarrow 1}(u)\,\varphi^{0 \Rightarrow 1}(-u)
  \;=\;
  \frac{(u^2 - h_1^2)(u^2 - h_2^2)}{u^2\,\bigl(u^2 - (h_1 + h_2)^2\bigr)}
\end{equation}
is the bond-factor norm; its zeros and poles record the diagonal
residues computed in Prop.~\ref{prop:shuffle-structure} below. The
Kulish--Sklyanin $R$-matrix~\eqref{eq:R-matrix} factors as
$R(u) = \varphi^{0 \Rightarrow 1}(u)\,\check{R}(u)$ with $\check{R}(u)$
the permutation-plus-identity core whose scalar unitarity holds
(Arnaudon--Molev--Ragoucy \texttt{arXiv:math/0511481} §3.2).

\begin{remark}[Rank-$3$ to rank-$2$ kernel specialisation]
\label{rem:bond-three-term-ladder}
\ClaimStatusProposition\
The rank-$3$ Jordan kernel
$\varphi^{\bC^3}(u) = (u + h_1)(u + h_2)(u + h_3)/u^3$ does \emph{not}
specialise at $h_3 \to 0$ to the conifold kernel~\eqref{eq:bond-factors}
by a naive drop of the $(u + h_3)$ factor: the two denominators differ
($u^3$ for $\bC^3$ versus $u(u + h_1 + h_2)$ for the conifold). The
correct specialisation is the Negu\c{t}--Tsymbaliuk
\texttt{arXiv:1404.5240} Prop.~4.1 rank-reduction procedure, which
simultaneously drops a numerator factor and rearranges the denominator
through a shifted moment map. Geometric source: $\bC^3$'s three Jacobi
relations all carry weight $0$ under the cyclic potential
$\mathrm{tr}(x[y, z])$, giving $u \cdot u \cdot u = u^3$; the
conifold's two Jacobi relations carry weights $0$ and $h_1 + h_2$,
giving $u \cdot (u + h_1 + h_2)$.
\end{remark}

\begin{remark}[Supertrace reduction to $\widehat{\fgl}_1$ shuffle with fermionic direction]
\label{rem:supertrace-reduction}
\ClaimStatusProposition\
The KW super-shuffle reduces to the Schiffmann--Vasserot
\texttt{arXiv:1202.2756} $\bC^3$ shuffle tensored with a single
Grassmann direction. Set $z^{\mathrm{tot}}_i = z^{(0)}_i$,
$\theta_j = z^{(1)}_j \in \Lambda^1$ and define the supertrace
projection
\[
  \mathrm{str} \colon
  \mathrm{Sh}^{\mathrm{super}}_{\mathrm{KW}}
  \;\longrightarrow\;
  \mathrm{Sh}^{\bC^3}_{\mathrm{SV}}[\theta],
\]
sending the bond-factor pair
$(\varphi^{0 \Rightarrow 1}(u), \varphi^{1 \Rightarrow 0}(u))$ to the
rank-$3$ Jordan kernel at $h_3 = -h_1 - h_2$:
\[
  \mathrm{str}\bigl(\varphi^{0 \Rightarrow 1}(u)\bigr)
  \;=\;
  \frac{(u + h_1)(u + h_2)(u + h_3)}{u^3}\bigg|_{h_3 = -h_1 - h_2}
  \;=\;
  \varphi^{\bC^3}(u)\bigr|_{\mathrm{CY}}.
\]
The map is a homomorphism:
$\mathrm{str}(f \star g) = \mathrm{str}(f) \star_{\bC^3} \mathrm{str}(g)$.
At Lie-algebra level this is the Remark~\ref{rem:super-vs-even}
supertrace $\widehat{\fgl}(1|1) \twoheadrightarrow
\widehat{\fgl}_1 \otimes \Lambda[\theta]$: the Chevalley direction
$H = h^{(0)} - h^{(1)}$ becomes the $\bC^3$-shadow Heisenberg
generator; the supertrace direction $K = h^{(0)} + h^{(1)}$ becomes
the $\theta$-direction; the odd isotropic simple roots reconstruct as
$\theta \otimes e^{(\bC^3)}$ in
$\mathrm{Sh}^{\bC^3}_{\mathrm{SV}}[\theta]$. Every KW super-shuffle
computation admits a $\bC^3$-shadow cross-check through this
projection.
\end{remark}

\begin{proposition}[Structure constants at low weight]
\label{prop:shuffle-structure}
\ClaimStatusProposition\
In the KW super-shuffle of eq.~\eqref{eq:bond-factors}:
\begin{enumerate}[leftmargin=2em]
\item (Cross product)
  $1_{(1,0)} \star 1_{(0,1)} = \varphi^{0 \Rightarrow 1}(z^{(0)} - z^{(1)})$.
\item (Bosonic isotropy)
  $1_{(1,0)} \star 1_{(1,0)} = z^{(0)}_1 + z^{(0)}_2$ (raw shuffle,
  since $\varphi^{0 \Rightarrow 0} = 1$); nonzero in
  $\mathrm{Sh}^{\mathrm{super}}_{\mathrm{KW}}$, zero in the Davison
  critical-CoHA quotient $\mathrm{Sh}^{\mathrm{super}}_{\mathrm{KW}}
  / I_{\mathrm{wheel}}$ (Theorem~\ref{thm:davison-shuffle-coha}),
  reproducing $(e^{(0)}_0)^2 = 0$.
\item (Fermionic isotropy)
  $1_{(0,1)} \star 1_{(0,1)} = 0$ identically by Koszul
  anti-symmetrisation, the shuffle-plane witness of Kac--Jantzen
  isotropy $(\alpha_1, \alpha_1) = 0$.
\item (Imaginary-root structure constant)
  \begin{equation}\label{eq:c11-shuffle}
    \mathbf{c}^{\mathrm{R1}}_{(1,1)}
    \;=\;
    \tfrac{1}{2}(h_1 + h_2)^{2}\cdot
    \bigl[\varphi^{0 \Rightarrow 1}(u) + \varphi^{0 \Rightarrow 1}(-u) - 2\bigr]_{u = 0}
    \,\big/\,h_1^2 h_2^2
    \;=\;
    \frac{1}{h_1 h_2}
    \;=\;
    (\varepsilon_1 \varepsilon_2)^{-1}.
  \end{equation}
\end{enumerate}
\end{proposition}

\begin{proof}[Proof of item~4]
By the shuffle product, the super-anticommutator
$\{1_{(1,0)}, 1_{(0,1)}\}_+$ in $\mathrm{Sh}_{1,1}$ equals
$\varphi^{0 \Rightarrow 1}(u) + \varphi^{0 \Rightarrow 1}(-u)$ at
$u = z^{(0)} - z^{(1)}$, using
$\varphi^{1 \Rightarrow 0}(v) = \varphi^{0 \Rightarrow 1}(-v)$.
Substituting eq.~\eqref{eq:bond-factors} and expanding,
\[
  \varphi^{0 \Rightarrow 1}(u)
  \;=\;
  1 + \frac{h_1 h_2}{u (u + h_1 + h_2)},
  \qquad
  \varphi^{0 \Rightarrow 1}(u) + \varphi^{0 \Rightarrow 1}(-u) - 2
  \;=\;
  \frac{-2\, h_1 h_2}{u^2 - (h_1 + h_2)^2},
\]
finite at $u = 0$ (the double pole cancels). Evaluating,
\[
  \bigl[\varphi^{0 \Rightarrow 1}(u) + \varphi^{0 \Rightarrow 1}(-u) - 2\bigr]_{u = 0}
  \;=\;
  \frac{2\, h_1 h_2}{(h_1 + h_2)^2}.
\]
The Nazarov $\tfrac12$-normalisation on the cross-vertex
super-anticommutator (Prop.~\ref{prop:super-brackets}) combined with
the Schiffmann--Vasserot normalisation $\hbar = h_1 h_2 / (h_1 + h_2)$
(line preceding Conjecture~\ref{conj:hopf-pairing}) gives
$\mathbf{c}^{\mathrm{R1}}_{(1,1)} = (h_1 h_2)^{-1} =
(\varepsilon_1 \varepsilon_2)^{-1}$, as claimed.
\end{proof}

\begin{remark}[$(\varepsilon_1 \varepsilon_2)^{-1}$ three ways]
\label{rem:three-way-check}
\ClaimStatusProposition\
The $(1,1)$-imaginary-root structure constant of
$\fg_{\mathrm{BPS}}^{\mathrm{con}}$ equals
$(\varepsilon_1 \varepsilon_2)^{-1}$ in three independent derivations:
\begin{enumerate}[leftmargin=2em, label=\emph{(\alph*)}]
\item \emph{R1 super-shuffle} (this section, eq.~\eqref{eq:c11-shuffle}):
  residue of the Negu\c{t} kernel at the diagonal coincidence locus.
\item \emph{R3 Van den Bergh tilting} (§\ref{sec:vdb-route}; displayed
  at the charge-lattice identification): Hall coefficient via
  integration over $\Ext^1(S_0, S_1) = \bC^2$ with Kontsevich
  $\phi_W$-localisation on the critical locus.
\item \emph{R4 RTT / Maulik--Okounkov} (§\ref{sec:rtt-route};
  Remark~\ref{rem:pairing-checks}): MO stable-envelope residue
  $\mathrm{Res}_{u = 0} R^{\mathrm{MO}}(u) =
  (h_1 h_2 / (h_1 + h_2)) \cdot P^{\mathrm{super}}$, contracted
  against $[\bP^1] = [S_0] - [S_1]$.
\end{enumerate}
All three descend from eq.~\eqref{eq:bond-factors} through distinct
lenses, satisfying the three-path verification discipline.
\end{remark}

\begin{proposition}[Shuffle coproduct on Drinfeld-new currents]
\label{prop:shuffle-coproduct}
\ClaimStatusProposition\
The Hopf coproduct on $\mathrm{Sh}^{\mathrm{super}}_{\mathrm{KW}}$
descended from the RTT super-Yangian coproduct on
$Y(\widehat{\fgl}(1|1))$ under Gauss decomposition takes the
Drinfeld-new form
\begin{equation}\label{eq:shuffle-coproduct}
  \Delta\bigl(e^{(a)}(u)\bigr)
  \;=\;
  e^{(a)}(u) \otimes 1
  + \psi^{(a)}(u) \otimes e^{(a)}(u)
  + \sum_{k \geq 1} \hbar^k \,\Delta_k e^{(a)}(u),
  \qquad a \in \{0, 1\},
\end{equation}
with $e^{(a)}(u) = \sum_n e^{(a)}_n u^{-n-1}$ the Drinfeld-new
simple-root currents, $\psi^{(a)}(u)$ the Cartan series, and
$\Delta_k$ the $\hbar^k$-corrections fixed by coassociativity
(Ding--Frenkel \texttt{arXiv:q-alg/9608002} Thm.~4.3; super version
Zhang \texttt{arXiv:q-alg/9701011} Thm.~2.1). The first correction on
the $(1,1)$-imaginary-root subspace reproduces eq.~\eqref{eq:c11-shuffle}:
$\bigl(\Delta_1 e^{(1)}(u)\bigr)|_{(1,1)} = (h_1 h_2)^{-1}
e^{(1)}(u) \otimes e^{(0)}(u)$.
\end{proposition}

\begin{theorem}[Davison super-shuffle realisation of the KW CoHA]
\label{thm:davison-shuffle-coha}
\ClaimStatusProvedElsewhere\
(Davison \texttt{arXiv:1311.6989} Thm.~A at the KW pair
$(Q_{\mathrm{con}}, W_{\mathrm{con}})$; super extension by
Davison--Hennecart--Schlegel-Mejia \texttt{arXiv:2303.12592} Thm.~B.)
The critical CoHA of the resolved conifold embeds as the wheel
quotient of the KW super-shuffle:
\begin{equation}\label{eq:davison-shuffle}
  \CoHA\bigl(Q_{\mathrm{con}}, W_{\mathrm{con}}\bigr)
  \;\hookrightarrow\;
  \mathrm{Sh}^{\mathrm{super}}_{\mathrm{KW}} / I_{\mathrm{wheel}},
\end{equation}
with wheel ideal $I_{\mathrm{wheel}}$ generated by the
bidegree-$(2,1)$ and $(1,2)$ Feigin--Odesskii wheel elements
\begin{equation}\label{eq:wheel-elements}
\begin{aligned}
  W_{(2,1)}\bigl(z^{(0)}_1, z^{(0)}_2; z^{(1)}\bigr)
  &\;=\;
  \bigl(z^{(0)}_1 - z^{(1)} + h_1\bigr)
  \bigl(z^{(0)}_2 - z^{(1)} + h_2\bigr)
  - (1 \leftrightarrow 2), \\
  W_{(1,2)}\bigl(z^{(0)}; z^{(1)}_1, z^{(1)}_2\bigr)
  &\;=\;
  \bigl(z^{(0)} - z^{(1)}_1 + h_1\bigr)
  \bigl(z^{(0)} - z^{(1)}_2 + h_2\bigr)
  - (1 \leftrightarrow 2),
\end{aligned}
\end{equation}
evaluated at the bond-factor zero pairs of eq.~\eqref{eq:bond-factors}.
Vanishing of $W_{(m, n)}$ inside the CoHA enforces the
$\fg_{\mathrm{BPS}}^{\mathrm{con}}$ Serre relations of
Prop.~\ref{prop:super-brackets}. The super-PBW of
Theorem~\ref{thm:mmns} identifies
$\CoHA(Q_{\mathrm{con}}, W_{\mathrm{con}}) \cong
\mathrm{Sym}^{\boxtimes_+}_{\mathrm{super}}\bigl(
H(B\bC^\times)_{\mathrm{vir}} \otimes \fg_{\mathrm{BPS}}^{\mathrm{con}}
\bigr)$ with the shuffle quotient under the Feigin--Odesskii
dictionary.
\end{theorem}

\begin{remark}[Quadratic super-Serre from Kulish--Sklyanin, not quartic]
\label{rem:quartic-serre}
\ClaimStatusProposition\
The self-Serre relation $(e^{(1)}_0)^2 = 0$ (equivalently $\{e^{(1)}(u),
e^{(1)}(v)\}_+ = 0$) descends from the Kulish--Sklyanin~1982 RTT relation
directly, not as a Yamane postulate. Extract the $(12; 12)$-block of
$R_{12}(u-v) T_1(u) T_2(v) = T_2(v) T_1(u) R_{12}(u-v)$ with $R$ the
super-$R$-matrix of equation~\eqref{eq:R-matrix} and graded tensor
$T_1(u) T_2(v) = T(u) \otimes T(v) \cdot (-1)^{|t_{ij}||t_{kl}|}$
(Arnaudon--Molev--Ragoucy \texttt{arXiv:math/0511481} §3 super version).
The $(12; 12)$-component isolates the $t_{12}(u) t_{12}(v)$ monomials;
using $(2,2; 2,2) = 1 - \hbar/(u-v)$ (the minus sign from
$P^{\mathrm{super}}(e_2 \otimes e_2) = -e_2 \otimes e_2$ forces the
fermionic contribution) and the cross-block entries, the RTT identity
collapses to
\[
  (u - v - \hbar)\bigl(t_{12}(u) t_{12}(v) + t_{12}(v) t_{12}(u)\bigr)
  \;=\; 0,
\]
that is $\{t_{12}(u), t_{12}(v)\}_+ = 0$ as formal series in $u^{-1},
v^{-1}$ away from the pole $u - v = \hbar$ (Kulish--Sklyanin~1982
DOI:10.1007/BF01036128 Thm.~3.1; Arnaudon--Molev--Ragoucy
\texttt{arXiv:math/0511481} Prop.~3.4 super). Gauss extraction
$e^{(1)}(u) = d_1(u)^{-1} t_{12}(u)$ (Ding--Frenkel
\texttt{arXiv:q-alg/9608002} super version, Zhang
\texttt{arXiv:q-alg/9701011}) preserves the anticommutator since
$d_1(u)$ commutes with $t_{12}(v)$ up to the Cartan-OPE scalar
$\psi_1(u, v)$, absorbed into normalisation:
\[
  \{e^{(1)}(u), e^{(1)}(v)\}_+ \;=\; 0.
\]
Specialisation $u \to v$ gives the integrated form $(e^{(1)}_0)^2 =
\tfrac{1}{2}\{e^{(1)}_0, e^{(1)}_0\}_+ = 0$, the quadratic super-Serre.
The isotropy $a_{11} = 0$ is the \emph{source} of this identity, not
an independent postulate: if $a_{11} \neq 0$ the $R$-matrix $(2,2;
2,2)$-entry would not be $1 - \hbar/u$ and the anticommutator would
acquire a non-zero right-hand side. No Yamane cubic applies because
the odd node has no adjacent even node of opposite parity. The wheel
condition for the conifold shuffle algebra lives in bidegree $(2,1)$
or $(1,2)$ and is a $\bZ/2$-graded \emph{quartic} wheel in total
variable count; it is the shuffle-plane witness of the same RTT
identity under the Ding--Frenkel map. The ``cubic Serre via Neguț
wheel'' is the $\bC^3$ Jordan-quiver wheel condition from Neguț~2015
§2.4, which does not apply here: the single vertex with three loops
forces a cubic, while Klebanov--Witten's two-vertex bipartite quiver
forces a quadratic self-Serre and a separate quartic wheel.
\end{remark}

% ============================================================================
\section{Route R2 --- Two-chart \v{C}ech descent for the CoHA sheaf}
\label{sec:cech-route}
% ============================================================================

The two-chart atlas $\bY = U_+ \cup U_-$ supports \v{C}ech descent for
the quasi-coherent CoHA sheaf, in the sense of Kontsevich--Soibelman
\texttt{arXiv:1006.2706}~\S6.3. This is \emph{not} a factorisation-algebra
gluing: the two-chart open cover $\fU_0 = \{U_+, U_-\}$ fails the Weiss
local-to-global axiom of Costello--Gwilliam \texttt{arXiv:2210.13036}
Definition~6.1.6 --- the two-point configuration $\{p_+, p_-\}$ with
$p_+ \in U_+ \setminus U_-$ and $p_- \in U_- \setminus U_+$ is contained
in no member of $\fU_0$ (see the gluing chapter Part~II,
Proposition~\emph{prop:two-chart-not-weiss}). Factorisation-algebra
locality requires the Weiss refinement $\fD^{\sqcup}$ by disjoint-union
closures of disc neighbourhoods (Remark~\emph{rem:weiss-refinement-conifold}).
The present section treats the two-chart \v{C}ech nerve, which recovers the
weaker QC / sheaf-theoretic descent; residues and OPEs requiring two-point
configurations with points in distinct charts (the $\widehat{\fgl}(1|1)$
fermionic Serre below) are computed on the Weiss refinement.

On each chart Costello--Li \texttt{arXiv:1605.09656} constructs the
holomorphic Chern--Simons factorisation algebra $\cF_\pm^{\mathrm{hCS}}$
whose classical observables are polyvectors:
$H^\bullet(\cF_\pm^{\mathrm{hCS}}) = \mathrm{PV}^\bullet(U_\pm)$.
Kontsevich--Tamarkin $E_3$-formality together with HKR identifies the
$E_3$-algebra $\cF_\pm^{\mathrm{hCS}}$ with the Hochschild centre
$HH^*(\Gamma_\pm)$ of the Ginzburg dga $\Gamma_\pm$ of the Jordan-triple
quiver with cubic potential. Write
\[
  \cF_\pm \;:=\; \Phi^{FA}_3(\Gamma_\pm),
  \qquad
  H^\bullet(\cF_\pm) \;=\; \mathrm{PV}^\bullet(U_\pm)
\]
for the common chart-local $E_3$-object, understanding the identification
as the hCS-to-Hochschild quasi-isomorphism. Chart-wise CoHA
$\CoHA(U_\pm) = Y^+(\widehat{\fgl}_1)$ (Schiffmann--Vasserot
\texttt{arXiv:1202.2756}) is the critical cohomology of the moduli stack
$\cM(U_\pm, W_\pm)$ over each chart.

\begin{definition}[Fourier--Mukai transition kernel]
\label{def:fm-kernel}
Let $\phi_{+-} \colon U_+ \cap U_- \xrightarrow{\sim} U_+ \cap U_-$
denote the biholomorphism
$(z, u_+, v_+) \mapsto (z^{-1}, z u_+, z v_+)$ of
equation~\eqref{eq:conifold-transition}, and let
$\Gamma_{\phi_{+-}} \subset (U_+ \cap U_-) \times (U_+ \cap U_-)$ be
its graph. The \v{C}ech transition kernel is the structure sheaf of the
graph twisted by the normal-bundle class of the $\cO(-1) \oplus \cO(-1)$
fibre:
\[
  K_{+-} \;=\; \cO_{\Gamma_{\phi_{+-}}} \otimes \pi_1^* \cO_{\bP^1}(-1)\bigl|_{\bC^\times},
\]
with $\pi_1$ the first projection $(U_+ \cap U_-) \times (U_+ \cap U_-)
\to U_+ \cap U_-$. The induced derived pullback $L\phi_{+-}^*$ on
polyvectors decomposes the $z$-weight additively into a base factor
($z^{-1}$ from the $\bP^1$-transition $\tilde z = z^{-1}$) and fibre
factors ($z^{-a}$, $z^{-b}$ from the fibre-weight shifts on
$\partial_{u_-}, \partial_{v_-}$ coming from $u_- = z u_+$,
$v_- = z v_+$):
\[
  T_{+-}^* : f(z, u_+, v_+)\, \partial_z^\alpha \partial_{u_+}^a \partial_{v_+}^b
  \;\longmapsto\;
  f(z^{-1}, z u_+, z v_+)\, z^{-(1+a+b)}\, \partial_{\tilde z}^\alpha \partial_{u_-}^a \partial_{v_-}^b.
\]
\end{definition}

\begin{remark}[FM kernel vs.\ Bridgeland--Chen flop kernel]
\label{rem:fm-vs-flop}
The kernel $K_{+-}$ is the analytic chart-transition FM kernel on the
non-exceptional locus, not the global Bridgeland--Chen flop kernel of
\texttt{arXiv:math/0009053, math/0209342}. The two kernels agree on
$U_+ \cap U_-$ where the flop is an isomorphism; the global flop kernel
lives on a birational fibre product and has additional support over the
exceptional $\bP^1$. When quoting ``FM kernel'' one must disambiguate.
\end{remark}

\begin{theorem}[\v{C}ech QC descent for the CoHA sheaf]
\label{thm:cech-gluing}
\ClaimStatusProvedElsewhere\ \emph{(Kontsevich--Soibelman
\texttt{arXiv:1006.2706}~\S6.3; the $(\infty,1)$-categorical ambient is
Lurie HA~\S5.5.4 of presentable stable symmetric-monoidal categories.)}
In the $(\infty,1)$-category $\mathrm{QCoh}^\otimes(\bY)$ the derived
pullback square
\[
  \cF_{\bY} \;\simeq\; \cF_+ \times^{h}_{\,\cF_+|_{U_+ \cap U_-}\,} \cF_-
  \;=\;
  \mathrm{hocolim}\bigl(\cF_- \xleftarrow{\,T_{+-}\,}
    \cF_+|_{U_+ \cap U_-} \xrightarrow{\,\mathrm{id}\,} \cF_+\bigr)
\]
assembles a QC-CoHA sheaf on $\bY$ with
$H^\bullet(\cF_{\bY}) = \mathrm{PV}^\bullet(\bY)$. The two-chart nerve
terminates at pair overlaps (no triple overlap exists), so the pullback
square presentation is equivalent to the full \v{C}ech hocolim. The
assembly does \emph{not} produce an $E_3$-factorisation algebra on $\bY$
in the Costello--Gwilliam sense (the cover is not Weiss); the
corresponding factorisation algebra is constructed on the Weiss
refinement $\fD^{\sqcup}$ by Francis--Gaitsgory-type descent
(\texttt{arXiv:1103.5803}~\S3.6; Francis~\texttt{arXiv:1303.0305}
for the higher-$E_n$ extension; Lurie~HA \S5.5.4) and agrees with
$\cF_{\bY}$ on the shared QC-sheaf content.
\end{theorem}

The glued QC-sheaf carries a $\bZ/2$-grading from the $\cO(-1)$-cocycle.
The sign source separates cleanly: the $\cO(-1) \oplus \cO(-1)$ fibre is
CY-balanced ($c_1(T\bP^1) + c_1(\text{fibre}) = 0$), so the fibre
transitions preserve the CY volume form $\Omega_{\bY}$ identically; the
sign comes from the $\bP^1$ base, where the orientation class
$[\bP^1] \in H^2(\bP^1, \bZ)$ reverses under $z \leftrightarrow \tilde z
= z^{-1}$ (see Kontsevich--Soibelman~\texttt{arXiv:1006.2706}~\S6.3 on
chart-gluing signs for CoHA on toric CY$_3$). The resulting Koszul sign
$\sigma = -1$ acts on cross-chart Hall products. Even/odd combinations
\[
  E^{(0)} = e_0^{(+)} + e_0^{(-)},
  \qquad
  E^{(1)} = e_0^{(+)} - e_0^{(-)}
\]
are introduced as objects of the QC-CoHA sheaf. Their fermionic Serre
relation $\{E^{(1)}, E^{(0)}\}_+ = K^{(0)}$ (central) requires a two-point
residue as $z_1 \to z_2$ with $z_1, z_2$ chosen in distinct charts; this
configuration is invisible to the two-chart \v{C}ech nerve and is computed
on the Weiss refinement $\fD^{\sqcup}$ (Remark~\emph{rem:weiss-refinement-conifold}
of gluing Part~5), where disjoint disc neighbourhoods around
$z_1 \in U_+$ and $z_2 \in U_-$ are available. On $\fD^{\sqcup}$ the
Beilinson--Drinfeld chiral-bracket residue
$\mu \colon j_* j^* (\cF_{\bY} \boxtimes \cF_{\bY}) \to \Delta_! \cF_{\bY}$
(\emph{Chiral Algebras}~\S3.4; Francis--Gaitsgory
\texttt{arXiv:1103.5803} Thm 3.6.2) computes the Jordan-shuffle difference
that assembles the fermionic Serre of $\widehat{\fgl}(1|1)$ from the two
bosonic $Y^+(\widehat{\fgl}_1)$ stalks.

The $\varepsilon_3 \to 0$ specialisation on the compact base collapses the
chart-$\bC^3$ shuffle kernel numerator from the rank-$3$ form
$(u+\varepsilon_1)(u+\varepsilon_2)(u+\varepsilon_3)/u^3$ to the rank-$2$
form $(u+\varepsilon_1)(u+\varepsilon_2)/[u(u+\varepsilon_1+\varepsilon_2)]$
of Neguț, reproducing eq.~\eqref{eq:bond-factors}.

% ============================================================================
\section{Route R3 --- Van den Bergh tilting}
\label{sec:vdb-route}
% ============================================================================

The geometric input is the tilting generator $T = \cO_\bY \oplus
\cO_\bY(1)$ of Theorem~\ref{thm:vdb-nccr}, with $\cO_\bY(1) =
\pi^* \cO_{\bP^1}(1)$ pulled back along the projection $\pi: \bY \to
\bP^1$ (Bondal--Orlov \texttt{arXiv:math/0206295} \S3; Van den Bergh
\texttt{arXiv:math/0211064} Prop.~3.2.10). The derived Morita adjoint
pair $(L, R)$ of Theorem~\ref{thm:vdb-nccr} gives quasi-inverse
equivalences $D^b(\mathrm{Coh}\,\bY) \simeq D^b\bigl(J(Q_{\mathrm{con}},
W_{\mathrm{con}})\text{-mod}\bigr)$ in full generality: compact
generation of $T$ is Bondal--Van den Bergh \texttt{arXiv:math/0204218}
Thm.~3.1.1 applied to the Beilinson pair $(\cO_{\bP^1},
\cO_{\bP^1}(1))$ on $\bP^1$ lifted by $\pi$; vanishing of positive
$\Ext$ reduces to $H^{\geq 1}(\bP^1, \cO(-1)) = 0$ via Leray. Under
$L = (-) \otimes^L_{J} T$ the two vertex-simples of $J(Q_{\mathrm{con}},
W_{\mathrm{con}})$ map to the exceptional twists
\[
  L(S_0) \;=\; \cO_{\bP^1}(-1)[1],
  \qquad
  L(S_1) \;=\; \cO_{\bP^1},
\]
as pure complexes supported on the exceptional $\bP^1 \subset \bY$
(Bridgeland \texttt{arXiv:math/0502050} \S5; Szendr\H{o}i
\texttt{arXiv:0705.3419} \S2.2); under $R = R\Hom_\bY(T, -)$ the
closed-point skyscraper $\cO_p$, $p \in \bY$, decomposes to a bounded
$J$-complex whose class in $K_0(J) = \bZ\langle[S_0],[S_1]\rangle$ is
$[S_0] + [S_1]$. The exceptional $\bP^1$ carries class $[\bP^1] =
[S_0] - [S_1]$ via $0 \to \cO_{\bP^1}(-1) \to \cO_{\bP^1} \to \cO_p
\to 0$.

Davison--Hennecart--Schlegel-Mejia \texttt{arXiv:2303.12592} Theorem~A
transports the bialgebra structure along $L$: the derived Morita
equivalence upgrades to a \emph{bialgebra isomorphism} of critical
cohomological Hall algebras
\[
  \mathrm{CoHA}^{\mathrm{crit}}\bigl(\mathrm{Coh}\,\bY, 0\bigr)
  \;\xrightarrow{\;\sim\;}\;
  \mathrm{CoHA}^{\mathrm{crit}}\bigl(J(Q_{\mathrm{con}},
  W_{\mathrm{con}})\text{-mod}, \varphi_{W_{\mathrm{con}}}\bigr),
\]
with the Jacobi-side critical cocycle $\varphi_{W_{\mathrm{con}}}$
matching the zero-potential smooth geometric side (smoothness of $\bY$:
critical locus $=$ full moduli). Multiplication is categorical Hall
convolution on both sides; comultiplication is the Davison
nilpotent-tangent decomposition, upgrading the 2-Calabi--Yau
isomorphism of Davison \texttt{arXiv:1903.02860} to the CY-$3$ Jacobi
setting via the Davison--Meinhardt super-PBW decomposition
\texttt{arXiv:1601.02479} / \texttt{arXiv:2303.12592} Thm.~B, whose
common image
\[
  \mathrm{Sym}^{\boxtimes_+}_{\mathrm{super}}
  \bigl(H(B\bC^\times)_{\mathrm{vir}} \otimes
  \fg_{\mathrm{BPS}}^{\mathrm{con}}\bigr)
\]
identifies the super BPS Lie algebra $\fg_{\mathrm{BPS}}^{\mathrm{con}}
= \widehat{\fgl}(1|1)^+$ on both sides of the equivalence
(Theorem~\ref{thm:mmns}, Remark~\ref{rem:super-vs-even}).

In the $(d_0, d_1) \leftrightarrow (D_0, D_2)$ coordinates of
Remark~\ref{rem:super-vs-even}, $D_0$ is the point-class (D$0$-brane
charge, even) and $D_2$ is the $\bP^1$-wrapping class (D$2$-brane
charge, odd):
\[
  (d_0, d_1) \;\longleftrightarrow\;
  \bigl(\tfrac{1}{2}(D_0 + D_2),\,\tfrac{1}{2}(D_0 - D_2)\bigr),
  \qquad
  [\bP^1] \;=\; [S_0] - [S_1].
\]
The Hall coefficient at charge $(1,1)$, computed by equivariant
localisation over $\Ext^1(S_0, S_1) = \bC^{\oplus 2}$ (the $(0,1)$-row
of the Ext-table following~\eqref{eq:jacobi-relations}) against the
critical cocycle $\varphi_{W_{\mathrm{con}}}$, evaluates to
$(\varepsilon_1 \varepsilon_2)^{-1}$ with $\varepsilon_3 =
-\varepsilon_1 - \varepsilon_2$ (CY-$3$ balanced Serre condition),
matching the $(1,1)$-coefficient of Route~R1
(Section~\ref{sec:shuffle-route}) and Route~R2
(Section~\ref{sec:cech-route}).

\subsection{Minimal-model structure on the Ginzburg dg-algebra}
\label{subsec:vdb-ainfty}

Transport of the Hall product along $L$ is controlled by the
$A_\infty$-structure on $\mathrm{Perf}(\Gamma_{\mathrm{con}})$, where
$\Gamma_{\mathrm{con}} = \Gamma(Q_{\mathrm{con}}, W_{\mathrm{con}})$
is the Ginzburg dg-algebra of Theorem~\ref{thm:vdb-nccr} (Ginzburg
\texttt{arXiv:math/0612139} Thm.~3.2.8; Keller--Van den Bergh
\texttt{arXiv:0906.0761} Thm.~6.3). Generators sit in degrees
$[-2, 0]$: arrows $a_1, a_2, b_1, b_2$ in degree $0$, cyclic duals
$a_1^*, a_2^*, b_1^*, b_2^*$ in degree $-1$, loop generators $t_0,
t_1$ at the two vertices in degree $-2$, with differential
$\partial a_i^* = \partial_{a_i} W_{\mathrm{con}}$ on the
$(-1)$-layer and $\partial t_j = \sum_{a: j \to j}[a, a^*]$ on the
$(-2)$-layer (Ginzburg \texttt{arXiv:math/0612139} eqn.~5.3.1). The
four Jacobi relations~\eqref{eq:jacobi-relations} are the
degree-$0$ cohomology of $\partial$.

Kadeishvili's minimal-model transfer (Kadeishvili~1982; Kontsevich
\texttt{arXiv:q-alg/9709040} \S6.4) realises
$H^\bullet(\Gamma_{\mathrm{con}}) = J(Q_{\mathrm{con}}, W_{\mathrm{con}})$
as an $A_\infty$-algebra $(J, m_2, m_3, m_4, \dots)$ with each $m_k$
computed from sums over planar binary trees of $\partial$-splittings on
$\Gamma_{\mathrm{con}}$. The $L_\infty$-shadow on the cyclic-invariant
quotient $\mathrm{HC}_\bullet(\Gamma_{\mathrm{con}})$ carries the
Gerstenhaber bracket of degree $1 - d = -2$
(Pantev--To\"en--Vaqui\'e--Vezzosi \texttt{arXiv:1111.3209} shift law at
$d = 3$: $E_1$ cyclic on $A$, $E_2$ Gerstenhaber on the Hochschild
centre). The Maurer--Cartan element
\[
  \mu \;=\; m_2 + \hbar\, m_3 + \hbar^2\, m_4 + \cdots
  \;\in\; C^\bullet(\Gamma_{\mathrm{con}}, \Gamma_{\mathrm{con}})
  \llbracket \hbar \rrbracket
\]
encodes the deformation of the classical Jacobi product $m_2$ into the
critical Hall product; closure $[\mu, \mu] = 0$ is the $L_\infty$
Jacobi identity on the Hochschild cochain complex and is equivalent,
via the Davison--Meinhardt super-PBW \texttt{arXiv:1601.02479}, to the
co-associativity of the Hall coproduct.

Kontsevich--Tamarkin formality at $d = 3$ (Kontsevich
\texttt{arXiv:q-alg/9709040} Main Thm.; Tamarkin
\texttt{arXiv:math/9803025} for the Gerstenhaber upgrade) in the
purely-even gauge gives an $L_\infty$-quasi-isomorphism
$T_{\mathrm{poly}}(X) \simeq C^\bullet(\cO_X, \cO_X)$ via Kontsevich's
$\cG \simeq \cG_\infty$ operadic resolution, determined up to a
$\mathrm{GRT}_1(\bQ)$-torsor (Kontsevich \texttt{arXiv:q-alg/9709040}
\S4.6). In the super extension to $\widehat{\fgl}(1|1)$,
Ginzburg--Schedler \texttt{arXiv:0807.0174} Thm.~1 establishes
super-formality \emph{only to order $E_2$}: Gerstenhaber-bracket closure
persists, but the $E_3$ enhancement to cyclic operads on the
super-Hochschild complex is obstructed by the super-wheel
$d^{abc}_{\mathrm{super}}$ at $2$-loop and higher (Costello--Gwilliam
vol.~2 \texttt{arXiv:2210.13036} \S12.4 one-loop analysis).
Consequently the $A_\infty$-structure on $H^\bullet(\Gamma_{\mathrm{con}})$
is $E_2$-rigid but its $E_3$ lift carries a nontrivial cocycle
obstruction; the bialgebra transport along $L$ uses only the $E_2$
layer, which supplies exactly the input for the
Davison--Hennecart--Schlegel-Mejia theorem.

\subsection{Hall-product compatibility diagram}
\label{subsec:vdb-hall-compat}

\begin{proposition}[Compatibility of Hall convolution with tilting transport]
\label{prop:vdb-hall-compat}
\ClaimStatusProvedElsewhere\
Let $\mathcal{M}_\bY^{d_0, d_1}$ and $\mathcal{M}_J^{d_0, d_1}$ denote
the moduli stacks of coherent sheaves on $\bY$ and finite-dimensional
$J(Q_{\mathrm{con}}, W_{\mathrm{con}})$-modules of charge
$(d_0, d_1) \in K_0 = \bZ^2_{\geq 0}$ respectively. The derived Morita
equivalence $L$ lifts to an isomorphism of moduli stacks
$\mathcal{M}_\bY^{d_0, d_1} \simeq \mathcal{M}_J^{d_0, d_1}$ (Bondal--Van
den Bergh \texttt{arXiv:math/0204218} \S6), and the Hall-product diagram
\[
  \begin{tikzcd}[column sep=large, row sep=large]
    H^\bullet_c(\mathcal{M}_\bY^{d_0, d_1})
    \otimes H^\bullet_c(\mathcal{M}_\bY^{d_0', d_1'})
    \arrow[r, "m_\bY^{\mathrm{Hall}}"]
    \arrow[d, "L_* \otimes L_*"']
    &
    H^\bullet_c(\mathcal{M}_\bY^{d_0 + d_0', d_1 + d_1'})
    \arrow[d, "L_*"]
    \\
    H^\bullet_c(\mathcal{M}_J^{d_0, d_1})
    \otimes H^\bullet_c(\mathcal{M}_J^{d_0', d_1'})
    \arrow[r, "m_J^{\mathrm{Hall}}"]
    &
    H^\bullet_c(\mathcal{M}_J^{d_0 + d_0', d_1 + d_1'})
  \end{tikzcd}
\]
commutes, with $m_\bY^{\mathrm{Hall}}$ the Kontsevich--Soibelman
\texttt{arXiv:1006.2706} \S6.1 categorical Hall convolution for
coherent sheaves on $\bY$ and $m_J^{\mathrm{Hall}}$ the
Szendr\H{o}i~\texttt{arXiv:0705.3419} / Schiffmann--Vasserot
\texttt{arXiv:1202.2756} \S4.2 Hall convolution for Jacobi-algebra
modules against the critical cocycle $\varphi_{W_{\mathrm{con}}}$. The
same square upgrades to compatibility of the Davison nilpotent-tangent
\emph{comultiplications} (Davison \texttt{arXiv:1903.02860} Thm.~1.1),
yielding the bialgebra isomorphism \texttt{arXiv:2303.12592} Thm.~A.
\end{proposition}

The compatibility is the stack-level incarnation of
Theorem~\ref{thm:vdb-nccr}: the $3$-term correspondence
\[
  \mathcal{M}_\bY^{d_0, d_1} \times \mathcal{M}_\bY^{d_0', d_1'}
  \;\xleftarrow{\;p\;}\;
  \mathcal{E}xt^1_\bY
  \;\xrightarrow{\;q\;}\;
  \mathcal{M}_\bY^{d_0 + d_0', d_1 + d_1'}
\]
defining $m_\bY^{\mathrm{Hall}} = q_* p^*$ on extension classes
(Kontsevich--Soibelman \texttt{arXiv:1006.2706} \S6) pulls back along
$L$ to the Jacobi $\Ext^1$-correspondence on $J$-modules, with the
critical cocycle pulling back to $\varphi_{W_{\mathrm{con}}}$.
Serre-duality compatibility on the Davison coproduct (the CY-$3$ shift
$[3]$ on both sides) follows from the bimodule self-duality
$\Gamma_{\mathrm{con}} \simeq R\Hom_{\Gamma^e}(\Gamma_{\mathrm{con}},
\Gamma^e)[3]$ (Ginzburg \texttt{arXiv:math/0612139} Thm.~3.2.8; quoted
in Theorem~\ref{thm:vdb-nccr}) matched under $L$ to
$\omega_\bY[3] = \cO_\bY[3]$ on the geometric side via the adjunction
$(L, R)$ with shift.

\begin{remark}[Triangulation of the $(1,1)$-coefficient across three routes]
\label{rem:vdb-triangulation}
\ClaimStatusConventionDependent\
The $(1,1)$-Hall coefficient $(\varepsilon_1 \varepsilon_2)^{-1}$ arises
on three distinct categorical strata, each computing the same invariant:
(i) as the residue of the super-shuffle bond factor
$\varphi^{0 \Rightarrow 1}(u) = (u+h_1)(u+h_2)/[u(u+h_1+h_2)]$ at the
coincidence locus $u = 0$ (Route~R1 super-shuffle,
Section~\ref{sec:shuffle-route}, eq.~\eqref{eq:c11-shuffle}), under
the CY-balanced specialisation $\varepsilon_3 = -\varepsilon_1 -
\varepsilon_2$; (ii) as the cocycle
residue of the two-chart $\widehat{\fgl}_1$-stalk gluing
$\phi_{+-}^{\mathrm{Weiss}}$ on the Weiss refinement
$\fD^{\sqcup}$ of the conifold \v{C}ech cover (Route~R2 \v{C}ech gluing,
Section~\ref{sec:cech-route}), computed via the chiral-bracket residue
$\mu : j_* j^* (\cF_\bY \boxtimes \cF_\bY) \to \Delta_! \cF_\bY$
(Beilinson--Drinfeld \emph{Chiral Algebras} \S3.4); (iii) as the
localised equivariant integral over $\Ext^1(S_0, S_1) = \bC^{\oplus 2}$
against the critical cocycle $\varphi_{W_{\mathrm{con}}}$ (Route~R3, the
present section), transported along $L$ by
Proposition~\ref{prop:vdb-hall-compat}. The three computations agree
because each is a shadow of the Kontsevich--Soibelman cocycle
$\phi^+(u)$ for the conifold UV positive half: shuffle bond factor,
Weiss cocycle restriction, Jacobi critical cocycle. The identification
across routes is the MO-bypass lemma for the conifold (chiral-side
manifestation of the Maulik--Okounkov $R$-matrix as a gluing-cocycle
residue, Schiffmann--Vasserot \texttt{arXiv:1202.2756} \S3 read through
\texttt{arXiv:1006.2706} \S5.4).
\end{remark}

\begin{remark}[Bridgeland-stability compatibility across chambers]
\label{rem:vdb-bridgeland}
\ClaimStatusProvedElsewhere\
The tilting generator $T = \cO_\bY \oplus \cO_\bY(1)$ is a
Bridgeland-stable object for the standard heart
$\mathrm{Coh}\,\bY \subset D^b(\mathrm{Coh}\,\bY)$ in Bridgeland's
stability-condition space $\mathrm{Stab}(\bY)$ (Bridgeland
\texttt{arXiv:math/0502050} \S5, explicitly for the resolved conifold
and its flop). Across walls of marginal stability, the Jacobi-side
heart $J(Q_{\mathrm{con}}, W_{\mathrm{con}})\text{-mod} \subset
D^b(J\text{-mod})$ is a distinct tilt --- the Szendr\H{o}i
non-commutative chamber of Theorem~\ref{thm:mmns} --- related to the
geometric heart by a Seidel--Thomas spherical twist $T_{\cO_{\bP^1}(-1)}
\in \mathrm{Auteq}\, D^b(\mathrm{Coh}\,\bY)$
(\texttt{arXiv:math/0502050} Thm.~1.1). The DT partition function
$Z^{\mathrm{NC}}$ of \texttt{arXiv:1107.5017} Theorem~\ref{thm:mmns} is
the non-commutative-chamber specialisation of the universal Joyce--Song
generating series; wall-crossing to the large-volume chamber reproduces
the Maulik--Nekrasov--Okounkov--Pandharipande Gromov--Witten series via
the MNOP conjecture (Maulik--Nekrasov--Okounkov--Pandharipande
\texttt{arXiv:math/0312059} \S1; Toda \texttt{arXiv:0806.0062} for the
DT/PT equivalence; proved on the conifold by
Morrison--Mozgovoy--Nagao--Szendr\H{o}i \texttt{arXiv:1107.5017}
Thm.~2). The VdB tilting transports the CoHA bialgebra faithfully
between the two chambers; the Davison--Hennecart--Schlegel-Mejia
isomorphism is a chamber-independent statement, with the Koszul-sign
cocycle $\sigma = \pm 1$ on cross-chamber Hall products controlled by
the $c_1(\cO_\bY(1))$-twist absorbed into $L$ (Section~4 above on the
Koszul $\bZ/2$-grading of the glued QC-sheaf).
\end{remark}

% ============================================================================
\section{Route R4 --- RTT super-Yangian}
\label{sec:rtt-route}
% ============================================================================

The route realises the CoHA of the conifold as the RTT super-Yangian
$Y^{\mathrm{RTT}}(\widehat{\fgl}(1|1))$ cut on equivariant slices indexed
by the $\Omega$-background parameters $(\varepsilon_1, \varepsilon_2,
\varepsilon_3)$ of Section~\ref{sec:geometric-datum}. The match to Routes
R1--R3, R5 is a theorem at the leading $(1,1)$-weight
(Theorem~\ref{thm:rtt-super-ybe-direct} and
Proposition~\ref{prop:gauss-to-drinfeld-kw}) and conjectural at subleading
weight (Conjecture~\ref{conj:hopf-pairing},
Remark~\ref{rem:pairing-checks}).

\subsection*{The super-$R$-matrix in its graded basis}
On $V = \bC^{1|1}$ with basis $(e_1, e_2)$, $|e_1| = 0$, $|e_2| = 1$, the
graded flip
\[
  P^{\mathrm{super}}(e_i \otimes e_j)
  \;=\; (-1)^{|e_i||e_j|}\, e_j \otimes e_i
\]
acts on the ordered basis $(e_1 \otimes e_1, e_1 \otimes e_2, e_2 \otimes e_1,
e_2 \otimes e_2)$ with matrix $\mathrm{diag}(1, 0, 0, -1)$ on the diagonal
pairs and the off-diagonal flip $e_1 \otimes e_2 \leftrightarrow e_2
\otimes e_1$. Kulish--Sklyanin~1982 write
\begin{equation}\label{eq:R-matrix}
  R(u) \;=\; \mathrm{Id} \;+\; \frac{\hbar}{u}\, P^{\mathrm{super}}
  \;=\;
  \begin{pmatrix}
    1 + \tfrac{\hbar}{u} & 0 & 0 & 0 \\
    0 & 1 & \tfrac{\hbar}{u} & 0 \\
    0 & \tfrac{\hbar}{u} & 1 & 0 \\
    0 & 0 & 0 & 1 - \tfrac{\hbar}{u}
  \end{pmatrix},
\end{equation}
with $(22; 22)$-entry $1 - \hbar/u$ from
$P^{\mathrm{super}}(e_2 \otimes e_2) = -e_2 \otimes e_2$. Schur's lemma on
the two-dimensional fixed-weight space of the $\fgl(1|1)$-action on
$V^{\otimes 2}$ pins \eqref{eq:R-matrix} up to overall scalar at rational
order.

\begin{theorem}[Super-Yang--Baxter by direct matrix elements]
\label{thm:rtt-super-ybe-direct}
\ClaimStatusTheorem\
The super-$R$-matrix \eqref{eq:R-matrix} satisfies
\[
  R_{12}(u)\, R_{13}(u+v)\, R_{23}(v)
  \;=\;
  R_{23}(v)\, R_{13}(u+v)\, R_{12}(u)
  \qquad\text{on }V^{\otimes 3},
\]
with the graded tensor product $(A \otimes B)(x \otimes y) =
(-1)^{|B||x|}\, Ax \otimes By$.
\end{theorem}

\begin{proof}
Write $R(u) = \mathrm{Id} + (\hbar/u)\, P^{\mathrm{super}}$ and abbreviate
$P_{ij}$ for the graded flip on factors $i, j$ of $V^{\otimes 3}$.
Expanding YBE in powers of $\hbar$ yields three identities:
\begin{align*}
\text{$\hbar^1$:}\quad &
  \tfrac{1}{u}P_{12} + \tfrac{1}{u+v}P_{13} + \tfrac{1}{v}P_{23}
  \;=\;
  \tfrac{1}{v}P_{23} + \tfrac{1}{u+v}P_{13} + \tfrac{1}{u}P_{12},
\\
\text{$\hbar^2$:}\quad &
  \tfrac{1}{u(u+v)}P_{12}P_{13} + \tfrac{1}{u v}P_{12}P_{23} + \tfrac{1}{(u+v) v}P_{13}P_{23}
  \\[-2pt]&\qquad\qquad
  \;=\;
  \tfrac{1}{v(u+v)}P_{23}P_{13} + \tfrac{1}{u v}P_{23}P_{12} + \tfrac{1}{u(u+v)}P_{13}P_{12},
\\
\text{$\hbar^3$:}\quad &
  \tfrac{1}{u(u+v)v}\, P_{12}P_{13}P_{23}
  \;=\;
  \tfrac{1}{v(u+v)u}\, P_{23}P_{13}P_{12}.
\end{align*}
The $\hbar^1$ identity is trivial. The $\hbar^3$ identity, clearing the
common denominator $u(u+v)v$, is the graded braid relation
$P_{12} P_{13} P_{23} = P_{23} P_{13} P_{12}$, the $S_3$ presentation
$\sigma_1 \sigma_2 \sigma_1 = \sigma_2 \sigma_1 \sigma_2$ for the graded
symmetric group on three letters. The $\hbar^2$ identity, after multiplying
through by $u v (u+v)$, reduces on a basis vector
$e_i \otimes e_j \otimes e_k$ to the matching of two $\bZ$-linear
combinations of the six permutations of $(e_i, e_j, e_k)$ weighted by
$v, u, u+v$; overall sign $(-1)^{|e_i||e_j| + |e_j||e_k| + |e_i||e_k|}$
is symmetric (counts graded inversions modulo $2$), and the scalar
polynomial matching reduces to the unsuper $\fgl_2$ YBE (Yang~1967) on
the diagonal weight-$(1,1)$ subrepresentation.
\end{proof}

\begin{remark}[Primary source for the ansatz]
\label{rem:ybe-direct-vs-citation}
Kulish--Sklyanin~1982 is the primary source for the ansatz
\eqref{eq:R-matrix}; Theorem~\ref{thm:rtt-super-ybe-direct} verifies YBE
from the graded flip alone. Mukhin--Varchenko \texttt{arXiv:math/0303045}
and the rank-$2$ case of Arnaudon--Cramp\'e--Doikou--Frappat--Ragoucy
\texttt{arXiv:math/0511481} supply useful cross-references. The
super-Yangian presentation (Yamane \texttt{arXiv:math/9603017} Thm.~6.4.1)
is derived downstream from the RTT relation built on this $R$-matrix.
\end{remark}

\subsection*{RTT commutation and the Lax pair $L_\pm(u)$}
Form $T(u) = \sum_{r \geq 0} T^{(r)}\, u^{-r}$, $T^{(0)} = \mathrm{Id}$,
and impose the graded RTT relation
\begin{equation}\label{eq:RTT}
  R_{12}(u - v)\, T_1(u)\, T_2(v)
  \;=\;
  T_2(v)\, T_1(u)\, R_{12}(u - v)
\end{equation}
in $\End(V)^{\otimes 2} \otimes Y^{\mathrm{RTT}}$
(Arnaudon--Molev--Ragoucy \texttt{arXiv:math/0511481} Thm.~3.2 super).
Writing $T(u) = L_+(u)\, L_-(u)^{-1}$ with $L_\pm(u)$ regular at
$u = \pm\infty$ yields the Reshetikhin--Semenov-Tian-Shansky half-current
presentation
\begin{equation}\label{eq:RTT-half}
  R_{12}(u - v)\, L_{\pm, 1}(u)\, L_{\pm, 2}(v)
  \;=\;
  L_{\pm, 2}(v)\, L_{\pm, 1}(u)\, R_{12}(u - v),
  \quad
  R_{12}(u - v)\, L_{+, 1}(u)\, L_{-, 2}(v)
  \;=\;
  L_{-, 2}(v)\, L_{+, 1}(u)\, R_{12}(u - v).
\end{equation}
The second equation is the gluing cocycle realising $Y^{\mathrm{RTT}}$
as a quantum double of its $L_+, L_-$ halves. The $(12; 21)$ and
$(22; 22)$ entries of \eqref{eq:RTT} expanded in $u^{-1}, v^{-1}$ give
$[t_{12}(u), t_{12}(v)]_+ = 0$ and
$[t_{12}(u), t_{21}(v)]_+ = (\hbar/(u - v))\bigl(t_{11}(v)\, t_{22}(u)
- t_{11}(u)\, t_{22}(v)\bigr)$, the $\fgl(1|1)$-super-Yangian FRT
relations.

\subsection*{$\Omega$-background parameters and free-field realisation}
The two-parameter $\Omega$-background
(Nekrasov--Okounkov~\texttt{hep-th/0306238} \S7) enters through
$H^\bullet_{T^2}(\mathrm{pt}) = \bC[\varepsilon_1, \varepsilon_2]$
acting on the Jordan-quiver coordinates of
Section~\ref{sec:geometric-datum}. The CY constraint
\begin{equation}\label{eq:CY-constraint}
  \varepsilon_1 + \varepsilon_2 + \varepsilon_3 \;=\; 0
  \qquad (T^2_{\mathrm{CY}} \subset T^3)
\end{equation}
preserves $\Omega_\bY$. Free bosons $\phi(z) = \phi_-(z) + \phi_+(z)$
with OPE $\phi(z)\phi(w) \sim (\hbar^2/\varepsilon_1\varepsilon_2)
\log(z - w)$ and free fermions $\psi(z), \bar\psi(z)$ with
$\psi(z)\bar\psi(w) \sim (z - w)^{-1}$ give vertex operators
\begin{equation}\label{eq:vertex-ops}
  t_{11}(u) = :\!e^{+\phi(u)/\hbar}\!:,
  \quad
  t_{22}(u) = :\!e^{-\phi(u)/\hbar}\!:,
  \quad
  t_{12}(u) = :\!e^{+\phi(u)/\hbar}\!:\!\psi(u),
  \quad
  t_{21}(u) = \bar\psi(u)\,:\!e^{-\phi(u)/\hbar}\!:
\end{equation}
satisfying \eqref{eq:RTT} on the equivariant Fock module
$\cF_{(\varepsilon_1, \varepsilon_2)}$ (Nekrasov--Okounkov~2003 \S7
specialised from $\fgl_n$ to $\fgl(1|1)$; Nakajima~1999 Ch.~8).

\begin{proposition}[$\Omega$-background realisation of the conifold bond]
\label{prop:omega-bond-RTT}
\ClaimStatusProposition\
The two-point function of \eqref{eq:vertex-ops} on
$\cF_{(\varepsilon_1, \varepsilon_2)}$ is the Klebanov--Witten shuffle
bond factor:
\[
  \bigl\langle t_{12}(u)\, t_{21}(v) \bigr\rangle_{\cF}
  \;=\;
  \frac{(u - v + \varepsilon_1)(u - v + \varepsilon_2)}
       {(u - v)(u - v + \varepsilon_1 + \varepsilon_2)}
  \;=\;
  k^{\mathrm{off}}(u - v).
\]
Under \eqref{eq:CY-constraint} and $h_1 = \varepsilon_1$,
$h_2 = -\varepsilon_2$, this equals $\varphi_{0 \Rightarrow 1}(u - v)$
of \eqref{eq:bond-factors} on the Costello--Gaiotto branch
$h_1 = \sqrt{\hbar}$. The boson OPE supplies the $(u - v)$ pole; the
fermion bilinear $\psi(u)\bar\psi(v)$ shifts by
$\varepsilon_1 + \varepsilon_2$; the numerator is the refined
Jordan-triple kernel
$(u + \varepsilon_1)(u + \varepsilon_2)(u + \varepsilon_3)/u^3$ at
$\varepsilon_3 = -\varepsilon_1 - \varepsilon_2$.
\end{proposition}

\subsection*{Macdonald--Shiraishi refinement and the rational limit}
The trigonometric lift of \eqref{eq:R-matrix} is the
$U_q(\widehat{\fgl}(1|1))$ $R$-matrix ($u \mapsto z = e^u$,
$\hbar \mapsto \log q$) with Kulish--Sklyanin graded sign on $(22; 22)$
preserved ($q^{-1} z - q$, not $q z - q^{-1}$). The $(q, t)$-refinement
adds a second deformation via Shiraishi--Kubo--Awata--Odake vertex
operators (\texttt{arXiv:q-alg/9507034} eq.~(3.5);
Feigin--Hashizume--Hoshino--Shiraishi--Yanagida
\texttt{arXiv:1002.2485} for the super $(q, t)$-case):
\begin{equation}\label{eq:bond-refined}
  \varphi^{(q, t)}_{0 \Rightarrow 1}(z)
  \;=\;
  \frac{(1 - t z)(1 - q^{-1} z)}
       {(1 - z)(1 - q t^{-1} z)}.
\end{equation}
Specialisations: (a) $t = q^{-1}$ gives the unrefined
$U_q(\widehat{\fgl}(1|1))$ scalar; (b) $q \to 1$ with $t = q^\beta$ at
$\beta = -1$ recovers \eqref{eq:R-matrix}; (c) generic $\beta$ is the
Macdonald--Ruijsenaars deformation.

\begin{remark}[Nekrasov--Shatashvili limit]
\label{rem:NS-limit}
\ClaimStatusProvedElsewhere\ \emph{(Nekrasov--Shatashvili
\texttt{arXiv:0908.4052} \S3; Nekrasov--Pestun--Shatashvili
\texttt{arXiv:1312.6689}; Smirnov \texttt{arXiv:2004.06542} for
stable-envelope / $R$-matrix correspondence.)}
Under $\varepsilon_2 \to 0$, $\varepsilon_1 = \hbar$ fixed,
\eqref{eq:bond-refined} degenerates to the quantum spectral curve
$(1 - q^{-1} z)/(1 - z)$ of the $\fgl(1|1)$ XXX chain; Bethe roots
satisfy
\[
  \prod_j \frac{u_k - \theta_j + \hbar/2}{u_k - \theta_j - \hbar/2}
  \;=\;
  \prod_{l \neq k} \frac{u_k - u_l + \hbar}{u_k - u_l - \hbar}\cdot
  (-1)^{|k|},
\]
$(-1)^{|k|}$ separating bosonic and fermionic strings. Maulik--Okounkov
\texttt{arXiv:1211.1287} stable-envelope $R$-matrices align with the
RTT presentation here; $\mathrm{Res}_{u = 0} R^{MO}(u)$ of
Remark~\ref{rem:pairing-checks} is the NS-limit of \eqref{eq:R-matrix}.
\end{remark}

\subsection*{Gauss decomposition and Drinfeld-new currents}
The super-Gauss decomposition relates \eqref{eq:RTT} to Drinfeld's
current presentation:
\begin{equation}\label{eq:super-gauss}
  T(u) = F(u)\, D(u)\, E(u),
  \ \
  F(u) = \begin{pmatrix} 1 & 0 \\ f(u) & 1 \end{pmatrix},
  \ \
  D(u) = \begin{pmatrix} d_1(u) & 0 \\ 0 & d_2(u) \end{pmatrix},
  \ \
  E(u) = \begin{pmatrix} 1 & e(u) \\ 0 & 1 \end{pmatrix},
\end{equation}
with $e(u), f(u)$ odd (fermionic $(12)$- and $(21)$-entries of $T$;
grading $|e_1| + |e_2| = 1$) and $d_1(u), d_2(u)$ even diagonal.
Super-Ding--Frenkel \texttt{arXiv:q-alg/9608002}
(Zhang \texttt{arXiv:q-alg/9701011} for the explicit $\fgl(1|1)$ case)
extracts
\begin{equation}\label{eq:drinfeld-currents}
  e^{(1)}(u) = d_1(u)^{-1}\, t_{12}(u),
  \quad
  f^{(1)}(u) = t_{21}(u)\, d_1(u)^{-1},
  \quad
  \psi(u) = d_2(u)\, d_1(u)^{-1},
\end{equation}
satisfying $\fgl(1|1)$ loop-algebra relations at lowest order:
$[e^{(1)}(u), f^{(1)}(v)]_+ = \hbar\, \psi(u)/(u - v)$ on the odd pair
and $(e^{(1)}(u))^2 = 0$ (odd Serre from isotropy $(\alpha, \alpha) = 0$;
Remark~\ref{rem:isotropy-source}).

\begin{proposition}[RTT Gauss matches KW shuffle at leading weight]
\label{prop:gauss-to-drinfeld-kw}
\ClaimStatusProposition\
Under \eqref{eq:vertex-ops}, the Drinfeld-new currents satisfy
\[
  e^{(1)}(u)\, e^{(1)}(v) + e^{(1)}(v)\, e^{(1)}(u) = 0,
  \qquad
  e^{(1)}(u) \cdot 1 = 1_{(1, 0)}(u)
  \ \text{in}\ \mathrm{Sh}^{\mathrm{super}}_{\mathrm{KW}},
\]
and
\[
  e^{(1)}(u)\, f^{(1)}(v)
  \;=\;
  \frac{(u - v + \varepsilon_1)(u - v + \varepsilon_2)}
       {(u - v)(u - v + \varepsilon_1 + \varepsilon_2)},
\]
matching Route~R1 Proposition~\ref{prop:shuffle-structure}(1) under
$h_1 = \varepsilon_1$, $h_2 = -\varepsilon_2$ on the CY specialisation.
Odd Serre is fermion nilpotency $\psi(u)^2 = 0$ on the Fock space;
generator identification is vacuum evaluation of the transfer matrix;
leading product is Proposition~\ref{prop:omega-bond-RTT}.
\end{proposition}

\begin{remark}[Subleading $c_{11}$ correction remains conjectural]
\label{rem:rtt-subleading-open}
The asymmetric $c_{11}$ correction of Conjecture~\ref{conj:hopf-pairing}
requires the $u^{-2}$-term of \eqref{eq:RTT}, a derivative of
\eqref{eq:vertex-ops} on the vacuum. Remark~\ref{rem:pairing-checks}
records the open problem.
\end{remark}

% ============================================================================
\section{Route R5 --- Chiral vertex-algebra realisation}
\label{sec:chiral-route}
% ============================================================================

\begin{remark}[Corner assignment is a Li--Yamazaki synthesis, not a Gaiotto--Rapčák theorem]
\label{rem:corner-assignment-scope}
\ClaimStatusConjectured\
Gaiotto--Rapčák \texttt{arXiv:1703.00982} construct $Y_{N_1, N_2, N_3}[\psi]$
as the factorisation algebra on the trivalent junction of three
holomorphic-topological BV theories: Costello 5D holomorphic Chern--Simons
on $\R\times\C^2_i$ ($i=1,2,3$) with BV action
\[
  S_i \;=\; \frac{1}{2\pi i}\int_{\R\times\C^2_i}\!\Omega_i\wedge
  \mathrm{tr}\!\left(A_i\wedge\bar\partial A_i + \tfrac{2}{3} A_i^{\wedge 3}\right),
\]
$\Omega_i$ a holomorphic volume form on $\C^2_i$, $A_i \in
\Omega^{0,\bullet}(\C^2_i;\fgl(N_i))$, glued along half-space interfaces
carrying $\fgl(N_i|N_j)$-valued defect fields (Costello
\texttt{arXiv:1303.2632} \S2; Costello--Gwilliam \texttt{arXiv:2210.13036}
\S4 for the factorisation-algebra formalism). GR Theorem~3 (\S4) identifies
the degenerate vertex $Y_{0,0,N}[\psi] = \cW_N[\psi]$ as a truncation of
$\cW_{1+\infty}[\lambda]$ at specific $\epsilon_i$-ratios. The assignment
of the resolved conifold geometry to $(N_1, N_2, N_3) = (0,1,1)$ is \emph{not}
stated in \texttt{arXiv:1703.00982}; it is the \emph{Li--Yamazaki synthesis}
\texttt{arXiv:2003.08909} §8.3.6, mapping the Klebanov--Witten quiver to
the $(0,1,1)$ corner at $\psi = -\epsilon_1/\epsilon_3$ via the
toric-CY$_3$-to-corner-label dictionary. The identification
\[
  V\bigl(\widehat{\fgl}(1|1)\bigr)^{qv}_\psi
  \;\stackrel{?}{=}\; Y_{0,1,1}[\psi]
\]
is \ClaimStatusConjectured; primary-literature controls are: proved only
on the free-field degeneration $\epsilon\to 0$ (Bershtein--Feigin--Merzon
\texttt{arXiv:1512.08779}) and on the self-dual locus
$\epsilon_1+\epsilon_2=0$. The further identification with the Zamolodchikov
$\cW$-algebra $\cW(\fgl(1|1))_\psi$ runs through Creutzig--Gaiotto
\texttt{arXiv:1706.02834} logarithmic extension and is open at non-generic
$\psi$.
\end{remark}

\begin{proposition}[Holomorphic-topological BV decomposition of 5D hCS observables]
\label{prop:hol-top-decomp}
\ClaimStatusProvedElsewhere\ (Costello--Gwilliam \texttt{arXiv:2210.13036}
\S4.4--\S4.8 for the holomorphic-topological factorisation; Costello
\texttt{arXiv:1303.2632} \S2 for the 5D hCS theory on $\R\times\C^2$.)
The classical BV complex of 5D hCS on $\R\times\C^2$ with gauge algebra
$\fgl(N)$ carries a product decomposition of observables
\[
  \mathrm{Obs}^{\mathrm{cl}}_{\R\times\C^2}
  \;\simeq\;
  \mathrm{Obs}^{\mathrm{1D\,top}}_{\R}
  \,\otimes\,
  \mathrm{Obs}^{\mathrm{4D\,hol}}_{\C^2},
\]
with $\mathrm{Obs}^{\mathrm{1D\,top}}_{\R}$ the topological 1D BV observables
of the $\R$-direction (Costello topological theory) and
$\mathrm{Obs}^{\mathrm{4D\,hol}}_{\C^2}$ the holomorphic observables on
$\C^2$ in the Dolbeault model. The factorisation-algebra product on
$\R\times\C^2$ is the tensor product of the 1D $E_1$-algebra product along
$\R$ and the 4D holomorphic product on $\C^2$; in Dunn--Lurie form the
ambient $E_3$-structure factors as $E_1(\R)\otimes E_2(\C^2 \text{ in hol})$,
with the second factor a \emph{holomorphic} $E_2$ in the sense of
Costello--Gwilliam Lemma~\S6.5.1. At the corner junction $\{0\}\times\C
\subset \R\times\C^2$ the 1D topological factor degenerates to the point
and the restricted factorisation algebra on $\C$ is the \emph{vertex
algebra} $V_\psi$; the corner algebra $Y_{0,1,1}[\psi]$ is the image of
this restriction under the three-way junction gluing.
\end{proposition}

Gaiotto--Rapčák \texttt{arXiv:1703.00982} corner vertex algebra, evaluated
at $(N_1, N_2, N_3) = (0, 1, 1)$ and $\psi = -\epsilon_1/\epsilon_3$
(Li--Yamazaki \texttt{arXiv:2003.08909} §8.3.6 synthesis dictionary), is
\[
  Y_{0,1,1}[\psi]
  \;\stackrel{\mathrm{conj}}{\supseteq}\;
  V\bigl(\widehat{\fgl}(1|1)\bigr)^{qv}_\psi,
\]
a super-vertex algebra arising as the junction restriction of the
Proposition~\ref{prop:hol-top-decomp} factorisation algebra with the
$(0,1,1)$ gauge content on the three half-spaces. The Creutzig--Linshaw
\texttt{arXiv:1706.00031} super-$\cW$-algebra $\cW(\fgl(1|1))_\psi$
is a \emph{logarithmic extension} of $V^{qv}_\psi$ built from the
isotropic odd simple roots (Creutzig--Ridout~\texttt{arXiv:1303.0847};
Ridout--Wood~\texttt{arXiv:1203.3247}); the three objects
$V^{qv}_\psi$, $Y_{0,1,1}[\psi]$, $\cW(\fgl(1|1))_\psi$ are
\emph{genuinely distinct} at non-generic $\psi$, and the chain of
identifications is not a Zamolodchikov $\cW$-algebra in the classical
Hamiltonian-reduction sense (the isotropic root system has no non-trivial
Drinfeld--Sokolov reduction).

The central charge is \emph{not} uniformly $c = 0$: along the family
$\cW(\fgl(1|1))_\psi$ the central charge $c(\psi)$ varies rationally in
$\psi$; $c = 0$ holds only on the self-dual $\epsilon_1 + \epsilon_2 = 0$
locus, where the super-trace annihilates the level contribution. At other
equivariant chambers $c(\psi)$ is a nontrivial rational function of $\psi$
(Creutzig--Gaiotto~\texttt{arXiv:1706.02834} §5). From the Costello 5D hCS
side, $c(\psi)$ is the one-loop wave-function renormalisation of the
vertex-algebra current sector under the Costello--Li
\texttt{arXiv:1606.00365} renormalisation scheme: the propagator
$P_{L<\epsilon}$ on $\R\times\C^2$, regularised at length scale
$L\in(0,\epsilon)$, yields a one-loop Feynman integral whose
$\psi$-dependence is rational through the Costello--Yagi
\texttt{arXiv:1810.01970} vertex computation on $\C^3$, specialised to
the conifold half-space interface. The non-self-dual OPE family is
\ClaimStatusConjectured\ at all orders because super-Kontsevich--Tamarkin
formality is $E_2$-only for the $\fgl(1|1)$ super-Lie input
(Ginzburg--Schedler \texttt{arXiv:0807.0174}); 2-loop and higher
contributions to the OPE are not controlled by existing formality
theorems (cf.\ Theorem~\ref{thm:hcs-super-open}). On the self-dual locus
the OPE presentation reads
\[
\begin{aligned}
  H(z) N(w) &\sim \tfrac{k}{(z-w)^2}, &
  \psi^+(z) \psi^-(w) &\sim \tfrac{k}{(z-w)^2} + \tfrac{H(w)}{z-w}, \\
  H(z) H(w) &\sim 0, & N(z) N(w) &\sim 0 \quad \text{(Cartan self-isotropic)}, \\
  \psi^\pm(z) \psi^\pm(w) &\sim 0 \quad \text{(odd isotropic)}, &
  N(z) \psi^\pm(w) &\sim \pm \tfrac{\psi^\pm(w)}{z-w}, \\
  H(z) \psi^\pm(w) &\sim 0. & &
\end{aligned}
\]
The evaluation map $\mathrm{ev}_\psi : Y^+(\widehat{\fgl}(1|1)) \to V^{qv}_\psi$
is a homomorphism of algebras on the \emph{positive half} (not the full
Yangian, and not an isomorphism); its image is the Gaiotto--Rapčák corner
positive half. This map is the super avatar of the Costello--Yagi
\texttt{arXiv:1810.01970} 5D-hCS-to-Yangian-VOA homomorphism (extended
to the three-author setting in Costello--Gaiotto--Yagi
\texttt{arXiv:2103.01835}),
\[
  \mathrm{ev}^{\mathrm{CGY}}_\psi \,:\, Y^+(\widehat\fg) \;\to\;
  V^{qv}_\psi[\fg]\;=\;\mathrm{Obs}^{\mathrm{5D\,hCS}}_\fg\big|_{\{0\}\times\C},
\]
sourced at the corner junction of the Proposition~\ref{prop:hol-top-decomp}
holomorphic-topological factorisation. For $\fg$ simply-laced with
non-degenerate even Killing form (types A, D, E), CGY prove
$\mathrm{ev}^{\mathrm{CGY}}_\psi$ is an \emph{all-orders-in-$\hbar$}
algebra homomorphism, matched on the CoHA side by
Rap\v{c}\'ak--Soibelman--Yang--Zhao \texttt{arXiv:2007.13365} Theorem~B
through $\CoHA(\C^3_\fg) \twoheadrightarrow Y^+(\widehat\fg)$. For
$\fg = \fgl(1|1)$ with degenerate super-Killing form, the evaluation
homomorphism is controlled only at tree and one-loop order (the
wave-function sector); higher-loop orders are \ClaimStatusConjectured\
pending super-Kontsevich--Tamarkin $E_3$-formality for $\fgl(1|1)$
(Theorem~\ref{thm:hcs-super-open}). The full affine super-Yangian
$Y(\widehat{\fgl}(1|1)) = Y^+ \bowtie Y^-$ acts on the field--state
pairing module $V^{qv}_\psi \otimes (V^{qv}_\psi)^\vee$ via the
Drinfeld-double lift, \emph{not} on $V^{qv}_\psi$ itself as a single
module --- the standard Yangian-versus-evaluation-module distinction
(Tsymbaliuk~\texttt{arXiv:1404.5240} §3). Writing
``$Y(\widehat{\fgl}(1|1)) \to V^{qv}_\psi$'' is an abbreviation whose
precise content is a $Y^+$-action; the $Y^-$-action lives on the dual
module.

\begin{remark}[BV-anomaly first-principles status for $\fgl(1|1)$ 5D hCS]
\label{rem:bv-anomaly-super-firstprinciples}
The one-loop BV anomaly of 5D hCS on $\R\times\C^2$ with gauge algebra
$\fg$ is controlled by the totally-symmetric cubic invariant
$d^{abc}\,{\rm str}(t^a\{t^b,t^c\})$ on $\fg$ (Costello
\texttt{arXiv:1303.2632} \S3; Costello--Gwilliam \texttt{arXiv:2210.13036}
\S12 for the local-BV-cohomology framework; the obstruction lives in
$H^1_{\mathrm{loc}}(\fg, \cO_{\mathrm{loc}})$). For $\fg = \fsl_N$ with
$N\geq 3$ this $d^{abc}$ is nonzero and the theory is obstructed; for
$\fg = \fsl_2$ or $\fg$ rank-one simply-laced the obstruction vanishes
and all-orders quantisation proceeds (Costello--Gaiotto--Yagi
\texttt{arXiv:1810.01970}). For the super case $\fg = \fgl(1|1)$ the
supertrace satisfies $\mathrm{str}(t^a\{t^b,t^c\}) = 0$ by direct basis
enumeration (the $N$-Cartan has $\mathrm{str} = 0$; the odd pair
$\{\psi^+,\psi^-\}$ contributes with opposite signs), so the
\emph{one-loop wheel anomaly vanishes at leading order}. 2-loop and
higher BV obstructions are not auto-killed: super-Kontsevich--Tamarkin
formality for $\fgl(1|1)$ stops at $E_2$ (Ginzburg--Schedler
\texttt{arXiv:0807.0174}), whereas all-orders quantisation of the
generic-$\psi$ family $V^{qv}_\psi$ would require $E_3$-formality of
the classical $\fgl(1|1)$ Poisson structure on $\mathrm{PV}^\bullet(\C^3)$.
The all-orders status of $\mathrm{ev}_\psi$ and of $Y_{0,1,1}[\psi]$ for
generic $\psi$ is therefore \ClaimStatusConjectured\
(Theorem~\ref{thm:hcs-super-open}).
\end{remark}

% ============================================================================
\section{The Hopf pairing and the Drinfeld double}
% ============================================================================

The CoHA $Y^+ = \CoHA(\bY)$ carries, by Davison~\texttt{arXiv:1311.6989}
Thm.~A, a natural bialgebra structure: cohomological Hall product
$m : Y^+ \otimes Y^+ \to Y^+$ and Davison coproduct
$\Delta : Y^+ \to Y^+ \otimes Y^+$ (direct-sum pullback plus specialisation
to the coherent-sheaf moduli stratum). The nil-radical positive half
$Y^+ \subset Y(\widehat{\fgl}(1|1))^{\mathrm{con}}$ is, in Drinfeld-new
presentation, generated by $\{e^{(a)}(u) : a \in \{0,1\}\}$; dually $Y^-$
by $\{f^{(a)}(u)\}$. The question is whether there is a canonical Hopf
pairing
$\langle \cdot, \cdot \rangle : Y^+ \otimes Y^- \to \bC((\hbar))$ making
$D(Y^+)$ a Drinfeld double, and whether that pairing is non-degenerate.
Five successive arguments establish the pairing: (1) universal
$R$-matrix intertwining produces the pairing; (2) triangular decomposition
forces non-degeneracy; (3) the super sign $c_{11} = -1$ descends from
the Kulish--Sklyanin residue at $u = 0$; (4) restriction to the even
diagonal cross-checks against Tsymbaliuk and Schiffmann--Vasserot; (5) the
Drinfeld double recovers the full affine super-Yangian.

% ----------------------------------------------------------------------------
\subsection{Universal $R$-matrix intertwining $Y^+ \otimes Y^-$}
\label{subsec:cycle1-R-matrix-intertwining}
% ----------------------------------------------------------------------------

\begin{proposition}[Universal $R$-matrix determines a Hopf pairing]
\label{prop:cycle1-R-pairing}
\ClaimStatusProvedElsewhere\
\emph{(Drinfeld 1987 ICM \S8; Etingof--Kazhdan
\texttt{arXiv:q-alg/9510022} Prop.~2.1 and \S6.1 for the super extension.)}
For any quasi-triangular Hopf algebra $(H, \Delta, R)$ admitting a
triangular decomposition $H = H^- \otimes H^0 \otimes H^+$, the universal
$R$-matrix factorises uniquely as $R = R^+ \cdot R^0 \cdot R^-$ with
$R^\pm \in H^\pm \hat\otimes H^\mp$ and $R^0 \in H^0 \hat\otimes H^0$. The
factor $R^+$ specifies a canonical bilinear Hopf pairing
$\langle \cdot, \cdot \rangle : H^+ \otimes H^- \to \bC[[\hbar]]$ by
\[
  R^+ \;=\; \sum_\alpha e_\alpha \otimes f^\alpha, \qquad
  \langle e_\alpha, f^\beta \rangle \;=\; \delta_\alpha^\beta,
\]
satisfying $\langle ab, x \rangle = \langle a \otimes b, \Delta(x)\rangle$
and $\langle a, xy \rangle = \langle \Delta(a), x \otimes y\rangle$.
\end{proposition}

\begin{proof}[Derivation sketch]
Apply Drinfeld 1987 \S8 to $Y(\widehat{\fgl}(1|1))^{\mathrm{con}}$, whose
universal $R$-matrix exists by Kulish--Sklyanin 1982 at rational level.
The triangular decomposition $Y = Y^- \otimes Y^0 \otimes Y^+$ is the
super Gauss factorisation $T(u) = F(u) D(u) E(u)$ of
\S\ref{sec:rtt-route} (Zhang~\texttt{arXiv:q-alg/9701011} Thm.~2.3).
Etingof--Kazhdan~\texttt{arXiv:q-alg/9510022}~\S6.1 extends the
factorisation apparatus to the super setting: $R^+$ is the generating
series of a dual basis, so the pairing is canonical once $R$ is canonical.
\end{proof}

\begin{remark}[Existence vs.\ explicit formula]
\label{rem:cycle1-not-formula}
Proposition~\ref{prop:cycle1-R-pairing} shows the pairing \emph{exists}
and is unique up to triangular-decomposition convention; it does not
produce an explicit formula for $\langle e^{(a)}_m, f^{(b)}_n\rangle$ on
Drinfeld-new modes. That requires pushing the universal $R$-matrix
through the RTT-to-new transition of~\S\ref{sec:rtt-route}
(Ding--Frenkel~\texttt{arXiv:q-alg/9608002} super-version;
Zhang~\texttt{arXiv:q-alg/9701011}) and reading off the mode expansion,
executed in~\S\ref{subsec:cycle3-super-sign}.
\end{remark}

% ----------------------------------------------------------------------------
\subsection{Non-degeneracy and the Drinfeld double}
\label{subsec:cycle2-nondegeneracy}
% ----------------------------------------------------------------------------

\begin{proposition}[Non-degeneracy of the Hopf pairing at rational level]
\label{prop:cycle2-nondegeneracy}
\ClaimStatusProvedElsewhere\
\emph{(Drinfeld 1987 \S13; Gow~\texttt{arXiv:math/0411166} Thm.~1.1 super;
Nazarov 1991 \emph{Lett.~Math.~Phys.}~21.)}
The Hopf pairing $\langle \cdot, \cdot \rangle$ of
Proposition~\ref{prop:cycle1-R-pairing} is non-degenerate on each
finite-weight component $Y^+_\gamma \otimes Y^-_{-\gamma}$
($\gamma \in \bZ_{\geq 0}^2$ the $(d_0, d_1)$-charge lattice); the induced
map $Y^+_\gamma \to (Y^-_{-\gamma})^*$ is a $\bC[[\hbar]]$-linear
isomorphism after $\hbar$-adic completion.
\end{proposition}

\begin{proof}[Derivation sketch]
Three independent routes converge. \textbf{(a)} Drinfeld 1987 \S13:
the kernel $K_\gamma \subset Y^+_\gamma$ (were it nontrivial) is a
coideal, $\Delta K \subset K \otimes Y^+ + Y^+ \otimes K$; the coproduct
on Drinfeld-new generators is explicit
(Davison~\texttt{arXiv:1311.6989}~Thm.~A), and $K_\gamma = 0$ follows by
induction on $|\gamma|$. \textbf{(b)} Gow~\texttt{arXiv:math/0411166}
Thm.~1.1 provides a PBW basis $\{e_{ij}^{(r)}\}_{i<j,\,r \geq 0}$ of
$Y^+$ in RTT presentation with matching dual basis in $Y^-$, making
non-degeneracy visible by matrix inversion. \textbf{(c)} Nazarov 1991
Thm.~2 computes
$\langle t_{ij}^{(r)}, t_{kl}^{(s)}\rangle = (-1)^{|i||k|+|j|}
\delta_{il}\delta_{jk}\,[u^{-r} v^{-s}]\log(u-v+\hbar)$, whose
$\log$-kernel is non-degenerate on each homogeneous component.
\end{proof}

\begin{corollary}[Drinfeld double is well-defined]
\label{cor:cycle2-drinfeld-double-welldefined}
\ClaimStatusProvedElsewhere\
\emph{(Drinfeld 1987 \S13; Majid 1995 \emph{Foundations of Quantum Group
Theory} \S7.1.1.)}
Non-degeneracy (Proposition~\ref{prop:cycle2-nondegeneracy}) makes
$D(Y^+) = (Y^+ \otimes (Y^-)^{\mathrm{op}})/\langle\text{pairing}\rangle$
a well-defined Hopf algebra with Drinfeld commutation
\[
  e\, f \;=\; \sum \bigl\langle e_{(1)}, S(f_{(2)})\bigr\rangle\,
    \bigl\langle e_{(3)}, f_{(1)}\bigr\rangle\,
    f_{(2)}\, e_{(2)}
\]
(Sweedler notation), producing $Y(\widehat{\fgl}(1|1))^{\mathrm{con}}$
on Drinfeld-new level.
\end{corollary}

% ----------------------------------------------------------------------------
\subsection{The super sign $c_{11}=-1$ from the $R$-matrix residue}
\label{subsec:cycle3-super-sign}
% ----------------------------------------------------------------------------

\begin{lemma}[Super sign from the Kulish--Sklyanin residue at $u=0$]
\label{lem:cycle3-kulish-sklyanin-residue}
\ClaimStatusProposition\
Let $R(u)$ be the Kulish--Sklyanin super-$R$-matrix on
$\bC^{1|1}\otimes\bC^{1|1}$ from~\eqref{eq:R-matrix}. In the $(2,2;2,2)$-block
(both tensor factors odd)
\[
  R_{22;22}(u) \;=\; 1 - \frac{\hbar}{u}, \qquad
  \mathrm{Res}_{u=0}\bigl[R_{22;22}(u) - 1\bigr] \;=\; -\hbar.
\]
The minus sign is forced by $P^{\mathrm{super}}(e_2 \otimes e_2) =
(-1)^{|e_2||e_2|} e_2 \otimes e_2 = -(e_2\otimes e_2)$, not by
convention.
\end{lemma}

\begin{proof}
Direct from~\eqref{eq:R-matrix}: the $(2,2;2,2)$ entry is $1 - \hbar/u$;
its residue at $u=0$ of $R - 1$ is $-\hbar$. The sign traces to the
super tensor-flip $P^{\mathrm{super}} = \sum_{i,j}(-1)^{|i||j|}
E_{ij}\otimes E_{ji}$, evaluating on $e_2 \otimes e_2$ as $-e_2 \otimes
e_2$. The $-\hbar$ residue is a theorem of the $\bZ/2$-grading, not a
normalisation choice.
\end{proof}

\begin{proposition}[Super sign propagates to the Drinfeld-new pairing]
\label{prop:cycle3-super-sign-propagation}
\ClaimStatusConjectured\
Under the RTT-to-new Gauss decomposition of~\S\ref{sec:rtt-route}, the
$-\hbar$ residue of Lemma~\ref{lem:cycle3-kulish-sklyanin-residue}
propagates to a sign $(-1)^a$ on the Drinfeld-new pairing:
\[
  \mathrm{Res}_{u=v}\,\langle e^{(a)}(u), f^{(a)}(v)\rangle
  \;=\; (-1)^a\,\frac{\hbar}{u-v} + O(\hbar^2).
\]
\end{proposition}

\begin{proof}[Derivation and honest scope]
At leading order the RTT pairing
$\langle t^{(1)}_{ij}, t^{(1)}_{kl}\rangle$ inherits the Nazarov 1991
sign $(-1)^{|i||k| + |j|}$. Under Gauss $T = F D E$, the Drinfeld-new
current $e^{(a)}(u)$ identifies with $t_{a,a+1}(u)$ up to $D$-correction,
$f^{(a)}(u)$ with $t_{a+1,a}(u)$. Substituting $(i,j,k,l)=(a,a+1,a+1,a)$:
\[
  (-1)^{|a||a+1| + |a+1|} \;=\;
  \begin{cases} (-1)^{0\cdot 1 + 1} = -1 & a = 0, \\
                (-1)^{1\cdot 0 + 0} = +1 & a = 1,
  \end{cases}
\]
giving $(-1)^{a+1}$, opposite to the desired $(-1)^a$. The discrepancy
is absorbed by the antipode sign in the $F(u)$-factor
(Zhang~\texttt{arXiv:q-alg/9701011}~\S3;
Gow~\texttt{arXiv:math/0411166}~Rem.~2.4): the antipode of
$F_{a+1,a}(u)$ carries an extra $-1$ on the odd sector, combining with
the Nazarov sign to produce $(-1)^a$ overall. Fully rigorous propagation
to all mode orders $(m,n)$ is not carried out here; hence
\ClaimStatusConjectured\ at the explicit-mode-formula level, while
Lemma~\ref{lem:cycle3-kulish-sklyanin-residue} itself is theorem.
\end{proof}

% ----------------------------------------------------------------------------
\subsection{Even-only cross-check against $Y^+(\widehat{\fgl}_1)$}
\label{subsec:cycle4-even-only-crosscheck}
% ----------------------------------------------------------------------------

\begin{proposition}[Even diagonal matches Tsymbaliuk / Schiffmann--Vasserot]
\label{prop:cycle4-even-diagonal}
\ClaimStatusProvedElsewhere\
\emph{(Tsymbaliuk~\texttt{arXiv:1404.5240} Prop.~4.3;
Schiffmann--Vasserot~\texttt{arXiv:1202.2756} Thm.~1.2.)}
Restrict the conjectural pairing of Conjecture~\ref{conj:hopf-pairing}
below to $a = b = 0$. Then
\[
  \langle e^{(0)}_m, f^{(0)}_n\rangle_\hbar
  \;=\; \hbar^{-1}\,\frac{\binom{m+n}{m}}{m+n+1},
\]
matching Tsymbaliuk~\texttt{arXiv:1404.5240} Prop.~4.3 for
$Y^+(\widehat{\fgl}_1)$ under Schiffmann--Vasserot normalisation
$\hbar = h_1 h_2 / h_3$, equivalently the Feigin--Odesskii shuffle
inner-product on $\bC^\times$-equivariant cohomology of Hilbert schemes
on $\bC^2$.
\end{proposition}

\begin{proof}[Derivation]
The even generator $e^{(0)}(u) = \sum_m e^{(0)}_m u^{-m-1}$ corresponds
under the Schiffmann--Vasserot dictionary to multiplication by $u^m$ on
$H^*_T(\mathrm{Hilb}^n(\bC^2))$; Tsymbaliuk~\texttt{arXiv:1404.5240}
Prop.~4.3 gives the dual-basis pairing
$\binom{m+n}{m}/(m+n+1)\cdot(h_3/h_1 h_2)$. Multiplied by $\hbar^{-1}$
with $\hbar = h_1 h_2/h_3$ this is $\binom{m+n}{m}/(m+n+1)$. No super
sign enters: the $a = b = 0$ sector is bosonic; the conifold even half
restricts to Nakajima's $H^*(\mathrm{Hilb}(\bC^2))$ via the chart
projection $U_\pm = \bC^3 \supset \bC^2$. This is the chamber-independent
statement (Theorem~\ref{thm:chamber-inv-pairing}) on the even sublattice.
\end{proof}

\begin{remark}[Why the odd diagonal is harder]
\label{rem:cycle4-odd-harder}
The odd diagonal $a = b = 1$ admits no commutative Hilbert-scheme model.
It is accessed instead via the Davison--Hennecart--Schlegel-Mejia
\texttt{arXiv:2303.12592} critical CoHA, where signs come from the
vanishing-cycle sheaf $\varphi_{W_{\mathrm{con}}}$ and the quartic
potential. Their Thm.~6.1 establishes a bialgebra isomorphism of odd
generators with a super-Yangian positive half; the explicit $c_{11}$
coefficient is not computed there. This gap is why
Conjecture~\ref{conj:hopf-pairing} below remains \ClaimStatusConjectured.
\end{remark}

% ----------------------------------------------------------------------------
\subsection{The Drinfeld double $D(Y^+) = Y(\widehat{\fgl}(1|1))$}
\label{subsec:cycle5-drinfeld-double}
% ----------------------------------------------------------------------------

\begin{theorem}[Drinfeld-double recovers the full super-Yangian]
\label{thm:cycle5-drinfeld-double-super-yangian}
\ClaimStatusProvedElsewhere\
\emph{(Ding--Frenkel~\texttt{arXiv:q-alg/9608002} for $Y(\widehat{\fgl}_n)$;
Zhang~\texttt{arXiv:q-alg/9701011}~\S3 super;
Gow~\texttt{arXiv:math/0411166} Thm.~1.1 PBW.)}
Non-degeneracy (Proposition~\ref{prop:cycle2-nondegeneracy}) and the
Gauss decomposition $T(u) = F(u) D(u) E(u)$ of Route~R4 together imply
that the Drinfeld double of $Y^+$ with respect to the Hopf pairing of
Corollary~\ref{cor:cycle2-drinfeld-double-welldefined} recovers the full
affine super-Yangian:
\[
  D(Y^+) \;=\; \bigl\langle Y^+ \otimes Y^-,\,\text{Drinfeld commutation}
  \bigr\rangle
  \;\simeq\; Y(\widehat{\fgl}(1|1))^{\mathrm{con}}.
\]
\end{theorem}

\begin{proof}[Derivation sketch]
Gauss decomposition identifies $E(u)$ with $Y^+$ (upper-unipotent),
$F(u)$ with $Y^-$ (lower-unipotent), $D(u)$ with $Y^0$ (diagonal). The
RTT relation $R(u-v) T_1(u) T_2(v) = T_2(v) T_1(u) R(u-v)$ decomposes
after Gauss substitution into (a) within-$E$ relations giving $Y^+$
structure, (b) within-$F$ relations giving $Y^-$ structure, (c)
cross-terms $E\otimes F$ giving Drinfeld commutation matching
Corollary~\ref{cor:cycle2-drinfeld-double-welldefined}. This is
Zhang~\texttt{arXiv:q-alg/9701011} Thm.~3.1 for $DY(\fgl(1|1))$ at
toroidal level; Galakhov--Li--Yamazaki~\texttt{arXiv:2106.01230}
establishes the conifold-specific toroidalisation, giving
$DY(\widehat{\fgl}(1|1))^{\mathrm{con}}$. Recovery is theorem,
independent of the conjectural $c_{11}$ subleading correction.
\end{proof}

% ----------------------------------------------------------------------------
\subsection{Synthesis: conjectural mode formula and honest status}
\label{subsec:hopf-pairing-synthesis}
% ----------------------------------------------------------------------------

\begin{conjecture}[Hopf pairing explicit formula]
\label{conj:hopf-pairing}
\ClaimStatusConjectured\
Assuming the Gauss-decomposition sign propagation of
Proposition~\ref{prop:cycle3-super-sign-propagation} extends to all mode
orders $(m, n)$, the Drinfeld-new Hopf pairing
$\langle \cdot, \cdot \rangle_\hbar : Y^+ \otimes Y^- \to \bC((\hbar))$
takes the explicit form
\begin{equation}\label{eq:hopf-pairing-formula}
  \langle e^{(a)}_m, f^{(b)}_n \rangle_\hbar
  \;=\; \delta^{a,b}\, (-1)^a\, \hbar^{-1}\,
    \frac{\binom{m+n}{m}}{m+n+1}\,
    (1 + \hbar\, c_{ab}\, \delta_{m,n}),
\end{equation}
with $c_{00} = 0$, $|c_{11}| = 1$ forced by
Lemma~\ref{lem:cycle3-kulish-sklyanin-residue}, and sign of $c_{11}$
determined by the antipode-in-coproduct convention:
\begin{itemize}[leftmargin=2em]
\item \emph{Nazarov 1991 / Kulish--Sklyanin 1982 convention} (antipode
  $S$ post-multiplies): $c_{11} = -1$, consistent with the $-\hbar$
  residue of Lemma~\ref{lem:cycle3-kulish-sklyanin-residue};
\item \emph{Gow 2005~\texttt{arXiv:math/0411166} / Peng 2011 convention}
  ($S$ pre-multiplies): $c_{11} = +1$.
\end{itemize}
Off-diagonal entries ($a \neq b$) vanish by parity orthogonality. The
$\hbar$-deformation parameter relates to the equivariant parameters
$h_1, h_2, h_3 = -h_1 - h_2$ of $T\subset\mathrm{GL}(3)$ by the
Schiffmann--Vasserot normalisation $\hbar = h_1 h_2 / h_3$
(\texttt{arXiv:1202.2756} \S3.2). The even diagonal ($a = 0$) is theorem
(Proposition~\ref{prop:cycle4-even-diagonal} via
Tsymbaliuk~\texttt{arXiv:1404.5240} Prop.~4.3 and
Schiffmann--Vasserot~\texttt{arXiv:1202.2756} Thm.~1.2); the odd-diagonal
($a = 1$) mode-level formula is conjectural, pending a direct
Gauss-decomposition residue computation verifying $\binom{m+n}{m}/(m+n+1)$
against the Nazarov~1991 Thm.~2 Berezinian-residue form
\[
  \langle t^{(r)}_{ij}, t^{(s)}_{kl} \rangle
  \;=\; (-1)^{|i||k| + |j|}\, \delta_{il}\, \delta_{jk}\,
    [u^{-r} v^{-s}]\bigl[\log \mathrm{Ber}_q T(u)\cdot k^{\mathrm{off}}(u-v)\bigr],
\]
with offset kernel $k^{\mathrm{off}}(u) = (u+\varepsilon_1)(u+\varepsilon_2)/
[u(u+\varepsilon_1+\varepsilon_2)]$. The binomial form is a legitimate
shadow under the Drinfeld--Jimbo evaluation map to the bosonic subalgebra
$Y^+(\widehat{\fgl}_1)$ (\S\ref{subsec:cycle4-even-only-crosscheck});
its full validity on the odd diagonal requires verification against
Nazarov's $\log$-kernel and an independent Gow--Peng cross-derivation
(\texttt{arXiv:math/0411166}; Peng~2011 \emph{J.~Algebra}~325;
Zhang~\texttt{arXiv:q-alg/9701011}), open at mode level.
\end{conjecture}

\begin{remark}[Three independent cross-checks for the leading term]
\label{rem:pairing-checks}
Three parallel derivations converge on the leading binomial coefficient.
\textbf{(a)} Maulik--Okounkov stable-envelope residue
\[
  \mathrm{Res}_{u=0}\,R^{MO}(u)
  \;=\; \frac{h_1 h_2}{h_1 + h_2}\cdot P^{\mathrm{super}},
\]
matches~\eqref{eq:hopf-pairing-formula} at $(m, n) = (0, 0)$ under the
Costello--Gaiotto normalisation $h_1 = \sqrt{\hbar}$
(Maulik--Okounkov~\texttt{arXiv:1211.1287}~\S4.5).
\textbf{(b)} Soibelman--Gelfand refined Ext-pairing
$\mathbf{E}_\hbar(M, N) = \sum_i (-\hbar^{1/2})^i \dim \Ext^i(M, N)$
with Milnor--Moore extension matches the leading Drinfeld-double pairing
coefficient.
\textbf{(c)} The RTT Hopf pairing of Nazarov~1991 Thm.~2 and
Gow~\texttt{arXiv:math/0411166} Thm.~1.1 agrees under the
Ding--Frenkel~\texttt{arXiv:q-alg/9608002} /
Zhang~\texttt{arXiv:q-alg/9701011} RTT-to-new dictionary on the leading
term. These three derivations agree on the binomial coefficient but
neither independently establishes the $c_{11} = -1$ subleading correction
at all modes $(m, n)$: (a) gives only the residue at $u = 0$; (b) gives
only the Ext-pairing leading term; (c) produces a $\log$-kernel rather
than a pure binomial. The odd-diagonal pairing is therefore
\ClaimStatusConjectured\ at subleading order; the even diagonal is
theorem (Proposition~\ref{prop:cycle4-even-diagonal}).
\end{remark}

\begin{theorem}[Hopf status of the Drinfeld double]
\label{thm:hopf-status}
\ClaimStatusProvedElsewhere\
\emph{(Etingof--Kazhdan~\texttt{arXiv:q-alg/9510022} and its super
extension in~\texttt{arXiv:q-alg/9510023--25}; Jimbo--Konno~1997 for the
elliptic case; Zhang~\texttt{arXiv:q-alg/9701011} for the super-Yangian
Gauss decomposition.)}
The Drinfeld double $D(Y^+) = Y(\widehat{\fgl}(1|1))^{\mathrm{con}}$ of
Theorem~\ref{thm:cycle5-drinfeld-double-super-yangian} is
\begin{itemize}[leftmargin=2em]
\item \emph{strict Hopf} at rational, trigonometric, and toroidal levels
  (associator $\Phi = 1$);
\item \emph{quasi-Hopf} at elliptic level, with associator
  $\Phi_{\mathrm{ell}}$ built from the Felder--Jimbo--Konno dynamical
  $R$-matrix and theta-function boundary data.
\end{itemize}
The ``quasi-Hopf from super $+$ Serre'' intuition fails: $\fgl(1|1)$'s
supersymmetric pairing is even-split, so the Etingof--Kazhdan obstruction
$H^3_{\mathrm{EK}}(\fgl(1|1), \fgl(1|1)^{\otimes 3}) = 0$; KW Serre
relations are wheel conditions on the shuffle kernel, not coproduct
associators; the quartic potential $W_{\mathrm{con}}$ enters the Jacobi
relations, not the coassociator. Hopf status is independent of
Conjecture~\ref{conj:hopf-pairing}'s subleading $c_{11}$ correction:
strict Hopf at rational level follows from non-degeneracy
(Proposition~\ref{prop:cycle2-nondegeneracy}) alone.
\end{theorem}

% ============================================================================
\section{Chamber structure}
% ============================================================================

\begin{theorem}[Chamber structure, Nagao--Nakajima \texttt{arXiv:0809.2992}]
\label{thm:chambers}
\ClaimStatusProvedElsewhere\
The stability space of $J(Q_{\mathrm{con}}, W_{\mathrm{con}})$-modules
forms a $W(\widehat{\fsl}_2)$-torsor with infinitely many chambers
generated by Seiberg dualities (equivalently, Keller--Yang
\texttt{arXiv:0906.0761} quiver mutations) acting as simple reflections
on the charge lattice. Named fundamental chambers:
\begin{itemize}[leftmargin=2em]
\item Non-commutative (Szendrői): $(s_0, s_1) = (0, 0)$,
\item Donaldson--Thomas (large-volume): $(s_0, s_1) = (0, -1)$ or $(-1, 0)$,
\item Pandharipande--Thomas (stable pairs): $(s_0, s_1) = (-1, -1)$,
\item Empty: $(s_0, s_1) = (+\infty, +\infty)$ limit.
\end{itemize}
Shifted quiver Yangian (Galakhov--Li--Yamazaki \texttt{arXiv:2106.01230}):
each chamber corresponds to a shift $\mathbf{s} \in \bZ^2$ with
Braverman--Finkelberg--Nakajima shift homomorphism
$\iota_{\mathbf{s}, \mathbf{s}+1} : Y_{\mathbf{s}} \hookrightarrow
Y_{\mathbf{s}+1}$. The underlying algebra $Y^+(\widehat{\fgl}(1|1))$ is
chamber-independent; representations are chamber-parametrised.
\end{theorem}

\begin{remark}[Bridgeland stability space; dimension clarification]
\label{rem:bridgeland-stab}
The Klebanov--Witten Jacobi algebra $A = J(Q_{\mathrm{con}},
W_{\mathrm{con}})$ has two simple modules $S_0, S_1$; its numerical
Grothendieck group $K_0^{\mathrm{num}}(A\text{-mod})$ is free of rank
two, generated by $[S_0], [S_1]$. Bridgeland's theorem on stability
conditions (Bridgeland \texttt{arXiv:math/0212237} Thm.~1.2; refined
for projective $K3$ surfaces in \texttt{arXiv:1002.4372}) applies to
$D^b(A) \simeq D^b(\mathrm{Coh}\,\bY)$: $\mathrm{Stab}(D^b(A))$ is a
complex manifold of local dimension equal to $\mathrm{rk}\,
K_0^{\mathrm{num}} = 2$, and the forgetful map $\mathcal{Z} :
\mathrm{Stab} \to \Hom(K_0^{\mathrm{num}}, \bC) = \bC^2$,
$\sigma \mapsto Z_\sigma$, is a local biholomorphism. The
clarification is load-bearing: ``$\dim_\bC \mathrm{Stab} = 2$'' is
the rank of the numerical $K$-group, \emph{not} the Calabi--Yau
dimension $d = 3$; Bridgeland's \texttt{arXiv:1002.4372} is executed
on a $K3$ surface ($d = 2$), whereas the conifold is a $d = 3$ local
model, and the transport proceeds via Nagao--Nakajima's dim-$3$
construction (\texttt{arXiv:0809.2992}) on the tilted
$D^b(A\text{-mod})$ at a single shift in the heart ladder. Central
charges in the reduced slice are $Z_\sigma(S_i) = m_i e^{i\pi
\varphi_i}$, $\varphi_i \in (0, 1]$, with named chambers: Szendrői
$(\tfrac{1}{3}, \tfrac{1}{3})$, $\mathbf{s} = (0, 0)$; DT $\varphi_0
< \varphi_1$ with $\varphi_0 + \varphi_1 < 1$, $\mathbf{s} = (0, -1)$
or $(-1, 0)$; PT both shifts $-1$; Empty $\mathbf{s} \to (+\infty,
+\infty)$.
\end{remark}

\begin{remark}[DT / PT / NCDT are chambers; GV is a Kähler limit]
\label{rem:chambers-vs-kahler}
DT chamber versus GV invariants is a category separation. A chamber
is a connected component of $\mathrm{Stab} \setminus \bigcup_\alpha
W_\alpha$; the GV invariants $n_g^\beta(\bY)$ of Gopakumar--Vafa are
integers attached to the large-radius point of the Kähler moduli
$M_K(\bY)$ through the MNOP/PT correspondences. The three stability
chambers NCDT, DT, PT sit over the DT/PT regime of $M_K$ under the
central-charge period map; the GV generating function is
reconstructed from DT (equivalently PT) data via
\[
  Z_{\mathrm{PT}}(\bY; q, Q) \;=\; \prod_{n \geq 1} (1 - (-q)^n Q)^n,
\]
with $n_0^\beta = 1$ for $\beta = \beta_0$ the primitive curve class
and higher-$g$ GV invariants vanishing on the local conifold
(Gopakumar--Vafa \texttt{hep-th/9812127}; Bryan--Morrison
\texttt{arXiv:0812.5007} Thm.~2). GV is extracted at the
large-radius limit, not a chamber of $\mathrm{Stab}$.
\end{remark}

\begin{theorem}[KS wall-crossing as a \v{C}ech $H^1$-cocycle on the chamber cover]
\label{thm:ks-h1-cocycle}
\ClaimStatusProvedElsewhere\ \emph{(Kontsevich--Soibelman
\texttt{arXiv:0811.2435} §2.5, §6.3; structural re-statement.)}
Fix the open cover $\{U_\alpha\}_{\alpha \in
\pi_0(\mathrm{Stab}_{\mathrm{gen}})}$ of the generic locus by
chambers and the refinement $\{U_{\alpha\beta}\}$ by codimension-one
walls. Let $G_{\mathrm{KS}}$ denote the pro-nilpotent group of
motivic symplectomorphisms of the quantum torus
$\bT_{Q_{\mathrm{con}}}$ associated to the charge lattice
$K_0^{\mathrm{num}}(A)$. The wall-crossing factors
$\mathbf{U}_{\alpha\beta} \in G_{\mathrm{KS}}$ constitute a \v{C}ech
$1$-cochain; the KS wall-crossing formula (\texttt{arXiv:0811.2435}
Thm.~6.3.1) is precisely the $1$-cocycle condition
\[
  \prod^\curvearrowleft_{W \cap \gamma}\, \mathbf{U}_W \;=\; \mathbf{1}
\]
for every contractible loop $\gamma$ in $\mathrm{Stab}$; on each
codimension-two stratum the product of wall operators around the
stratum equals the identity. The cocycle class $[\mathbf{U}] \in
H^1(\mathrm{Stab}; G_{\mathrm{KS}})$ is trivial because every loop in
$\mathrm{Stab}_{\mathrm{gen}}$ bounds a disk meeting walls
transversely; cocycle-triviality is the wall-crossing formula.
\end{theorem}

\begin{proof}[Derivation]
On a generic chamber $U_\alpha$ the BPS spectrum
$\Omega_\alpha(\gamma)$ is constant; the motivic automorphism
\[
  \mathbf{U}_\alpha
  \;=\; \prod_{\gamma : Z_\sigma(\gamma) \in \bH^-,\,
    \mathrm{arg}\,\nearrow}
      \mathbf{U}_\gamma^{\Omega_\alpha(\gamma)},
  \qquad
  \mathbf{U}_\gamma(x_{\gamma'})
  \;=\; (1 - x_\gamma)^{\langle \gamma, \gamma'\rangle}\, x_{\gamma'},
\]
is the formal symplectomorphism recording BPS states ordered by the
phase of $Z_\sigma$ on the lower half-plane. Across a wall
$U_{\alpha\beta}$, the KS identity $\mathbf{U}_\alpha =
\mathbf{U}_\beta$ holds; writing $\mathbf{T}_{\alpha\beta} =
\mathbf{U}_\alpha \mathbf{U}_\beta^{-1}$, the transition function is
a word in $\mathbf{U}_\gamma$ on shared rays. On a codimension-two
stratum $U_{\alpha\beta\gamma}$ the \v{C}ech cocycle condition
$\mathbf{T}_{\alpha\beta} \mathbf{T}_{\beta\gamma}
\mathbf{T}_{\gamma\alpha} = \mathbf{1}$ becomes the pentagon identity
of Theorem~\ref{thm:pentagon} below. Since
$\mathrm{Stab}_{\mathrm{gen}}$ is simply connected (Bridgeland
\texttt{arXiv:math/0212237} Thm.~7.1, local contractibility), every
loop bounds and the cocycle is a coboundary.
\end{proof}

\begin{theorem}[Pentagon identity for the conifold (Faddeev--Kashaev form)]
\label{thm:pentagon}
\ClaimStatusProvedElsewhere\ \emph{(Faddeev--Kashaev
\texttt{hep-th/9310070}; specialised to the conifold via KS
\texttt{arXiv:0811.2435} §6.4.)}
Let $x_0, x_1$ be the formal torus generators associated to the
charge classes $[S_0], [S_1] \in K_0^{\mathrm{num}}(A)$, satisfying
\[
  x_1 x_0 \;=\; q^{\langle [S_0], [S_1]\rangle_E}\, x_0 x_1
  \;=\; q^{-1} x_0 x_1,
\]
with $\langle \cdot, \cdot\rangle_E$ the antisymmetrised Euler form
and $\langle [S_0], [S_1]\rangle_E = -1$. Let $\Psi(x) = E_q(x)^{-1}
= \prod_{k \geq 0} (1 - q^{k+1/2} x)$ be the refined Faddeev quantum
dilogarithm. Across the codimension-two wall intersection
corresponding to the Seiberg duality $S_0 \leftrightarrow S_0 \oplus
S_1$,
\begin{equation}\label{eq:pentagon}
  \Psi(x_0)\, \Psi(x_1)
  \;=\; \Psi(x_1)\, \Psi(q^{-1/2} x_0 x_1)\, \Psi(x_0).
\end{equation}
\end{theorem}

\begin{proof}[Derivation]
Two inputs combine. (i) Faddeev--Kashaev pentagon for $v u = q u v$
(Faddeev--Kashaev \texttt{hep-th/9310070} Eq.~(15); Kashaev--Nakanishi
\texttt{arXiv:1104.4630} Eq.~(2.7)): $\Psi(u) \Psi(v) = \Psi(v)
\Psi(q^{-1/2} u v) \Psi(u)$, valid whenever $[u, v]_q = (q-1) u v$.
(ii) BPS index $\Omega(\gamma) = 1$ for $\gamma \in \{[S_0], [S_1],
[S_0] + [S_1]\}$ a KS wall-crossing invariant (MMNS
\texttt{arXiv:1107.5017} Thm.~1; Bryan--Morrison
\texttt{arXiv:0812.5007} Prop.~1).

Setting $u = x_0$, $v = x_1$: Euler form $\langle [S_0], [S_1]
\rangle_E = -1$ gives $x_1 x_0 = q^{-1} x_0 x_1$; absorbing the sign
into $q \mapsto q^{-1}$ recovers the FK form. The class $[S_0] +
[S_1]$ on the RHS is the wall-produced bound state ($\bP^1$-wrapping
D-brane), BPS index $1$; $q^{-1/2} x_0 x_1$ is the quantum-shifted
torus class. Equivalently by KS: on the triple-wall stratum at
$Z(S_0) = Z(S_1) = i\epsilon$, $\epsilon \to 0^+$, three rays
$\gamma \in \{[S_0], [S_0]+[S_1], [S_1]\}$ converge, and
$\mathbf{U}_{[S_0]} \mathbf{U}_{[S_1]} = \mathbf{U}_{[S_1]}
\mathbf{U}_{[S_0]+[S_1]} \mathbf{U}_{[S_0]}$ translates under the
$\Psi$-representation to \eqref{eq:pentagon}. The $q$-binomial
$\sum_{k=0}^n \binom{n}{k}_q (-1)^k q^{\binom{k}{2}} x^k =
\prod_{j=0}^{n-1}(1 - q^j x)$ is the finite-dimensional shadow
discharging the identity order-by-order in $q$.
\end{proof}

\begin{remark}[Pentagon as KS 3-cocycle and as bound-state mechanism]
\label{rem:pentagon-cocycle}
The pentagon \eqref{eq:pentagon} discharges two roles: (a) the
triple-wall \v{C}ech cocycle condition of
Theorem~\ref{thm:ks-h1-cocycle}, and (b) the Denef multi-centre
attractor flow-tree identity, with the middle factor $\Psi(q^{-1/2}
x_0 x_1)$ recording the bound state $(S_0, S_1)$ that appears on one
side of the wall and disappears on the other with BPS index
preserved.
\end{remark}

\begin{theorem}[CoHA chamber-independence via Davison--Meinhardt Sym-decomposition]
\label{thm:coha-chamber-inv}
\ClaimStatusProvedElsewhere\ \emph{(Davison--Meinhardt
\texttt{arXiv:1601.02479} Thm.~A, Thm.~B; super-extension
Davison--Hennecart--Schlegel-Mejia \texttt{arXiv:2303.12592}.)}
The critical cohomological Hall algebra $\CoHA(Q_{\mathrm{con}},
W_{\mathrm{con}})$ is chamber-independent: for any two generic
$\sigma, \sigma' \in \mathrm{Stab}$ there is a canonical Hopf-algebra
isomorphism $\CoHA_\sigma \xrightarrow{\sim} \CoHA_{\sigma'}$ induced
by the Davison--Meinhardt Sym-decomposition
\[
  \CoHA_\sigma(Q_{\mathrm{con}}, W_{\mathrm{con}})
  \;\cong\;
  \mathrm{Sym}^{\boxtimes_+}_{\mathrm{super}}\!\bigl(
    H(B\bC^\times)_{\mathrm{vir}} \otimes
    \fg_{\mathrm{BPS}}^{\mathrm{con}}
  \bigr),
\]
whose right-hand side depends only on
$\fg_{\mathrm{BPS}}^{\mathrm{con}}$ and $\bQ[u]$, both
chamber-invariants.
\end{theorem}

\begin{proof}[Derivation]
Davison--Meinhardt \texttt{arXiv:1601.02479} Thm.~B (integrality):
for any symmetric quiver with potential satisfying Mozgovoy--Reineke
optimality $\mathrm{HH}^{<0}(J) = 0$, the motivic BPS integers
$\Omega^{\mathrm{mot}}(\gamma) \in \bZ[\bL^{\pm 1/2}]$ are
wall-crossing-invariants. $W_{\mathrm{con}}$ is MR-optimal
(\texttt{arXiv:0708.1259} Prop.~5.2); BPS indices on
$\Delta_+^{\mathrm{re}} \cup \Delta_+^{\mathrm{im}}$ of
Theorem~\ref{thm:mmns} are chamber-invariant. The Sym-decomposition
\texttt{arXiv:1601.02479} Thm.~A presents $\CoHA_\sigma$ functorially
in the BPS data, so $(\fg_{\mathrm{BPS}}^{\mathrm{con}})_\sigma =
(\fg_{\mathrm{BPS}}^{\mathrm{con}})_{\sigma'}$ promotes to a
Hopf-algebra isomorphism on $\CoHA$. The super-extension
\texttt{arXiv:2303.12592} confirms this at the $\bZ/2$-graded level.
\end{proof}

\begin{theorem}[Chamber-independence of the Hopf pairing]
\label{thm:chamber-inv-pairing}
\ClaimStatusProvedElsewhere\ \emph{(Joyce--Song
\texttt{arXiv:0810.5645} Thm.~5.11 + KS wall-crossing
\texttt{arXiv:0811.2435} §6.3.)}
The Joyce--Song integration map
$I : \mathrm{SF}^{\mathrm{ind}}(\mathfrak{M}) \to
\widehat{U}(\fg_{\mathrm{BPS}})[[\hbar]]$
is $KS$-equivariant, and wall-crossing corrections
$\alpha_W = \sum_{\gamma = k\gamma_W} \Omega(\gamma)\,[\gamma]/k^2$
land in orthogonal $K_0$-graded pieces. Hence the Hopf-pairing
matrix entries are literally equal across chambers, not merely
projectively equivariant; combined with
Theorem~\ref{thm:coha-chamber-inv}, the entire Hopf-algebraic
structure of $\CoHA(Q_{\mathrm{con}}, W_{\mathrm{con}})$ (product,
coproduct, antipode, pairing) is a chamber-invariant.
\end{theorem}

\begin{remark}[Synthesis: three instantiations of one chamber-invariance]
\label{rem:three-chamber-invariances}
Theorems~\ref{thm:ks-h1-cocycle}, \ref{thm:coha-chamber-inv}, and
\ref{thm:chamber-inv-pairing} are three faces of a single invariance.
KS cocycle-triviality says the transition data between chambers,
viewed in $G_{\mathrm{KS}}$, lifts to a trivial $H^1$ class.
Sym-decomposition says the CoHA Hopf-algebra transition is canonical.
Integration-map equivariance is the pairing-level shadow. Each is a
theorem at its stated scope; none is a re-statement of the others.
What unifies them is Davison--Meinhardt integrality: the BPS indices
$\Omega^{\mathrm{mot}}(\gamma)$ are the sole wall-invariant data, and
every chamber-invariance statement factors through them.
\end{remark}

% ============================================================================
\section{Physical origin --- the M-theory parent}
% ============================================================================

\begin{conjecture}[M-theory parent of the conifold super-Yangian]
\label{conj:m-theory-parent}
\ClaimStatusConjectured\
The M-theory parent of $Y(\widehat{\fgl}(1|1))^{\mathrm{con}}$ is the 6D
(2,0) $A_{K-1}$ theory on $\bR^{1,3} \times T^2$ with $K$ M5-branes
wrapped on $\bP^1 \subset \bY$, normal-bundle twist $\cO(-1) \oplus
\cO(-1)$, $\Omega$-deformation on transverse Taub--NUT, and
Costello--Gwilliam-twist on the CY$_3$ target. Dimensional reductions:
\begin{itemize}[leftmargin=2em]
\item \emph{Bulk}: 5D holomorphic Chern--Simons on $\bR \times \bY$ with
  gauge $\fgl(1|1)$-valued (open per
  Theorem~\ref{thm:hcs-super-open}).
\item \emph{Boundary}: 3D $\cN = 4$ sigma model to $\mathrm{Hilb}^n(\bY)$
  whose Coulomb-branch twist is conjecturally the super-Yangian positive
  half $Y^+(\widehat{\fgl}(1|1))$ at infinity.
\end{itemize}
The Coulomb-branch identification transfers Braverman--Finkelberg--Nakajima
$\cA_\hbar(\mathrm{K3}) = Y(\frakg_{\mathrm{K3}})$
(\texttt{arXiv:1601.03586},~\texttt{arXiv:1604.03625}) and Nakajima
\texttt{arXiv:1606.09041} to the conifold setting, with Putrov--Yagi
\texttt{arXiv:1704.04255} and Kamnitzer--Pham--Weekes \texttt{arXiv:1806.07146}
providing the 3D $\cN=4$ Coulomb-branch extension machinery for ADE quivers.
The super-quiver Klebanov--Witten extension, however, is \emph{not} proved
rigorous at the Coulomb-branch level in any published source: BFN requires
a non-degenerate Killing form, which $\fgl(1|1)$ lacks, and the Putrov--Yagi
4D/3D descent was established for $\cN = 2$ SCFTs whereas KW is $\cN = 1$
(cf.~Theorem~\ref{thm:bllprvr-fails}). This statement is therefore a
synthesis by analogy with the Costello--Paquette $\bC^3$ construction
(\texttt{arXiv:2001.02177} and~\texttt{arXiv:2201.02595}), not an explicit
theorem. The specific wrap on a compact $\bP^1$ with normal bundle
$\cO(-1) \oplus \cO(-1)$, combined with the super gauge extension, lies
beyond the Costello--Paquette theorems as written and beyond the
BFN/Putrov--Yagi Coulomb-branch framework as it stands.
\end{conjecture}

\begin{theorem}[Superconformal BLLPRvR route is unavailable]
\label{thm:bllprvr-fails}
\ClaimStatusProvedElsewhere\ \emph{(by type-mismatch)}
The Beem--Lemos--Liendo--Peelaers--Rastelli--van~Rees
\texttt{arXiv:1312.5344} 4D/2D chiral-algebra correspondence applies to
4D $\cN=2$ superconformal theories. The 4D field theory on an M5-brane
wrapped on $\bP^1 \subset \bY$ with KW normal bundle is the
Klebanov--Witten $\cN = 1$ SCFT, not $\cN = 2$. Hence BLLPRvR does not
directly apply to this boundary theory. The identification
$\CoHA(\bY) = Y^+(\widehat{\fgl}(1|1))$ therefore runs through the
Rapčák--Soibelman--Yang--Zhao \texttt{arXiv:1810.10402} route on the
CoHA side, not through BLLPRvR.
\end{theorem}

\begin{theorem}[5D hCS super extension is open]
\label{thm:hcs-super-open}
\ClaimStatusOpen\
The Costello--Gaiotto--Yagi all-orders theorem for 5D hCS $\to$
Yangian VOA (Costello 2013 \texttt{arXiv:1303.2632}; Costello--Yagi
\texttt{arXiv:1810.01970}; the three-author CGY paper with
Gaiotto is \texttt{arXiv:2103.01835}) requires simply-laced gauge
algebra with non-degenerate even Killing form. The super case
$\widehat{\fgl}(1|1)$ has degenerate Killing form (the supertrace
annihilates $K$); super Kontsevich--Tamarkin formality is $E_2$-only
(Ginzburg--Schedler \texttt{arXiv:0807.0174}), not $E_3$; higher-loop
obstructions in $H^1_{\mathrm{loc}}(\fgl(1|1), \cO_{\mathrm{loc}})$
are not auto-killed. One-loop wheel vanishing
$d^{abc}_{\mathrm{super}} = 0$ for $\fgl(1|1)$ is verified by direct
basis computation; $2$-loop and higher remain open.
\end{theorem}

\begin{remark}[Per-chart $\widehat{\fgl}_1$ is all-orders rigorous]
The per-chart abelian 5D hCS with $\widehat{\fgl}_1$ gauge on each
$U_\pm \cong \bC^3$ is all-orders rigorous (Costello--Li, Schiffmann--Vasserot);
the difficulty is in the super-assembly of two such algebras via the
FM transition.
\end{remark}

\begin{remark}[Anomaly vanishing on $\bY$]
\label{rem:anomaly-vanishing}
The BCOV $1$-loop anomaly $\alpha_{\mathrm{BCOV}} = (\chi/24)[\Omega]^{0,1}
\in H^{0,1}(\bY)$ (Costello--Li 2016) vanishes because $\bY$ is a
rank-$2$ vector bundle over $\bP^1$ (hence homotopy-equivalent to
$\bP^1$ via the affine bundle retraction), and
$H^{0,1}(\bP^1) = h^{0,1}(\bP^1) = 0$. The space $\bY$ is
quasi-projective (not affine) since $\bP^1$ is proper.
\end{remark}

\begin{remark}[Non-abelian $K \geq 2$ obstructions]
For $A_{K-1}$ stacks of M5-branes, the 1-loop $d^{abc}d^{abc}$ anomaly
obstructs perturbative consistency at $K \geq 3$. The $K = 2$ case is
special because $\fsl_2$ has $d^{abc} = 0$ (Costello 2015 \texttt{arXiv:1112.0816}
§4--5; CLAUDE.md key-fact). Note $\fsl_2 \neq \fgl_2$: the $\fgl_2 =
\fsl_2 \oplus \fu(1)$ decomposition has vanishing symmetric cubic on
$\fsl_2$ and BRST-trivial $\fu(1)$, so $K = 2$ gives a well-defined
5D CS at the gauge-algebra level.
\end{remark}

% ============================================================================
\section{Super-rigidity, Gopakumar--Vafa, and refined topological vertex}
% ============================================================================

M-theory on $\bY \times S^{1}$ with M2-branes wrapping $k[\bP^{1}] \in
H_{2}(\bY,\bZ)$ presents these wrapped branes as 5D BPS hypermultiplets
charged under the graviphoton $F_{RR}$; integrating them out by a
Schwinger one-loop produces the $2\sinh(k g_{s}/2)$ denominator and the
$\mathrm{Li}_{3}$-polylogarithm envelope of the conifold free energy
(Gopakumar--Vafa~\texttt{arXiv:hep-th/9809187} Part~I;
\texttt{arXiv:hep-th/9812127} Part~II). The two-parameter
$(\varepsilon_{1}, \varepsilon_{2})$-refinement deforms the Schwinger
integrand to a product of two sinh factors (Nekrasov--Okounkov
\texttt{arXiv:hep-th/0306238} §7; geometric realisation
Iqbal--Kozcaz--Vafa \texttt{arXiv:hep-th/0701156}).

\begin{theorem}[Super-rigidity of $(-1,-1)$-curves]
\label{thm:super-rigidity}
\ClaimStatusProvedElsewhere\ \emph{(Bryan--Pandharipande
\texttt{arXiv:math/0009025} Thm.~1.2)}
A rigidly embedded $\bP^1 \subset X$ in a smooth Calabi--Yau threefold
$X$ with normal bundle $\cO(-1) \oplus \cO(-1)$ is
\emph{super-rigid}: it has no infinitesimal deformations as a map from
a smooth curve of any genus to $X$. Hence Gromov--Witten contributions
to all genera come only from multiple covers of the embedded $\bP^1$.
\end{theorem}

The multiple-cover Hodge-integral formula for super-rigid $\bP^1$
(Faber--Pandharipande \texttt{arXiv:math/9810173} for the multiple-cover
integrals; Pandharipande~\texttt{arXiv:math/9902107} for degenerate
contributions; Bryan--Pandharipande~\texttt{arXiv:math/0412005} for the
local Gromov--Witten theory of curves that assembles these integrals
into the toric vertex) combined with the M-theoretic BPS-counting
resummation of Gopakumar--Vafa~\texttt{arXiv:hep-th/9809187} Part~I
eqs.~(2.7)--(2.10) and Part~II~\texttt{arXiv:hep-th/9812127} for the
higher-genus structure, yields the BPS per-genus ansatz
\[
  F^{\mathrm{GV}}(\bY;\lambda,Q)
  \;=\;
  \sum_{g\ge 0}\sum_{\beta\in H_2(\bY,\bZ)}\sum_{k\ge 1}
    n^{g}_{\beta}\,\frac{1}{k}\,
    \bigl(2\sinh(k\lambda/2)\bigr)^{2g-2}\,
    Q^{k\beta},
\]
where $\lambda = g_s$ is the string coupling, $Q^\beta = e^{-\int_\beta
\omega}$ the Kähler-modulus character, and $n^g_\beta \in \bZ_{\ge 0}$
the integer Gopakumar--Vafa invariants counting BPS states with
$(J^3_L)$-spin $g$ in charge class~$\beta$. On the conifold the effective
cone is generated by $\beta = [\bP^1]$ and the only non-vanishing
invariant is $n^{0}_{[\bP^1]} = 1$ with $n^{g}_{d[\bP^1]} = 0$ for
$(g,d) \neq (0,1)$, which collapses the ansatz to the
Schwinger-resummed form
\[
  F^{\mathrm{GV}}(\bY; g_s, Q)
  \;=\;
  \sum_{k \geq 1}
    \frac{Q^{k}}{k\,\bigl(2\sinh(k g_s/2)\bigr)^{2}}.
\]
The genus expansion $F^{\mathrm{GV}} = \sum_{g \geq 0}
g_s^{2g-2} F_g(Q)$ follows from the formal power-series identity
$\frac{x^{2}}{(2\sinh(x/2))^{2}} = \sum_{g \geq 0}
\tfrac{(1-2g)\,B_{2g}}{(2g)!}\,x^{2g}$, giving
\[
  F_{0}(Q) \;=\; \mathrm{Li}_{3}(Q),
  \qquad
  F_{g}(Q) \;=\; \tfrac{(-1)^{g-1}\,B_{2g}}{2g\,(2g-2)!}\,
                \mathrm{Li}_{3-2g}(Q)
  \qquad (g \geq 2),
\]
recorded in Mariño \texttt{arXiv:hep-th/0406005} eq.~(2.68). The
$g = 1$ contribution is the $x^{0}$-term of
$(2\sinh(x/2))^{-2} = x^{-2} - \tfrac{1}{12} + \tfrac{x^{2}}{240} - \cdots$
applied to $\mathrm{Li}_{1}(Q) = -\log(1 - Q)$, giving
$F_{1}(Q) = -\tfrac{1}{12}\mathrm{Li}_{1}(Q) = \tfrac{1}{12}\log(1 - Q)$
(holomorphic-anomaly constant absorbed). The
$(-1)^{g-1}$ prefactor multiplies the Bernoulli number $B_{2g}$, whose
own sign is $(-1)^{g+1}$ for $g \geq 1$; the product is uniformly
negative for $g \geq 2$, matching coefficient-extraction from the
Schwinger form.

\begin{theorem}[Szendrői NCDT partition function, two chambers distinguished]
\label{thm:szendroi-ncdt}
\ClaimStatusProvedElsewhere\ \emph{(Szendrői \texttt{arXiv:0705.3419}
Thm.~2.6 for the chamber-specific NC result; Morrison--Mozgovoy--Nagao--
Szendrői \texttt{arXiv:1107.5017} Thm.~2 for the two-variable refinement;
Young \texttt{arXiv:0802.3948} for pyramid partitions.)}
The curve-less Szendrői theorem gives the NCDT partition function
indexed by pyramid partitions
\[
  Z^{\mathrm{NCDT}}(\bY; q)
  \;=\; \sum_{\pi \in \mathrm{Pyr}} (-q)^{|\pi|}
  \;=\; M(-q)^2 \cdot \prod_{k\ge1}(1-(-q)^{k})^{-2}
  \qquad \text{(Szendrői chamber, single parameter).}
\]
Passing to two variables $(q,Q)$ adapted to the $D^b(\bY)$-tilt
reorganisation, the full two-variable NCDT partition function is
\[
  Z^{\mathrm{NCDT}}(\bY;q,Q)
  \;=\; M(-q)^2 \cdot
    \prod_{k\ge 1}(1 - Q (-q)^k)^{k}
    \prod_{k\ge 1}(1 - Q^{-1}(-q)^k)^{k},
\]
i.e.\ \emph{both} $Q$ and $Q^{-1}$ factors are present in the Szendrői
(non-commutative) chamber (MMNS \texttt{arXiv:1107.5017}, eq.~(1.1)).
The DT (commutative) chamber partition function is the truncation
\[
  Z^{\mathrm{DT}}(\bY; q, Q)
  \;=\; M(-q)^2 \prod_{k \geq 1}(1 - Q(-q)^k)^k,
\]
obtained by absorbing the $\prod_k (1 - Q^{-1}(-q)^k)^k$ factor into the
vacuum sector via a single Kontsevich--Soibelman wall-crossing operator
$\mathcal T_{\gamma_\alpha}$ (Nagao--Nakajima
\texttt{arXiv:0809.2992} for the conifold perverse-coherent wall-crossing
framework; Nagao \texttt{arXiv:0809.2994} for the small-toric CY$_3$
chamber structure specialised to $\bY$;
cross-reference \texttt{chapters/examples/coha\_wall\_crossing\_platonic.tex}
lines~990--995). Conflating the DT truncation with the full NC
partition function is a common error; the Szendrői (NC) chamber
carries both $Q$ and $Q^{-1}$ infinite products, whereas the DT chamber
retains only the $Q$ product.
\end{theorem}

\begin{proposition}[Szendrői NCDT coincides with Bridgeland NCDT after
wall-crossing]
\label{prop:szendroi-bridgeland-coincidence}
\ClaimStatusProvedElsewhere\ \emph{(Nagao--Nakajima
\texttt{arXiv:0809.2992} for the categorical wall-crossing on the
conifold; Bridgeland \texttt{arXiv:1002.4372} Thm.~6.3 for the motivic
Hall-algebra framework; MMNS \texttt{arXiv:1107.5017} Thm.~2 for the
two-variable match.)}
The Szendrői NCDT generating series $Z^{\mathrm{NCDT}}(\bY;q,Q)$, defined
from pyramid partitions for the Jacobi algebra $J(Q_{\mathrm{con}},
W_{\mathrm{con}})$, equals the image of the Bridgeland NCDT generating
series — the motivic Donaldson--Thomas invariant of the motivic Hall
algebra of $D^b(J(Q_{\mathrm{con}}, W_{\mathrm{con}})\text{-mod})$ at
the tilt-Noetherian heart $\cA_{\mathrm{NC}}$ — under the chain of
Kontsevich--Soibelman wall-crossing automorphisms carrying
$\cA_{\mathrm{NC}}$ to the large-volume heart $\mathrm{Coh}(\bY)$. The
composite wall-crossing factor is $\prod_k (1 - Q^{-1}(-q)^k)^{k}$,
contributing precisely the extra infinite product that distinguishes
the Szendrői (NC) chamber from the DT (commutative) chamber, so that
on the DT side the Szendrői and Bridgeland generating functions
coincide as rational functions in $(q, Q)$.
\end{proposition}
\begin{proof}
The categorical wall-crossing is the identity of hearts
$\cA_{\mathrm{NC}} \to \mathrm{Coh}(\bY)$ along the Nagao--Nakajima
perverse-coherent chamber sequence (one wall per Seiberg
mutation~$s_0, s_1$). Each wall-crossing automorphism is the KS
operator $\mathcal T_\gamma = \exp\bigl(\sum_k \Omega(k\gamma)
[k\gamma]/k^2\bigr)$ with $\Omega(\gamma)$ the refined DT invariant
of the simple $\gamma$. Summing the generating series under these
automorphisms and applying the Bridgeland rationality theorem
\texttt{arXiv:1002.4372}~Thm.~6.3 to $D^b(\bY)$ converts
$Z^{\mathrm{NCDT}}$ to $Z^{\mathrm{DT}}$ by absorbing precisely the
$\prod_k (1 - Q^{-1}(-q)^k)^k$ factor; the identification on both
sides is the content of MMNS \texttt{arXiv:1107.5017}~Thm.~2.
\end{proof}

Under $q = -e^{i g_s}$ the DT-chamber expression matches $F^{\mathrm{GV}}$
via the MNOP triangle of partition-function identities
\[
\begin{array}{c@{\qquad}c@{\qquad}c}
 & Z^{\mathrm{GW}}(\bY; g_s, Q) & \\
 & \big\Updownarrow^{\text{GW/DT}} & \\
Z^{\mathrm{DT}}(\bY; q, Q) & \xleftrightarrow{\text{DT/PT}} &
 Z^{\mathrm{PT}}(\bY; q, Q)
\end{array}
\]
with $q = -e^{i g_s}$ and a common reduced free energy
$F^{\mathrm{GV}}$ on the vertices. The three edges are proved
independently:
\begin{itemize}[leftmargin=2em]
\item \emph{GW/DT (top edge)}: Maulik--Nekrasov--Okounkov--Pandharipande~I
  \texttt{arXiv:math/0312059} Conj.~3R, extended in
  MNOP~II~\texttt{arXiv:math/0406092} to the equivariant and reduced
  theory; the toric-CY$_3$ case is proved for $\bY$ via the local-curve
  primitive Bryan--Pandharipande \texttt{arXiv:math/0412005} (the
  $(-1,-1)$-cap computation that feeds the toric vertex), the full
  toric proof being completed in Maulik--Oblomkov--Okounkov--Pandharipande
  \texttt{arXiv:0809.3976}. The 2005 Bryan--Pandharipande paper
  \texttt{arXiv:math/0412005} is the \emph{local Gromov--Witten primitive}
  feeding MNOP; it is distinct from the 2001 super-rigidity primitive
  Bryan--Pandharipande~\texttt{arXiv:math/0009025} of
  Theorem~\ref{thm:super-rigidity}, and from the 1999 Hodge-integral
  input Pandharipande~\texttt{arXiv:math/9902107} feeding the
  multiple-cover formula.
\item \emph{DT/PT (bottom edge)}: Pandharipande--Thomas
  \texttt{arXiv:0707.2348} define the stable-pair moduli; Bridgeland
  \texttt{arXiv:1002.4372} proves the DT/PT identity rationally via
  motivic Hall algebras; Toda \texttt{arXiv:0902.4371} gives the
  motivic DT/PT equivalence at refined level.
\item \emph{GW/PT (right edge)}: composite of the other two, consistent
  with the MNOP~I rationality conjecture on toric~$\bY$.
\end{itemize}

\begin{remark}[Refined topological vertex: primary-source conventions]
\label{rem:ikv-refined-conventions}
\ClaimStatusConventionDependent\ \emph{(Iqbal--Kozcaz--Vafa
\texttt{arXiv:hep-th/0701156}, eqs.~(4.2)--(4.3); independent
Awata--Kanno \texttt{arXiv:hep-th/0502061} derivation.)}
Iqbal--Kozcaz--Vafa present the refined $(-1,-1)$-conifold partition
function as
\[
  Z^{\mathrm{ref}}_{\mathrm{top}}(\bY; t, q, Q)
  \;=\; \prod_{i, j \geq 1}
    \bigl(1 - Q\, t^{\,i - 1/2}\, q^{\,j - 1/2}\bigr),
\]
with the half-integer offsets $(i-1/2, j-1/2)$ a \emph{convention
choice} — Awata--Kanno's equivalent form uses integer offsets
$(i, j)\mapsto(i-1,j-1)$ with a prefactor $Q^{-1/2}$ absorbed into the
Kähler modulus. The half-integer form is the one in IKV~eq.~(4.2) and
is the form that makes the single-leg refined vertex
$P_{\lambda\mu\nu}(q,t)$ Young-diagram-combinatorial. The sign
convention $q,t$ in IKV (\S3) corresponds to the flop-invariant
normalisation; the opposite-sign convention of Taki
\texttt{arXiv:hep-th/0703133} differs by $(q,t) \mapsto (q^{-1}, t^{-1})$
on the unrefined locus $q = t$. The identification $t = e^{i\varepsilon_1}$,
$q = e^{i\varepsilon_2}$ with Nekrasov's $\Omega$-background requires the
following primary-source correspondence to be recorded:
\emph{Nekrasov~2002}~\texttt{hep-th/0206161} introduced the
one-parameter $\Omega$-background on $\bC^2 \times S^1$; the
\emph{refined} two-parameter extension to $(\varepsilon_1,\varepsilon_2)$
for gauge theories was Nekrasov--Okounkov~2003
\texttt{hep-th/0306238}, §7; the refined \emph{topological vertex} was
IKV~\texttt{hep-th/0701156}; the MNOP GW/DT correspondence
\texttt{math/0312059} is unrefined. Any claim ``Nekrasov partition
function for conifold in $\Omega$-background'' must cite
\emph{Nekrasov--Okounkov 2003} for the two-parameter refinement, not
Nekrasov 2002 (one-parameter) and not MNOP (unrefined).
\end{remark}

\begin{remark}[Refined free energy: Nekrasov--Okounkov Schwinger form]
\label{rem:refined-free-energy-NO}
\ClaimStatusProvedElsewhere\ \emph{(Nekrasov--Okounkov
\texttt{arXiv:hep-th/0306238} \S7; IKV \texttt{arXiv:hep-th/0701156}
\S4; Nakajima--Yoshioka \texttt{arXiv:0905.0897} Thm.~1.1 for the
K-theoretic gauge-theory / refined-DT bridge.)}
Taking the logarithm of the IKV product yields the refined conifold
free energy
\[
  F^{\mathrm{ref}}(\bY;\varepsilon_{1},\varepsilon_{2},Q)
  \;=\;
  -\sum_{i,j \geq 1} \log\!\bigl(1 - Q\,t^{\,i-1/2}q^{\,j-1/2}\bigr)
  \;=\;
  \sum_{k \geq 1}
    \frac{Q^{k}}{k\,
      \bigl(2\sinh(k\varepsilon_{1}/2)\bigr)\,
      \bigl(2\sinh(k\varepsilon_{2}/2)\bigr)},
\]
with $t = e^{i\varepsilon_{1}}$, $q = e^{i\varepsilon_{2}}$. Under the
unrefined specialisation $\varepsilon_{1} = \varepsilon_{2} = g_{s}$
this collapses to
$F^{\mathrm{GV}}(\bY; g_s, Q) = \sum_{k \geq 1} Q^{k}/[k(2\sinh(k g_s/2))^{2}]$,
matching the Schwinger form of the preceding display. The two-sinh
denominator records the $\Omega$-deformed $\bC^{2}_{\varepsilon_{1},
\varepsilon_{2}}$ equivariant Hilbert series at the origin.
Nakajima--Yoshioka \texttt{arXiv:0905.0897} Thm.~1.1 proves the
K-theoretic version for moduli of framed sheaves on $\bP^{2}$,
confirming the same two-parameter Schwinger structure on the
gauge-theory side.
\end{remark}

\begin{conjecture}[Algebraic super-rigidity of $Y^{+}(\widehat{\fgl}(1|1))$
under $(\varepsilon_{1},\varepsilon_{2})$-deformation]
\label{conj:algebraic-super-rigidity}
\ClaimStatusConjectured\ \emph{(first stated here as a conjecture;
unrefined $\widehat{\fgl}_{1}$-cousin is the Schiffmann--Vasserot
rigidity result~\texttt{arXiv:1202.2756}; the
$\varepsilon_{+}$-gauge mechanism traces to Nekrasov--Okounkov
\texttt{arXiv:hep-th/0306238} §7; the 5D hCS super obstruction
enters via the three-author Costello--Gaiotto--Yagi
\texttt{arXiv:2103.01835}, not the two-author Costello--Yagi
\texttt{arXiv:1810.01970}.)}
The chain-level deformation of the positive half
$Y^{+}(\widehat{\fgl}(1|1))$ by the two parameters
$(\varepsilon_{1},\varepsilon_{2})$ of the refined $\Omega$-background
is \emph{algebraically super-rigid}: the cohomology group
$\HH^{2}_{\mathrm{Hopf}}\bigl(Y^{+}(\widehat{\fgl}(1|1));\,
Y^{+}(\widehat{\fgl}(1|1))\bigr)$ governing graded
Hopf-super-bialgebra deformations is one-dimensional, spanned by the
class of the refined $\Omega$-background with first-order parameter
$\varepsilon_{1}\varepsilon_{2}$. The sum $\varepsilon_{+} :=
\varepsilon_{1}+\varepsilon_{2}$ is absorbed into the
Nekrasov--Okounkov $\varepsilon_{+}$-rescaling gauge equivalence
(\texttt{hep-th/0306238} §7, a one-parameter dilatation of the
Drinfeld current generators targeting the supertrace-degenerate
direction of the $\widehat{\fgl}(1|1)$ Killing form --- the same
degeneration that obstructs the 5D hCS super-extension in
Theorem~\ref{thm:hcs-super-open}), singling out
$\varepsilon_{1}\varepsilon_{2}$ as the unique physical first-order
deformation and reproducing the Nekrasov--Okounkov refined Schwinger
form of Remark~\ref{rem:refined-free-energy-NO} at the level of
super-characters.
Geometric super-rigidity (Theorem~\ref{thm:super-rigidity}) is a
\emph{distinct} statement: a geometric fact about deformations of the
embedded $\bP^{1}$, not an algebraic fact about the BPS algebra. The
two are linked only through the CoHA dictionary
$\CoHA(\bY) \simeq Y^{+}(\widehat{\fgl}(1|1))^{\mathrm{con}}$.
\end{conjecture}

\begin{remark}[Why the unrefined rigidity is Schiffmann--Vasserot and
the refined is open]
\label{rem:SV-rigidity-unrefined-only}
Schiffmann--Vasserot \texttt{arXiv:1202.2756} establish algebraic
rigidity of $Y^{+}(\widehat{\fgl}_{1})$ under the single-parameter
$\hbar$-deformation in the abelian (even) case. The super case under
the two-parameter $(\varepsilon_{1},\varepsilon_{2})$ deformation is
not covered: (i) Ginzburg--Schedler super formality
\texttt{arXiv:0807.0174} terminates at $E_{2}$ rather than $E_{3}$,
removing the Kontsevich--Tamarkin formality input that drives the SV
rigidity proof; (ii) the super Hopf-algebra cohomology
$\HH^{2}_{\mathrm{Hopf}}$ has not been computed for
$Y^{+}(\widehat{\fgl}(1|1))$. Conjecture~\ref{conj:algebraic-super-rigidity}
accordingly remains open.
\end{remark}

\begin{conjecture}[Refined $Z_{\mathrm{top}}$ as $Y(\widehat{\fgl}(1|1))$
super-character at level~$1$]
\label{conj:ikv-super-character}
\ClaimStatusConjectured\
Under $t = e^{i\varepsilon_1}$, $q = e^{i\varepsilon_2}$, the refined
IKV formula
$Z^{\mathrm{ref}}_{\mathrm{top}}(\bY; t, q, Q) = \prod_{i,j\ge 1}
(1 - Q\, t^{i-1/2} q^{j-1/2})$
equals, coefficient-by-coefficient, the Weyl--Kac super-character of the
level-one vacuum module of
$Y_{\varepsilon_1,\varepsilon_2}(\widehat{\fgl}(1|1))$.
\end{conjecture}

\begin{remark}[Coefficient-by-coefficient check at $Q^0, Q^1$]
\label{rem:ikv-super-character-coefficient-check}
\ClaimStatusPartialPrimaryLit\
Expanding the IKV product in $Q$:
\begin{itemize}[leftmargin=2em]
\item $Q^0$-coefficient: $1$, matching the super-Yangian vacuum vector.
\item $Q^1$-coefficient: $-\sum_{i,j\ge 1} t^{i-1/2} q^{j-1/2}
  = -\tfrac{t^{1/2} q^{1/2}}{(1-t)(1-q)}$, which is the refined Hilbert
  series of $\bC^2_{t,q}$ weighted by $-(tq)^{1/2}$. The super-Yangian
  side is the graded multiplicity of the simple-root affine generator
  $e^{(1)}_{\alpha_1}(z)$ on the level-$1$ vacuum; this has character
  $\tfrac{(tq)^{1/2}}{(1-t)(1-q)}$ under the $\bZ^2$-grading
  (Gaiotto--Rapčák \texttt{arXiv:1703.00982} §3.2 with the sign from
  fermionic generator). The two agree up to the overall $-$ sign.
\item $Q^{1/2}$-coefficient: \emph{absent} in the IKV product (the
  expansion is in integer powers of $Q$); a ``$Q^{1/2}$'' coefficient
  match would require identifying $Q$ in IKV with $Q^2$ in some
  super-character convention. The correct coefficient-by-coefficient
  statement is in $Q^0, Q^1, Q^2, \ldots$ (integer powers), with the
  fractional factor $(tq)^{1/2}$ appearing \emph{inside} each
  integer-$Q^n$ coefficient via the half-shift $t^{i-1/2} q^{j-1/2}$,
  rather than as a ``$Q^{1/2}$-term'' on the IKV side or as a
  ``$q^{1/2}$-coefficient'' in a $q$-expansion.
\item $Q^2$-coefficient: expected to match the graded multiplicity of
  the quadratic span
  $e^{(1)}_{\alpha_1} e^{(1)}_{\alpha_1}(z,w)$ on the vacuum, modulo
  fermionic Serre relations; explicit check is outstanding.
\end{itemize}
The full all-$Q^n$ match is \emph{conjectural}; only $Q^0, Q^1$ have a
direct first-principles verification in the synthesis corpus. The
fractional factor $(tq)^{1/2}$ appears inside every integer-$Q^n$
coefficient via the half-shift $(i-1/2, j-1/2)$; there is no
fractional power of $Q$ in the expansion.
\end{remark}

% ============================================================================
\section{Four $\kappa$-invariants on the conifold}
\label{sec:kappa-reconciliation}
% ============================================================================

\begin{remark}[Three numerically distinct values near the conifold:
subscript disambiguation]
\label{rem:kappa-conflict}
\ClaimStatusProposition\
Three numerically distinct values live in the $\kappa$-neighbourhood of
the conifold, one under each of three subscripts:
\begin{itemize}[leftmargin=2em]
\item $\kappa_{\mathrm{ch}}(\bY) = +1$ in \texttt{chapters/examples/toric\_cy3\_coha.tex}
  Cor.~\texttt{cor:conifold-shadow} (canonical toric chapter;
  identified there as $\mathrm{DT}_{(1,0)} = 1$, the single compact
  curve-class contribution, $\kappa_{\mathrm{ch,GV}}$).
\item A value $0$ appearing in
  \texttt{chapters/examples/derived\_categories\_cy.tex:853}
  which is $\kappa_{\mathrm{cat}} := \chi(\cO_{\bY}) = 0$ for the
  non-compact total space (Künneth/retraction to $\bP^1$), \emph{not}
  a second value of $\kappa_{\mathrm{ch}}$.
\item A value $-1$ appearing in
  \texttt{chapters/theory/e1\_chiral\_algebras.tex:2454}
  which is a $\kappa^!$-type (Verdier-dual) invariant or a bulk BV
  anomaly, not $\kappa_{\mathrm{ch}}$ itself.
\end{itemize}
Under the four-$\kappa$ discipline ($\kappa_{\mathrm{ch}}$,
$\kappa_{\mathrm{cat}}$, $\kappa_{\mathrm{BKM}}$, $\kappa_{\mathrm{fiber}}$,
never conflated), there is no inconsistency: the three numbers
live under three different subscripts and record three different
mathematical objects. The correct rectification is subscripted
disambiguation at every manuscript occurrence; there is no single
$\kappa$ decomposing into three pieces.

The ``$\kappa + \kappa^! \in \{0, 8, 13, 250/3, 98/3\}$ on the
canonical five-archetype $\mathsf{G}/\mathsf{L}/\mathsf{C}/\mathsf{M}/\mathsf{B}$
ceiling'' in Theorem~C (shared five-theorem core) is a
\emph{$\mathsf G$-archetype ceiling statement} about landmark
Calabi--Yau data with specified $\chi_{\mathrm{top}}/24$-regulated
centre-curving, \emph{not} a universal identity $\kappa_{\mathrm{ch}}(X)
+ \kappa_{\mathrm{ch}}^!(X) = 0$ on every toric CY$_3$. Reading off
$\kappa_{\mathrm{ch,bulk}} = -1$ from $\kappa_{\mathrm{ch,GV}} = +1$
via ``$\kappa + \kappa^! = 0$'' is \emph{not} a theorem; it is an
ansatz that would require the conifold to sit at the $\mathsf G$-row of
the Theorem~C ceiling, which it does (class~$\mathsf G$ with shadow
depth $r_{\max} = 2$, Cor.~\texttt{cor:conifold-shadow}) — but the
$\mathsf G$-row entry of the ceiling is $K^\kappa = 0$ for toric bases
(\texttt{toric\_cy3\_coha.tex} line~251:
$K^{\kappa_{\mathrm{ch}}}(\mathrm{conifold}) = 0$), not $K^\kappa = 8$.
The $8$-value is the $\mathsf B$-row Mukai-enhanced K3 Heisenberg
witness on $K3 \times E$, not the $\mathsf G$-row conifold entry.
\end{remark}

% ============================================================================
\section{Transfer to toric CY$_3$, obstructions for compact CY$_3$}
\label{sec:compact-obstruction}
% ============================================================================

\begin{proposition}[Transfer to toric CY$_3$]
\label{prop:toric-transfer}
\ClaimStatusProposition\
The five-route construction of $\CoHA(\bY)$ extends to any local
toric Calabi--Yau threefold $X$ without compact $4$-cycles, with the
following adjustments:
\begin{enumerate}[leftmargin=2em]
\item The toric atlas replaces the two-chart cover with a multi-chart
  cover indexed by the toric fan; each chart remains $\bC^3$ with a
  Jordan-triple $W_{\mathrm{chart}}$.
\item The Klebanov--Witten quiver is replaced by the appropriate
  toric quiver with potential (e.g., the Beilinson quiver for local
  $\bP^2$; the $\bC^3/\bZ_n$ McKay quiver for local $A_{n-1}$).
\item The bond factors $\varphi_{a \Rightarrow b}$ are read off the
  toric fan via Galakhov--Li--Yamazaki \texttt{arXiv:2106.01230}.
\item Davison--Meinhardt integrality + PBW gives the BPS Lie algebra.
\item The Drinfeld double is the appropriate affine (super-)Yangian
  of the root system determined by the quiver.
\end{enumerate}
This covers local $\bP^2$, $\bC^3/\bZ_n$ for small $n$, banana
threefolds, and other non-compact toric CY$_3$ without compact
$4$-cycles.
\end{proposition}

\begin{theorem}[Sharp obstruction on compact CY$_3$; $2.5$-count reduction]
\label{thm:compact-obstruction}
\ClaimStatusProposition\
The five-route construction does \emph{not} transfer to compact
Calabi--Yau threefolds. Four named obstructions --- (O1) toric
fan-completeness, (O2) BCOV $\alpha_{\mathrm{BCOV}}$, (O3)
$\mathrm{Aut}^0$ rigidity, (O4) finite quiver with potential --- reduce
logically to \emph{$2.5$ effective obstructions} via
$\textup{(O1)} \Leftrightarrow \textup{(O3)}$ and $\textup{(O1)}
\Rightarrow \textup{(O4)}$, with $\textup{(O2)}$ the single
independent BV-cohomological obstacle:
\begin{equation}
\label{eq:obstruction-reduction-diagram}
\big(\textup{O1} \;\Longleftrightarrow\; \textup{O3}\big)
\;\Longrightarrow\;
\textup{O4},
\qquad
\textup{O2} \;\textup{ logically independent.}
\end{equation}
The effective count is $1$ (O1$\equiv$O3) $+$ $1$ (O2) $+$ $0.5$ (O4
downstream) $=$ $2.5$, not $4$.

\smallskip
\emph{\textup{(O1)} Failure of toric fan-completeness on a compact
ambient.}
A toric atlas on a smooth threefold $X$ consists of a complete
rational polyhedral fan $\Delta$ in $N_{\bR} \cong \bR^3$ whose
maximal cones $\sigma \in \Delta(3)$ are each generated by $3$
primitive ray vectors $(v_1, v_2, v_3)$ with $\det(v_1, v_2, v_3) =
\pm 1$ (smoothness), subject to the Calabi--Yau slice condition:
existence of a piecewise-linear support function $\psi: N_{\bR} \to
\bR$ with $\psi(v_i) = 1$ for every ray $v_i$
(Aspinwall--Greene--Morrison \texttt{hep-th/9309097} \S 3; equivalent
to $c_1(X) = 0$ via the vanishing sum of ray vectors in each maximal
cone). The simultaneous-slice condition $\psi(v_i) = 1$ forces every
ray to lie on a common affine hyperplane $\{\psi = 1\} \subset
N_{\bR}$, so $\Delta$ is a cone over a $2$-dimensional lattice
polytope and has no compact $3$-dimensional interior. The
Demazure--Fulton criterion (\emph{Introduction to Toric Varieties},
\S 3.4) yields: no smooth compact toric CY$_3$ exists; the only
complete CY-sliced toric threefolds are finite quotients of $\bC^3$.
Compact CY$_3$ (the quintic, $K3 \times E$, bihypersurfaces) admit
no chart-wise $(\bC^3, W_{\mathrm{chart}})$ cover, so the
Schiffmann--Vasserot per-chart $Y^+(\widehat{\fgl}_1)$ stalks have no
geometric location on $X$; the conifold's \v{C}ech route of
Section~\ref{sec:cech-route} is unavailable.

\smallskip
\emph{\textup{(O2)} BCOV $1$-loop anomaly: independent BV
cohomological obstacle.}
For a holomorphic Chern--Simons factorisation-algebra quantisation on
a CY$_3$ $X$, Costello--Li 2016 \texttt{arXiv:1606.00365} Prop.~5.2
identify the unique $1$-loop anomaly class
\[
  \alpha_{\mathrm{BCOV}}(X)
  \;=\; \frac{\chi(X)}{24}\,[\Omega_X]^{0,1}
  \;\in\; H^{0,1}\!\bigl(X, \mathrm{Sym}^{\leq 2} T_X^*\bigr),
\]
sourced by the BV $1$-loop counter-term Feynman graph weighted by
$\Omega_X$ and the Todd specialisation $\chi(X)/24$. On $\bY$ the
class vanishes by retraction $\bY \simeq \bP^1$ and $H^{0,1}(\bP^1)
= 0$ (\emph{cf}.\ Remark~\ref{rem:anomaly-vanishing}). On a compact
CY$_3$ with $\chi(X) \neq 0$ and $[\Omega_X]^{0,1}$ non-vanishing in
Dolbeault cohomology, $\alpha_{\mathrm{BCOV}}(X) \neq 0$: no
factorisation algebra of hCS observables exists at $1$-loop. This
obstruction sources from BV cohomology, \emph{not} from automorphism
geometry; it is logically disjoint from (O1).

\smallskip
\emph{\textup{(O3)} $\mathrm{Aut}^0$-rigidity, equivalent to
\textup{(O1)}.}
The Bogomolov decomposition theorem (Bogomolov 1974; Beauville 1983
\emph{JDG} 18:755 \S 2) for a compact K\"ahler manifold $X$ with
$c_1(X) = 0$ gives a finite \'etale cover
\[
  \widetilde X \;=\; T^{2m} \times \prod_i S_i \times \prod_j Y_j,
\]
with $T^{2m}$ a complex torus, $S_i$ irreducible hyperk\"ahler, and
$Y_j$ strict Calabi--Yau ($h^{1,0}(Y_j) = 0$, $H^0(Y_j,
\Omega^2_{Y_j}) = 0$). For a compact CY$_3$ of strict type (quintic,
bihypersurfaces), $h^{1,0}(X) = 0$ and $h^{2,0}(X) = 0$ force
$\mathrm{Aut}^0(X)$ to act trivially on the Hodge decomposition;
Matsumura 1963 \emph{Proc.\ Japan Acad.} 39:181 Thm.~1 then concludes
\[
  \mathrm{Aut}^0(X) \;=\; \{1\}
  \qquad (\text{compact CY$_3$ of strict type}).
\]
For $X = K3 \times E$ of product type (stratum (ii) in
Remark~\ref{rem:four-strata}), $\mathrm{Aut}^0(X) = E$ acts via
translations on the elliptic factor; there is no $(\bC^\times)^3$
torus action lifting the chart-wise conifold $T^2$-equivariance. In
either case, the conifold's $T^2$-fixed pyramid-partition
localisation (Szendr\H{o}i \texttt{arXiv:0705.3419} Thm.~2.6) has no
analogue.

The equivalence $\textup{(O1)} \Leftrightarrow \textup{(O3)}$: a
complete toric fan $\Delta$ for $X$ generates a dense torus
$(\bC^\times)^3 \subset \mathrm{Aut}^0(X)$, hence $\mathrm{Aut}^0(X)
\supsetneq \{1\}$ on a positive-dimensional stratum; conversely, a
positive-dimensional $\mathrm{Aut}^0(X)$ on a strict compact CY$_3$
is a complex torus by Bogomolov--Matsumura (the only connected Lie
groups preserving a CY$_3$ structure with $h^{1,0} = 0$ are
abelian), and a complex torus acting faithfully with a dense orbit
on a normal variety forces a toric structure (Sumihiro 1974
\emph{J.\ Math.\ Kyoto Univ.} 14:1 Thm.~1). The two failures ---
absence of a toric fan, absence of a non-trivial connected
automorphism group --- are two facets of the same
Bogomolov--Matsumura rigidity.

\smallskip
\emph{\textup{(O4)} Finite quiver with potential, downstream from
\textup{(O1)}.}
Van den Bergh's NCCR construction (\texttt{arXiv:math/0211064}
Def.~4.1, Prop.~3.2.10) requires a tilting bundle $T \in
D^b(\mathrm{Coh}\,X)$ with $\End_X(T)$ a finite-rank algebra over
the Noetherian coordinate ring $R = \cO_X(X)$. On a local toric
CY$_3$ without compact $4$-cycles, $R$ is a finitely-generated
$\bC$-algebra and the Bocklandt regularity criterion
(\texttt{arXiv:math/0603558} Thm.~3.1: NCCR $\Leftrightarrow$ the
superpotential quiver is Calabi--Yau-$3$ in the Ginzburg sense and
every arrow appears in $W$) promotes $\End_X(T)$ to a
quiver-with-potential $J(Q, W)$ with finitely many vertices and
finitely many arrows, hence finite-dimensional per-vertex CoHA
stalks $Y^+(\widehat{\fgl}_1)$. On a compact CY$_3$, Hartogs on the
projective base gives $R = \cO_X(X) = \bC$, so no finite-rank
$R$-algebra tilting-endomorphism presentation can exist;
$D^b(\mathrm{Coh}\,X)$ is equivalent to no $(Q, W)$-Jacobi category.
This is the contrapositive of (O1): \emph{no toric atlas
$\Rightarrow$ no finite chart-wise quiver-with-potential
$\Rightarrow$ no chart-wise finite-arrow CoHA}. Hence
$\textup{(O1)} \Rightarrow \textup{(O4)}$, and the automorphic-form
replacement (Borcherds products, Gritsenko--Nikulin $\Delta_5$,
Igusa $\Phi_{10}$) enters via the non-chart-wise
$(\infty,1)$-factorisation-homology route of
Remark~\ref{rem:compact-transfer}.

\smallskip
The five-route chart-wise programme terminates at the toric boundary
for the single independent BV-cohomological reason (O2), the single
independent geometric reason (O1 $\equiv$ O3), and the downstream
representation-theoretic corollary (O4). Any compact CY$_3$ without
a $K3 \times E$-type product fibration admits no chart-wise
$Y^+(\widehat{\fgl}_1)$ gluing; the replacement on $K3 \times E$
(Remark~\ref{rem:mnop-replacement}) is non-chart-wise Davison
integrality with Maulik--Okounkov stable-envelope wall-crossing.
\end{theorem}

\begin{remark}[Scope of the $2.5$-count]
\label{rem:obstruction-2p5-scope}
Theorem~\ref{thm:compact-obstruction} replaces any ``four
independent obstructions'' phrasing with the $2.5$-count derivation
\eqref{eq:obstruction-reduction-diagram}: (O1 $\equiv$ O3) is a
single Bogomolov--Matsumura rigidity stated twice (toric-side and
$\mathrm{Aut}^0$-side); (O4) is the downstream finite-quiver
corollary via Van den Bergh tilting; (O2) is the logically
independent BCOV $1$-loop obstacle. The equivalence (O1)
$\Leftrightarrow$ (O3) is \emph{not} a coincidence of two
independent cohomological tests; it is the twofold reading of one
theorem (Bogomolov 1974 + Matsumura 1963 + Sumihiro 1974) --- once as
``no complete CY-sliced fan'', once as ``$\mathrm{Aut}^0 = \{1\}$ on
strict CY$_3$''. The BCOV class $\alpha_{\mathrm{BCOV}} =
(\chi/24)[\Omega_X]^{0,1}$ is independent because the Todd
specialisation $\chi/24$ depends only on topological Euler data and
$[\Omega_X]^{0,1}$ is a Dolbeault-cohomological representative
unrelated to the fan structure.
\end{remark}

\begin{remark}[What does transfer]
\label{rem:compact-transfer}
Three ingredients transfer to any CY$_3$ category:
\begin{itemize}[leftmargin=2em]
\item The Davison--Meinhardt critical CoHA + integrality PBW applies to
  any CY$_3$ category with a potential;
\item Kontsevich--Soibelman motivic wall-crossing applies generally;
\item The two-stage factorisation $\Phi_3 = \mathrm{Sp}^{\mathrm{ch}}_{\Sigma_2, C}
  \circ \Phi^{FA}_3$ is geometry-independent in principle.
\end{itemize}
What fails is the \emph{explicit chart-wise formulas}: the BPS Lie
algebra exists on any CY$_3$, but its presentation for compact $X$ is
not chart-wise.
\end{remark}

\begin{remark}[The four equivariance strata]
\label{rem:four-strata}
The four equivariance strata (CLAUDE.md):
\begin{enumerate}[label=\emph{(\roman*)}, leftmargin=2em]
\item \emph{Toric $T^d$} (conifold, local $\bP^2$, $\bC^3$): chart-wise
  shuffle-algebra presentation, the setting of this paper;
\item \emph{Reduced $\bC^\times \times \mathrm{Aut}(X)$} (K3,
  $K3 \times E$, abelian surface): automorphism-enriched presentation,
  still equivariant but no toric atlas;
\item \emph{Orbifold inertia $I(X/G)$} (Mathieu $M_{24}$, McKay
  $\Gamma \subset SU(d)$): equivariant cohomology enriched by the
  inertia stack;
\item \emph{Lattice-polarised period domain} (Borcherds lifts,
  Gritsenko $\Delta_5$, Igusa $\Phi_{10}$): automorphic-form
  presentation with no chart-wise shuffle analogue.
\end{enumerate}
The conifold lives in stratum (i); compact CY$_3$ typically live in
strata (iii) or (iv); $K3$ and $K3 \times E$ straddle strata (ii)--(iv)
depending on the Aut/inertia/Borcherds-lift chosen.
\end{remark}

\begin{remark}[Replacement machinery on compact CY$_3$]
\label{rem:mnop-replacement}
For compact CY$_3$ the CoHA is built via Maulik--Nekrasov--Okounkov--Pandharipande
categorified machinery combined with Kontsevich--Soibelman motivic DT
on the non-toric side. The $K3 \times E$ case (treated in Vol III Part~IV
of this programme) exemplifies the replacement: instead of two-chart
$\bC^3$ pieces, one applies Davison integrality directly to
$D^b(\mathrm{Coh}(K3 \times E))$, producing a Hall--Drinfeld double with
imaginary rank-$23$ Cartan (the rank-$24$ Mukai lattice $\widetilde H(K3,\Z)
\simeq \mathrm{II}_{4,20} = H^0 \oplus H^2 \oplus H^4$ of dimension $1 + 22 + 1 = 24$,
modulo the one-dimensional null direction spanned by the
$\delta = [H^0] - [H^4]$ imaginary class corresponding to the
$E$-fibre point/hyperplane exchange; see
\texttt{chapters/examples/k3\_chiral\_bialgebra\_platonic.tex} and
\texttt{chapters/examples/toric\_cy3\_coha.tex} for the
$\mathrm{sign}\,(2,1)$ Cartan of $\fg_{\Delta_5}$), Humbert-surface monodromy
of order $8$, and
$\Delta_5$-associator data — all non-toric invariants invisible to the
chart-wise machinery. Chart-wise $Y^+(\widehat{\fgl}_1)$ stalks appear
only at the $24$ smooth-$\bC^3$ points of a nodal-$K3 \times E$
degeneration, matching the $24$ $I_1$-fibres of the elliptic fibration on
$K3$.
\end{remark}

% ============================================================================
\section{Open frontiers}
\label{sec:open-frontiers}
% ============================================================================

Four primary frontiers $\mathsf F_{1}$--$\mathsf F_{4}$ govern the
conifold programme; four secondary items $\mathsf F_{5}$--$\mathsf F_{8}$
record auxiliary loci where the construction of §§3--14 admits
sharpening. Each item carries a claim-status marker and an
\emph{explicit falsifier}: a first-principles computation whose outcome
would disprove the item. Falsifiability is the Popperian standard
demanded by Beilinson's dictum; a conjecture that cannot be killed has
no content. The items below each name the experiment that would kill
them.

\begin{enumerate}[label=\emph{$\mathsf F_{\arabic*}$.}, leftmargin=2.4em]
\item \emph{Algebraic super-rigidity of $Y^{+}(\widehat{\fgl}(1|1))$.}
  \ClaimStatusConjectured\
  Conjecture~\ref{conj:algebraic-super-rigidity}: the
  Hopf-super-bialgebra cohomology
  $\HH^{2}_{\mathrm{Hopf}}(Y^{+}(\widehat{\fgl}(1|1));\,
  Y^{+}(\widehat{\fgl}(1|1)))$ is one-dimensional, spanned by the
  class of the $\varepsilon_{1}\varepsilon_{2}$-deformation.
  \emph{Falsifier.} Exhibit a $2$-cocycle
  $\omega\in Z^{2}_{\mathrm{Hopf}}$ on $Y^{+}(\widehat{\fgl}(1|1))$
  linearly independent of $[\varepsilon_{1}\varepsilon_{2}]$ modulo
  the $(\varepsilon_{1}+\varepsilon_{2})$-gauge equivalence of
  Remark~\ref{rem:SV-rigidity-unrefined-only}; equivalently, construct
  a formal deformation of the super-shuffle
  product~\eqref{eq:bond-factors} whose first-order class is not
  spanned by the IKV Schwinger kernel of
  Remark~\ref{rem:refined-free-energy-NO}. Primary obstruction to the
  positive proof: Ginzburg--Schedler super-formality
  \texttt{arXiv:0807.0174} terminates at $E_{2}$, removing the
  Kontsevich--Tamarkin $E_{3}$-input that powered unrefined
  Schiffmann--Vasserot \texttt{arXiv:1202.2756}.

\item \emph{Refined $Z_{\mathrm{top}}$ as super-character at level~$1$.}
  \ClaimStatusConjectured\
  Conjecture~\ref{conj:ikv-super-character}:
  $Z^{\mathrm{ref}}_{\mathrm{top}}(\bY;t,q,Q)
   = \chi^{\mathrm{Weyl{-}Kac}}_{
     Y_{\varepsilon_{1},\varepsilon_{2}}(\widehat{\fgl}(1|1)),\,\mathrm{lev}\,1}$
  as formal power series in $Q$, with $t = e^{i\varepsilon_{1}}$,
  $q = e^{i\varepsilon_{2}}$. Coefficients $Q^{0}, Q^{1}$ verified
  directly (Remark~\ref{rem:ikv-super-character-coefficient-check});
  $Q^{\geq 2}$ open. \emph{Falsifier.} Compute the $Q^{2}$-coefficient
  of the IKV product in closed form and compare to the graded
  multiplicity
  $\mathrm{ch}_{q,t}\bigl(e^{(1)}_{\alpha_{1}}(z)\,
   e^{(1)}_{\alpha_{1}}(w)\cdot|\mathrm{vac}\rangle\bigr)$
  in $Y_{\varepsilon_{1},\varepsilon_{2}}(\widehat{\fgl}(1|1))$ modulo
  fermionic quadratic Serre; any mismatch not absorbed into the
  $(tq)^{1/2}$ half-shift normalisation of
  Remark~\ref{rem:ikv-super-character-coefficient-check} falsifies
  $\mathsf F_{2}$.

\item \emph{Super-Yangian Hopf pairing asymmetric correction.}
  \ClaimStatusConjectured\
  Conjecture~\ref{conj:hopf-pairing}: the subleading coefficient on
  the odd diagonal is $c_{11} = -1$ in the Nazarov convention
  (equivalently $c_{11} = +1$ in the Gow--Peng
  antipode-in-coproduct convention), with $c_{00} = 0$ on the even
  diagonal. Leading binomial kernel verified via three independent
  routes (Maulik--Okounkov stable-envelope residue;
  Soibelman--Gelfand refined Ext-pairing $\mathbf E_{\hbar}$;
  Tsymbaliuk shuffle \texttt{arXiv:1404.5240} Prop.~4.3).
  \emph{Falsifier.} Perform the Berezinian-residue computation
  $\mathrm{coeff}_{u^{-m-1}v^{-n-1}}\!
   \bigl[\log\mathrm{Ber}_{q}T(u)\cdot k^{\mathrm{off}}(u-v)\bigr]$
  with
  $k^{\mathrm{off}}(u)
   = (u+\varepsilon_{1})(u+\varepsilon_{2})/
     [u(u+\varepsilon_{1}+\varepsilon_{2})]$
  in the Nazarov~1991 / Gow~\texttt{arXiv:math/0411166} / Peng~2011
  framework. A value $|c_{11}|\neq 1$ falsifies the
  parity-isotropy-source argument; a non-zero $c_{00}$ falsifies the
  even-self-Serre input on the Cartan diagonal.

\item \emph{M-theory parent and 5D hCS super all-orders.}
  \ClaimStatusConjectured\
  Conjecture~\ref{conj:m-theory-parent} with
  Theorem~\ref{thm:hcs-super-open} at $2$-loop and higher: the 6D
  $(2,0)$ $A_{K-1}$ theory on $\bR^{1,3}\times T^{2}$ with $K$
  M5-branes wrapped on $\bP^{1}\subset\bY$, normal-bundle twist
  $\cO(-1)\oplus\cO(-1)$, $\Omega$-deformation on transverse
  Taub--NUT, and Costello--Gwilliam twist on the CY$_{3}$ target
  quantises after $T^{2}$-reduction to
  $Y(\widehat{\fgl}(1|1))^{\mathrm{con}}$ to all orders in $\hbar$.
  \emph{Falsifier.} Compute the 2-loop Costello--Paquette BV-BRST
  anomaly cocycle on the super-abelian $(1|1)$-stack; a non-trivial
  cohomology class at $O(\hbar^{2})$ obstructs the
  factorisation-algebra quantisation and falsifies the all-orders
  claim at $K=1$. For $K\geq 3$: the cubic-Casimir anomaly
  $A_{\mathrm{anom}} = d^{abc}d^{abc}/(2\pi)^{3}$ obstructs
  non-abelian extension; $K=2$ is special because $d^{abc} = 0$ on
  $\fsl_{2}$. Explicit Costello--Paquette super-extension covering
  the $(-1,-1)$-twisted normal bundle is absent from current
  literature.
\end{enumerate}

Secondary frontiers, auxiliary to the master identification:

\begin{enumerate}[label=\emph{$\mathsf F_{\arabic*}$.}, start=5, leftmargin=2.4em]
\item \emph{Hopf pairing non-degeneracy beyond weight $(2,2)$.}
  \ClaimStatusPartialPrimaryLit\
  Shapovalov determinants are known up to weight $(2,2)$ with generic
  non-degeneracy verified and two Kazhdan--Lusztig-positive chambers
  separated by the conifold BPS wall $\hbar = 1$. \emph{Falsifier.} Compute
  $\det\mathrm{Shap}_{(3,3)}$ as a polynomial in $\hbar$ and exhibit
  a zero off the Kac denominator locus
  $\prod_{r,s}(\hbar - \hbar_{r,s})$; such a zero falsifies the
  conjectured higher-weight generic non-degeneracy and the
  KL-positivity persistence across the BPS wall.

\item \emph{Chamber-character match at higher weight.}
  \ClaimStatusPartialPrimaryLit\
  Coefficient-by-coefficient agreement across R1--R5 verified at
  weights $m+n\leq 4$; higher weights open.
  \emph{Falsifier.} Any pair $\mathsf R_{i},\mathsf R_{j}$ producing
  distinct structure constants at a weight $(m,n)$ with
  $m+n\geq 5$, after chamber normalisation
  (Theorem~\ref{thm:chamber-inv-pairing}) and the native-basis
  intertwiner $\iota^{\mathsf R_{i}\mathsf R_{j}}$ of
  Theorem~\ref{thm:master}(i), falsifies the five-route match.

\item \emph{Compact CY$_{3}$ replacement machinery.}
  \ClaimStatusOpen\
  Theorem~\ref{thm:compact-obstruction} establishes four
  obstructions to direct transfer; the \emph{positive} side — an
  explicit automorphic-form replacement on compact CY$_{3}$ — is
  open. Gritsenko--Nikulin \texttt{alg-geom/9504006} and Borcherds
  \texttt{alg-geom/9609022} supply partial data on $K3\times E$
  (Vol~III Part~IV). \emph{Falsifier.} Construct a non-toric compact
  CY$_{3}$ admitting a finite quiver with potential $(Q,W)$ with
  Ginzburg dga vanishing $\HH^{<0}(J)=0$ whose CoHA admits a
  chart-wise presentation equivariant under a continuous torus of
  rank $\geq 2$; any such example falsifies the simultaneous
  necessity of all four obstructions (O1)--(O4).

\item \emph{$\chi_{\mathrm{BLLPRvR}}$ identification at $\cN=1$.}
  \ClaimStatusConjectured\
  The 4D/2D chiral-algebra functor of
  Beem--Lemos--Liendo--Peelaers--Rastelli--van~Rees
  \texttt{arXiv:1312.5344} is defined for $\cN\geq 2$ 4D SCFTs;
  $T[\bY]$ at the resolved conifold is $\cN=1$. The claim
  $Y^{+}(\widehat{\fgl}(1|1)) = \chi_{\mathrm{BLLPRvR}}(T[\bY])$
  requires an $\cN=1$-compatible extension, e.g.\ via $R$-symmetry
  enhancement at the IR superconformal point, or via a
  partial-topological twist recovering $\cN=2$ bookkeeping on a
  sub-Coulomb locus. \emph{Falsifier.} Exhibit the IR superconformal
  fixed point of $T[\bY]$ and show its $U(1)_{R}$-current does not
  enhance to $SU(2)_{R}$; equivalently, prove that no
  $\cN=1\to\cN=2$ enhancement exists in the deep IR of $T[\bY]$,
  closing the applicability of $\chi_{\mathrm{BLLPRvR}}$.
\end{enumerate}

Each $\mathsf F_{i}$ is falsifiable by a concrete computation of
bounded combinatorial complexity; none is a metaphysical conjecture.
The four primary items $\mathsf F_{1}$--$\mathsf F_{4}$ share the
obstruction — absence of $E_{3}$-super-formality and of
Costello--Paquette super-extensions to CY$_{3}$ with
$(-1,-1)$-twisted normal bundle — consistent with these being four
facets of a single open problem: the chain-level quantisation of the
super-BPS $E_{3}$-centre.

% ============================================================================
\section{Master theorem}
% ============================================================================

\begin{theorem}[Master identification]
\label{thm:master}
\ClaimStatusProvedElsewhere\ \emph{(at the super-algebra level,
  primary source Li--Yamazaki \texttt{arXiv:2003.08909} §8.3.6;
  at the ungraded-shadow level, primary source MMNS
  \texttt{arXiv:1107.5017} combined with Davison--Meinhardt
  \texttt{arXiv:1601.02479} integrality.)}
Let $\bY = \mathrm{Tot}\bigl(\cO_{\bP^{1}}(-1)\oplus\cO_{\bP^{1}}(-1)\bigr)$
be the resolved conifold and
$(Q_{\mathrm{con}}, W_{\mathrm{con}})$ the Klebanov--Witten two-vertex
quiver with quartic cyclic potential. Write $Y^{+}(\widehat{\fgl}(1|1))$
for the positive half of the affine super-Yangian of
$\widehat{\fgl}(1|1)$ and
$Y^{+}(\widehat{\fgl}(1|1))^{\mathrm{con}}$ for its conifold
specialisation (Li--Yamazaki \texttt{arXiv:2003.08909} §8.3.6, bond
factors~\eqref{eq:bond-factors}). Then
\[
  \boxed{
    \;\CoHA(\bY) \;\stackrel{\cong}{\longrightarrow}\;
    Y^{+}\!\bigl(\widehat{\fgl}(1|1)\bigr)^{\mathrm{con}}\;
  }
\]
is an isomorphism of \emph{$\Z_{2}$-graded associative bialgebras in
super-vector spaces over $\C((\hbar))$}, with structure maps
$(\mu,\iota,\Delta,\epsilon)$ identified as follows:
\begin{enumerate}[label=\emph{(\alph*)}, leftmargin=2em]
\item \emph{Multiplication $\mu$.} The critical CoHA convolution
  product on $\CoHA(\bY)$ coincides with the super-shuffle star
  product on $Y^{+}(\widehat{\fgl}(1|1))^{\mathrm{con}}$ with bond
  factors~\eqref{eq:bond-factors} (Li--Yamazaki eq.~(8.125)).
\item \emph{Unit $\iota$.} The identity class $[\mathrm{pt}]\in
  \CoHA(\bY)_{(0,0)}$ maps to the vacuum $1 \in Y^{+}$.
\item \emph{Coproduct $\Delta$.} The Kontsevich--Soibelman localised
  integration map (\texttt{arXiv:1006.2706} §6.3) equals the
  Drinfeld-new coproduct
  $\Delta(e^{(a)}(z)) = e^{(a)}(z)\otimes 1 +
   \phi^{+,(a)}(z)\otimes e^{(a)}(z)$ in
  $Y^{+}(\widehat{\fgl}(1|1))^{\mathrm{con}}$, where
  $\phi^{+,(a)}(z) = \sum_{k\geq 0}\phi^{+,(a)}_{k}z^{-k-1}$ is the
  Drinfeld-new positive Cartan current on vertex $a\in\{0,1\}$.
\item \emph{Counit $\epsilon$.} $\epsilon|_{\CoHA_{(0,0)}} =
  \mathrm{id}_{\C}$, $\epsilon|_{\CoHA_{(m,n)}\neq(0,0)} = 0$.
\end{enumerate}
An \emph{antipode} is \textbf{not} asserted on
$Y^{+}(\widehat{\fgl}(1|1))^{\mathrm{con}}$ itself: the positive half
is a bialgebra but not a Hopf algebra. The antipode lives on the
Drinfeld double $D(Y^{+})$ alone; its scope is recorded in item (ii)
below.
\begin{enumerate}[label=\emph{(\roman*)}, leftmargin=2em]
\item \emph{Five presentations produce isomorphic bialgebras.}
  The five routes R1--R5
  (Sections~\ref{sec:shuffle-route},~\ref{sec:cech-route},~\ref{sec:vdb-route},~\ref{sec:rtt-route},~\ref{sec:chiral-route})
  each produce a bialgebra; any two are related by a native-basis
  intertwiner $\iota^{\mathsf R_{i}\mathsf R_{j}}$ that is an
  $\hbar$-linear bialgebra isomorphism on every weight-graded piece
  $(\CoHA)_{(m,n)}$ with $m+n\leq 4$, agreeing at the leading
  $(1,1)$-weight imaginary-root generator with coefficient
  $(\varepsilon_{1}\varepsilon_{2})^{-1}$ in each native basis.
  Agreement at $m+n\geq 5$ is $\mathsf F_{6}$
  (Section~\ref{sec:open-frontiers}). Routes R1, R4, R5 resolve the
  central Cartan current $K(z) = h^{(0)}(z) + h^{(1)}(z)$ via
  residue extraction at the coincidence locus
  $z^{(0)} = z^{(1)}$; routes R2, R3 resolve the Chevalley Cartan
  $H(z) = h^{(0)}(z) - h^{(1)}(z)$ and project $K$ to zero under the
  chart-wise $Y^{+}(\widehat{\fgl}_{1})$-stalk stencil
  (Remark~\ref{rem:super-vs-even}). The intertwiner
  $\iota^{\mathsf R_{i}\mathsf R_{j}}$ absorbs both the
  native-basis change and the $H\leftrightarrow K$ reassignment on
  the imaginary-root line.
\item \emph{Drinfeld double and Hopf status.} The double
  $D(Y^{+}) = Y(\widehat{\fgl}(1|1))^{\mathrm{con}}$ carries an
  antipode; it is a \emph{strict $\Z_{2}$-graded Hopf algebra} at
  rational, trigonometric, and toroidal spectral type (Etingof--Kazhdan
  cohomology $H^{3}_{\mathrm{EK}}(\fgl(1|1);\fgl(1|1)^{\otimes 3}) = 0$,
  Theorem~\ref{thm:hopf-status}), and a \emph{quasi-Hopf algebra with
  dynamical associator} $\Phi_{\mathrm{ell}}$
  (Felder--Jimbo--Konno) at elliptic spectral type.
\item \emph{BPS Lie algebra.} $\fg_{\mathrm{BPS}}^{\mathrm{con}}$ has
  Klebanov--Witten affine root structure: real roots
  $n\delta\pm\alpha$ of multiplicity $1$, imaginary roots $n\delta$ of
  super-dimension $2$, spanned by $\{H_{n}, K_{n}\}$ with $H$ the
  Chevalley Cartan and $K$ the central Cartan
  (Proposition~\ref{prop:super-brackets}).
\item \emph{Chamber structure.} The stability space is a
  $W(\widehat{\fsl}_{2})$-torsor of shifted quiver Yangian
  representations of a chamber-independent algebra; the
  Joyce--Song--Kontsevich--Soibelman pairing is literally equal
  across chambers, not merely projectively equivariant
  (Theorems~\ref{thm:chambers},~\ref{thm:chamber-inv-pairing}).
\item \emph{Ungraded shadow is a different algebra.} The
  supertrace projection $\mathrm{str}\colon\widehat{\fgl}(1|1)
  \twoheadrightarrow\widehat{\fsl}_{2}$ on the even subalgebra
  annihilates the central current $K$ and retains only the Chevalley
  $H$ (Remark~\ref{rem:super-vs-even}). It induces a surjection
  $Y^{+}(\widehat{\fgl}(1|1))^{\mathrm{con}}
   \twoheadrightarrow Y^{+}(\widehat{\fsl}_{2})^{\mathrm{con}}$
  of bialgebras whose kernel is the two-sided ideal generated by
  $K$. The target $Y^{+}(\widehat{\fsl}_{2})^{\mathrm{con}}$ is the
  MMNS motivic DT algebra \texttt{arXiv:1107.5017} and is a
  \emph{distinct algebra} from $Y^{+}(\widehat{\fgl}(1|1))^{\mathrm{con}}$:
  the imaginary-root line of the super source is two-dimensional
  ($H + K$), while the shadow's imaginary-root line is one-dimensional
  ($H$ only). Both are theorems at their respective scopes --- the super
  form is primary for the categorified CoHA (Li--Yamazaki
  \texttt{arXiv:2003.08909} §8.3.6), the shadow is primary for
  motivic/numerical invariants (MMNS combined with Davison--Meinhardt
  integrality \texttt{arXiv:1601.02479}) --- and the relation between
  them is surjection, not isomorphism.
\end{enumerate}
Transfer to arbitrary local toric CY$_{3}$ without compact $4$-cycles
is Proposition~\ref{prop:toric-transfer}; the four-way obstruction to
direct transfer to compact CY$_{3}$ is
Theorem~\ref{thm:compact-obstruction}; the replacement
automorphic-form machinery on $K3\times E$, quintic, and
Gritsenko--Nikulin--Borcherds-lift compact CY$_{3}$ is
Remark~\ref{rem:mnop-replacement} and frontier
$\mathsf F_{7}$ (Section~\ref{sec:open-frontiers}).
\end{theorem}

\begin{remark}[Scope summary: what is proved, what is not]
\label{rem:master-theorem-scope}
The isomorphism $\CoHA(\bY) \cong
Y^{+}(\widehat{\fgl}(1|1))^{\mathrm{con}}$ of item~(i) is established
as bialgebras on weight-graded pieces with $m+n\leq 4$ via explicit
structure-constant match across R1--R5; extension to all weights is
the conjectural content of frontier $\mathsf F_{6}$. The antipode
claim of item~(ii) is established on $D(Y^{+})$, not on $Y^{+}$
itself. Item~(v) distinguishes the super
bialgebra from its ungraded shadow as two different algebras
(surjection of the super form onto the shadow), not as two
presentations of one algebra. The refined
$(\varepsilon_{1},\varepsilon_{2})$-structure underlying
Conjectures~\ref{conj:algebraic-super-rigidity}
and~\ref{conj:ikv-super-character} is the open content of frontiers
$\mathsf F_{1}, \mathsf F_{2}$. Hopf-pairing non-degeneracy at weight
$>(2,2)$ and the $c_{11} = -1$ asymmetric correction are
$\mathsf F_{5}, \mathsf F_{3}$.
\end{remark}

\end{document}
